\chapter{Computer Networking and Subnetting}
\label{ch:networking_subnetting}

This chapter covers the fundamental computer networking principles, architectural models, transmission media, and core protocols tested in Bangladesh government technical recruitment examinations (including Assistant Programmer, Assistant Network Engineer, Assistant Maintenance Engineer, and related IT/ICT positions).

The focus is on architectural frameworks (OSI and TCP/IP), protocol mechanics (TCP versus UDP, sliding window flow control, congestion avoidance), network hardware operation (switches, routers, and hubs), logical addressing (IPv4, IPv6, CIDR, private address spaces), and examination-oriented comparative analyses.

\newpage
\begin{center}
\vspace{0.5em}
{\color{primary}\rule{0.45\textwidth}{1.2pt}\\[4pt]}
{\Large\bfseries\color{primary} Section I: Fundamentals of Computer Networks\\[2pt]}
{\color{primary}\rule{0.45\textwidth}{1.2pt}}
\vspace{0.5em}
\end{center}

\section{Fundamentals of Computer Networks}

\subsection{1.1 What is a computer network?}

A \textbf{computer network} is a collection of interconnected computing devices that communicate and exchange data using transmission media and predefined communication protocols.

A network enables devices to share:

\begin{itemize}
    \item Data and files.
    \item Hardware resources, such as printers and storage.
    \item Software and services.
    \item Internet connectivity.
    \item Communication services, such as email, messaging, and video conferencing.
\end{itemize}

\textbf{Example:} In an office network, multiple computers connect to a central switch. The switch connects to a router, which provides access to the Internet. Employees can share files, access a common printer, and communicate with external systems.

\subsection{1.2 Basic components of a computer network}

\begin{table}[htbp]
\centering
\begin{tabularx}{\textwidth}{p{4.2cm} X X}
\toprule
\textbf{Component} & \textbf{Purpose} & \textbf{Example} \\
\midrule
Sender & Generates and sends data & Client computer \\
Receiver & Accepts transmitted data & Server \\
Message & Information being transmitted & File, email, packet \\
Transmission medium & Carries signals & Fiber optic, copper cable, radio \\
Network interface & Connects a device to a network & Ethernet NIC, Wi-Fi adapter \\
Protocol & Defines communication rules & TCP, IP, HTTP \\
Network device & Connects or manages network traffic & Switch, router, access point \\
\bottomrule
\end{tabularx}
\end{table}

\subsection{1.3 Types of computer networks}

Networks are commonly classified according to their geographical coverage.

\textbf{A. Personal Area Network (PAN)}

A PAN connects devices around an individual, usually within a few metres.

Examples:

\begin{itemize}
    \item Smartphone connected to wireless earbuds through Bluetooth.
    \item Laptop connected to a mobile hotspot.
    \item Smartwatch communicating with a smartphone.
\end{itemize}

\textbf{B. Local Area Network (LAN)}

A LAN connects devices within a limited geographical area, such as a home, office, school, or laboratory.

Characteristics:

\begin{itemize}
    \item Usually privately owned and managed.
    \item High data transmission speeds.
    \item Relatively low latency.
    \item Commonly uses Ethernet or Wi-Fi.
\end{itemize}

Example: A university computer laboratory where all computers connect to a central switch.

\textbf{C. Metropolitan Area Network (MAN)}

A MAN connects multiple LANs across a city or metropolitan area.

Example: A network connecting several branches of a bank throughout Chattogram.

\textbf{D. Wide Area Network (WAN)}

A WAN connects networks across large geographical areas, such as countries or continents.

Characteristics:

\begin{itemize}
    \item Covers large geographical distances.
    \item May use leased lines, fiber-optic links, satellite connections, or public infrastructure.
    \item Often involves multiple service providers.
\end{itemize}

Example: A bank connecting its branches across Bangladesh.

\textbf{The Internet is the largest example of an interconnected global network of networks.}

\subsection{1.4 Comparison of network types}

\begin{table}[htbp]
\centering
\begin{tabularx}{\textwidth}{p{2.5cm} X X}
\toprule
\textbf{Network} & \textbf{Coverage} & \textbf{Example} \\
\midrule
PAN & Personal space & Bluetooth \\
LAN & Building or campus & Office network \\
MAN & City & City-wide institutional network \\
WAN & Country or continent & Nationwide bank network \\
\bottomrule
\end{tabularx}
\end{table}

\subsection{1.5 Network topologies}

A \textbf{network topology} describes how devices are physically or logically connected.

\textbf{A. Bus topology}

All devices share a single main communication cable called the backbone.

Characteristics:

\begin{itemize}
    \item Simple and inexpensive for small networks.
    \item Requires less cable than some other topologies.
    \item Failure of the backbone can disrupt the entire network.
    \item Collisions may occur when multiple devices transmit simultaneously.
    \item Terminators are used at both ends of the traditional bus cable.
\end{itemize}

\textbf{B. Star topology}

All devices connect to a central device, typically a switch.

Characteristics:

\begin{itemize}
    \item Easy to install and manage.
    \item Failure of one connecting cable usually affects only one device.
    \item Failure of the central switch can disrupt the entire network.
    \item Commonly used in modern Ethernet LANs.
\end{itemize}

\textbf{C. Ring topology}

Each device connects to two neighbouring devices, forming a closed loop.

Characteristics:

\begin{itemize}
    \item Data travels around the ring.
    \item Traditional implementations may use token passing.
    \item A failure can disrupt communication unless redundancy or bypass mechanisms exist.
\end{itemize}

\textbf{D. Mesh topology}

Devices are interconnected through multiple paths.

In a full mesh network, every device has a direct connection to every other device.

For $n$ devices, the number of links in a full mesh is:


\[
\boxed{L=\frac{n(n-1)}{2}}
\]


Example: For 5 devices:


\[
L=\frac{5(5-1)}{2}=\frac{20}{2}=10
\]


Characteristics:

\begin{itemize}
    \item High redundancy and reliability.
    \item Multiple alternative communication paths.
    \item Expensive and complex to implement.
\end{itemize}

\textbf{E. Tree topology}

A hierarchical topology where multiple star networks connect through a backbone or higher-level devices.

Commonly used in large institutional networks.

\textbf{F. Hybrid topology}

A combination of two or more topologies, such as star and mesh.

\subsection{1.6 Network topology comparison}

\begin{table}[htbp]
\centering
\begin{tabularx}{\textwidth}{p{2.6cm} X X}
\toprule
\textbf{Topology} & \textbf{Main advantage} & \textbf{Main disadvantage} \\
\midrule
Bus & Low cost and simple & Backbone failure affects all devices \\
Star & Easy troubleshooting & Central device is a failure point \\
Ring & Predictable access in token-based systems & Ring disruption can affect communication \\
Mesh & High redundancy & High cost \\
Tree & Scalable hierarchy & Dependence on higher-level devices \\
Hybrid & Flexible design & More complex management \\
\bottomrule
\end{tabularx}
\end{table}

\subsection{1.7 Communication modes}

Communication mode describes the direction in which data can travel.

\textbf{Simplex:} Communication occurs in only one direction.

Example: Traditional television broadcasting.

\textbf{Half-duplex:} Communication can occur in both directions, but not simultaneously.

Example: Walkie-talkie.

\textbf{Full-duplex:} Communication occurs in both directions simultaneously.

Example: A telephone conversation.

\begin{table}[htbp]
\centering
\begin{tabularx}{\textwidth}{p{3.2cm} X X}
\toprule
\textbf{Mode} & \textbf{Direction} & \textbf{Simultaneous transmission?} \\
\midrule
Simplex & One-way & No \\
Half-duplex & Two-way & No \\
Full-duplex & Two-way & Yes \\
\bottomrule
\end{tabularx}
\end{table}

\subsection{1.8 Data transmission techniques}

\textbf{Unicast:} One sender communicates with one specific receiver.

\textbf{Broadcast:} One sender transmits to all devices within a broadcast domain.

\textbf{Multicast:} One sender transmits to a selected group of receivers.

\textbf{Anycast:} Data is delivered to one suitable member of a group, commonly the nearest according to routing policy.

Examples:

\begin{itemize}
    \item Unicast: A client requests a webpage from a server.
    \item Broadcast: An IPv4 ARP request on a LAN.
    \item Multicast: Streaming data to subscribed receivers.
    \item Anycast: DNS services using distributed server locations.
\end{itemize}

\subsection{1.9 Circuit switching vs packet switching}

\begin{table}[htbp]
\centering
\begin{tabularx}{\textwidth}{p{4.2cm} X X}
\toprule
\textbf{Feature} & \textbf{Circuit switching} & \textbf{Packet switching} \\
\midrule
Connection & Dedicated communication path & Data divided into packets \\
Resource allocation & Reserved for the connection & Shared dynamically \\
Efficiency & Can waste reserved capacity & Efficient resource sharing \\
Delay & Setup delay, then predictable path characteristics & Variable queuing and transmission delay \\
Example & Traditional telephone networks & Internet \\
\bottomrule
\end{tabularx}
\end{table}

\textbf{Important:} The Internet primarily uses packet switching.


\newpage

\newpage
\begin{center}
\vspace{0.5em}
{\color{primary}\rule{0.45\textwidth}{1.2pt}\\[4pt]}
{\Large\bfseries\color{primary} Section II: The OSI Reference Model\\[2pt]}
{\color{primary}\rule{0.45\textwidth}{1.2pt}}
\vspace{0.5em}
\end{center}

\section{The OSI Reference Model}

\subsection{2.1 What is the OSI model?}

The \textbf{Open Systems Interconnection (OSI) model} is a conceptual networking framework developed by the International Organization for Standardization (ISO).

It divides network communication into seven logical layers. Each layer has specific responsibilities and provides services to the layer above it.

The OSI model helps engineers:

\begin{itemize}
    \item Understand network communication.
    \item Design interoperable systems.
    \item Troubleshoot network problems.
    \item Separate networking responsibilities into manageable layers.
\end{itemize}

\textbf{The seven OSI layers}

\begin{quote}
\textbf{Layer 7: }\\[2pt]
\textit{**Application**}
\end{quote}

Network services for applications

\begin{quote}
\textbf{Layer 6: }\\[2pt]
\textit{**Presentation**}
\end{quote}

Data format, encryption and compression

\begin{quote}
\textbf{Layer 5: }\\[2pt]
\textit{**Session**}
\end{quote}

Session establishment and management

\begin{quote}
\textbf{Layer 4: }\\[2pt]
\textit{**Transport**}
\end{quote}

End-to-end delivery and reliability

\begin{quote}
\textbf{Layer 3: }\\[2pt]
\textit{**Network**}
\end{quote}

Logical addressing and routing

\begin{quote}
\textbf{Layer 2: }\\[2pt]
\textit{**Data Link**}
\end{quote}

Frames, MAC addressing and local delivery

\begin{quote}
\textbf{Layer 1: }\\[2pt]
\textit{**Physical**}
\end{quote}

Transmission of raw bits

\subsection{2.2 Layer 1: Physical Layer}

The Physical Layer is responsible for transmitting raw binary bits (0s and 1s) through a physical transmission medium.

\textbf{Main responsibilities:}

\begin{itemize}
    \item Transmission of electrical, optical, or radio signals.
    \item Defining voltage levels and signal characteristics.
    \item Bit synchronization.
    \item Data transmission rates.
    \item Physical connectors and cable specifications.
    \item Physical network topology.
\end{itemize}

\textbf{Examples of technologies and components:}

\begin{itemize}
    \item Ethernet physical standards.
    \item Fiber-optic cables.
    \item UTP cables.
    \item Repeaters.
    \item Hubs.
    \item Physical connectors such as RJ45.
\end{itemize}

\textbf{PDU (Protocol Data Unit):} Bits.

\textbf{Exam question:} Which OSI layer is responsible for transmitting raw bits?

Answer: Physical Layer.

\subsection{2.3 Layer 2: Data Link Layer}

The Data Link Layer provides node-to-node data delivery over a directly connected network link.

\textbf{Main responsibilities:}

\begin{itemize}
    \item Framing: Dividing data into frames.
    \item Physical addressing using MAC addresses.
    \item Error detection, commonly using CRC.
    \item Media access control.
    \item Link-level flow control in technologies that support it.
\end{itemize}

It is commonly divided into two sublayers:

\begin{enumerate}
    \item LLC (Logical Link Control).
    \item MAC (Media Access Control).
\end{enumerate}

\textbf{Examples:}

\begin{itemize}
    \item Ethernet (IEEE 802.3).
    \item Wi-Fi MAC (IEEE 802.11).
    \item PPP.
    \item Network bridges and Layer 2 switches.
\end{itemize}

\textbf{PDU:} Frame.

\textbf{Addressing:} MAC address.

\textbf{Exam question:} Which OSI layer uses MAC addresses?

Answer: Data Link Layer.

\subsection{2.4 Layer 3: Network Layer}

The Network Layer is responsible for delivering packets across different interconnected networks.

\textbf{Main responsibilities:}

\begin{itemize}
    \item Logical addressing.
    \item Routing.
    \item Packet forwarding.
    \item Path selection.
    \item Fragmentation in IPv4, where applicable.
\end{itemize}

\textbf{Important protocols:}

\begin{itemize}
    \item IPv4.
    \item IPv6.
    \item ICMP.
    \item Routing protocols such as OSPF and BGP.
\end{itemize}

\textbf{Devices:} Routers and Layer 3 switches.

\textbf{PDU:} Packet (or datagram in IP terminology).

\textbf{Addressing:} IP address.

\textbf{Example:} A computer in Chattogram communicates with a server located in Dhaka. The Network Layer handles logical addressing and routing across networks.

\subsection{2.5 Layer 4: Transport Layer}

The Transport Layer provides communication between application processes running on different hosts.

Its responsibilities include:

\begin{itemize}
    \item Segmentation and reassembly.
    \item Process-to-process delivery.
    \item Port addressing.
    \item Flow control.
    \item Error recovery and retransmission in reliable protocols.
    \item Connection management.
\end{itemize}

\textbf{Important protocols:}

\begin{itemize}
    \item TCP (Transmission Control Protocol).
    \item UDP (User Datagram Protocol).
\end{itemize}

\textbf{PDU:}

\begin{itemize}
    \item TCP segment.
    \item UDP datagram.
\end{itemize}

\textbf{Addressing:} Port numbers.

Examples:

\begin{itemize}
    \item TCP port 443: HTTPS.
    \item TCP port 80: HTTP.
    \item UDP port 53: DNS commonly uses UDP, though DNS can also use TCP.
\end{itemize}

\subsection{2.6 Layer 5: Session Layer}

The Session Layer manages communication sessions between applications.

Responsibilities:

\begin{itemize}
    \item Establishing sessions.
    \item Maintaining sessions.
    \item Terminating sessions.
    \item Dialog control.
    \item Synchronization and checkpoints in relevant systems.
\end{itemize}

Example: Managing an ongoing communication session between two application systems.

\textbf{Exam note:} In practical TCP/IP implementations, session-related functionality is often handled by application-layer protocols and libraries rather than a distinct OSI layer.

\subsection{2.7 Layer 6: Presentation Layer}

The Presentation Layer deals with how data is represented and interpreted.

Responsibilities:

\begin{itemize}
    \item Data translation.
    \item Character encoding.
    \item Encryption and decryption.
    \item Compression and decompression.
    \item Data serialization and formatting.
\end{itemize}

Examples:

\begin{itemize}
    \item UTF-8 character encoding.
    \item JPEG image representation.
    \item PNG.
    \item Data serialization formats such as JSON (as an application-level format).
\end{itemize}

\textbf{Exam question:} Which OSI layer is associated with encryption, compression, and data translation?

Answer: Presentation Layer.

\subsection{2.8 Layer 7: Application Layer}

The Application Layer provides network services directly to application software.

It does not mean the user application itself; rather, it provides the protocols and services through which applications communicate over a network.

\textbf{Important protocols:}

\begin{table}[htbp]
\centering
\begin{tabularx}{\textwidth}{p{2.6cm} X X}
\toprule
\textbf{Protocol} & \textbf{Full form} & \textbf{Purpose} \\
\midrule
HTTP & Hypertext Transfer Protocol & Web communication \\
HTTPS & HTTP Secure & Secure web communication \\
FTP & File Transfer Protocol & File transfer \\
SMTP & Simple Mail Transfer Protocol & Sending email \\
POP3 & Post Office Protocol v3 & Retrieving email \\
IMAP & Internet Message Access Protocol & Accessing and synchronizing email \\
DNS & Domain Name System & Resolving domain names \\
DHCP & Dynamic Host Configuration Protocol & Automatic IP configuration \\
SNMP & Simple Network Management Protocol & Network monitoring and management \\
SSH & Secure Shell & Secure remote access \\
Telnet & Telecommunication Network & Remote terminal access without encryption by default \\
\bottomrule
\end{tabularx}
\end{table}

\textbf{PDU:} Data.

\subsection{2.9 OSI layers: complete summary}

\begin{table}[htbp]
\centering
\begin{tabularx}{\textwidth}{p{1.2cm} p{2.5cm} p{2.8cm} p{2.2cm} X}
\toprule
\textbf{Layer} & \textbf{Name} & \textbf{PDU} & \textbf{Addressing} & \textbf{Key responsibility} \\
\midrule
7 & Application & Data & Application-specific & Network services \\
6 & Presentation & Data & — & Translation, encryption, compression \\
5 & Session & Data & — & Session management \\
4 & Transport & Segment/Datagram & Port & End-to-end transport \\
3 & Network & Packet & IP & Routing \\
2 & Data Link & Frame & MAC & Local delivery \\
1 & Physical & Bits & — & Signal transmission \\
\bottomrule
\end{tabularx}
\end{table}

\subsection{2.10 Data encapsulation and decapsulation}

When data travels from a sender to a receiver, each layer adds or removes protocol information.

\textbf{Encapsulation:} The process of adding headers (and sometimes trailers) as data moves down the protocol stack.

\textbf{Decapsulation:} The process of removing and interpreting those headers as data moves up the stack at the receiver.

\textbf{Sender: Encapsulation}

Application data

TCP/UDP header + Data

IP header + Segment/Datagram

Frame header/trailer + Packet

Bits transmitted through the medium

\textbf{Receiver: Decapsulation}

The receiving device processes the bits, extracts the frame, packet and transport data, and delivers the application data to the appropriate application.

\textbf{Important exam point:} A router primarily processes Layer 3 information, while a traditional Layer 2 switch primarily processes Layer 2 information.

\subsection{2.11 Mnemonic for remembering the OSI layers}

From Layer 7 to Layer 1:

\textbf{A}ll \textbf{P}eople \textbf{S}eem \textbf{T}o \textbf{N}eed \textbf{D}ata \textbf{P}rocessing.

\begin{itemize}
    \item Application
    \item Presentation
    \item Session
    \item Transport
    \item Network
    \item Data Link
    \item Physical
\end{itemize}

From Layer 1 to Layer 7:

\textbf{P}lease \textbf{D}o \textbf{N}ot \textbf{T}hrow \textbf{S}ausage \textbf{P}izza \textbf{A}way.


\newpage

\newpage
\begin{center}
\vspace{0.5em}
{\color{primary}\rule{0.45\textwidth}{1.2pt}\\[4pt]}
{\Large\bfseries\color{primary} Section III: The TCP/IP Protocol Architecture\\[2pt]}
{\color{primary}\rule{0.45\textwidth}{1.2pt}}
\vspace{0.5em}
\end{center}

\section{The TCP/IP Protocol Architecture}

\subsection{3.1 What is the TCP/IP model?}

The TCP/IP model is a practical networking architecture used by the Internet and most modern IP-based networks.

TCP/IP stands for \textbf{Transmission Control Protocol / Internet Protocol}.

Unlike the seven-layer OSI reference model, the commonly used TCP/IP model has four layers.

\begin{table}[htbp]
\centering
\begin{tabularx}{\textwidth}{p{4.2cm} X X}
\toprule
\textbf{TCP/IP layer} & \textbf{Corresponding OSI layers} & \textbf{Main responsibilities} \\
\midrule
Application & Application, Presentation, Session & Application services and data representation \\
Transport & Transport & End-to-end process communication \\
Internet & Network & IP addressing and routing \\
Network Access (Link) & Data Link, Physical & Local network delivery and physical transmission \\
\bottomrule
\end{tabularx}
\end{table}

Some textbooks use a five-layer Internet model by separating the Link layer into Data Link and Physical. For examinations, follow the model specified in the question.

\subsection{3.2 TCP/IP layers in detail}

\begin{quote}
\textbf{Question 1: Application Layer}\\[4pt]
\vspace{4pt}
\textbf{Analysis / Answer:} Provides services to network applications. Examples: - HTTP/HTTPS. - FTP. - SMTP. - DNS. - DHCP. - SSH.
\end{quote}
\begin{quote}
\textbf{Question 2: Transport Layer}\\[4pt]
\vspace{4pt}
\textbf{Analysis / Answer:} Provides process-to-process communication. Protocols: - TCP. - UDP. Responsibilities: - Port addressing. - Segmentation. - Reliability where supported. - Flow control.
\end{quote}
\begin{quote}
\textbf{Question 3: Internet Layer}\\[4pt]
\vspace{4pt}
\textbf{Analysis / Answer:} Responsible for logical addressing and routing between networks. Protocols: - IPv4. - IPv6. - ICMP.
\end{quote}
\begin{quote}
\textbf{Question 4: Network Access Layer}\\[4pt]
\vspace{4pt}
\textbf{Analysis / Answer:} Handles communication over the local network and physical transmission. Examples: - Ethernet. - Wi-Fi. - PPP.
\end{quote}
\subsection{3.3 OSI vs TCP/IP model}

\begin{table}[htbp]
\centering
\begin{tabularx}{\textwidth}{p{4.2cm} X X}
\toprule
\textbf{Feature} & \textbf{OSI} & \textbf{TCP/IP} \\
\midrule
Full form & Open Systems Interconnection & Transmission Control Protocol / Internet Protocol \\
Number of layers & 7 & 4 (commonly taught) \\
Nature & Reference model & Practical protocol architecture \\
Development & ISO & DARPA research origins \\
Application functions & Three separate upper layers & One Application Layer \\
Physical and Data Link & Separate & Combined into Link/Network Access \\
Usage & Learning, design, troubleshooting & Real-world Internet communication \\
\bottomrule
\end{tabularx}
\end{table}

\textbf{Important:} OSI and TCP/IP are not competing protocols. OSI is a reference model; TCP/IP is the protocol architecture used in real networks.


\newpage

\newpage
\begin{center}
\vspace{0.5em}
{\color{primary}\rule{0.45\textwidth}{1.2pt}\\[4pt]}
{\Large\bfseries\color{primary} Section IV: Transport Layer Protocols: TCP versus UDP\\[2pt]}
{\color{primary}\rule{0.45\textwidth}{1.2pt}}
\vspace{0.5em}
\end{center}

\section{Transport Layer Protocols: TCP versus UDP}

\subsection{4.1 Transmission Control Protocol (TCP)}

TCP is a \textbf{connection-oriented, reliable transport-layer protocol} that provides ordered delivery of data between two applications.

Before transmitting application data, TCP establishes a connection between the sender and receiver through a process called the \textbf{three-way handshake}.

Its major characteristics are:

\begin{enumerate}
    \item Connection-oriented communication.
    \item Reliable data delivery.
    \item Ordered delivery of data.
    \item Error detection and recovery.
    \item Acknowledgment of received data.
    \item Retransmission of lost segments.
    \item Flow control.
    \item Congestion control.
    \item Full-duplex communication.
\end{enumerate}

\textbf{Example:} When downloading a document, missing data must be retransmitted to ensure the downloaded file is complete. TCP provides this reliability.

\subsection{4.2 TCP three-way handshake}

TCP establishes a connection through three messages.

\begin{table}[htbp]
\centering
\begin{tabularx}{\textwidth}{p{3.0cm} p{3.2cm} X X}
\toprule
\textbf{Step} & \textbf{Sender} & \textbf{Receiver} & \textbf{Description} \\
\midrule
1 & SYN & — & Client requests a connection \\
2 & — & SYN-ACK & Server acknowledges and responds \\
3 & ACK & — & Client acknowledges the server's response \\
\bottomrule
\end{tabularx}
\end{table}

Suppose a client wants to connect to a web server.

\begin{enumerate}
    \item Client sends \texttt{SYN} with initial sequence number 100.
    \item Server responds with \texttt{SYN-ACK}, using its own sequence number 500 and acknowledging 101.
    \item Client sends \texttt{ACK}, acknowledging 501.
\end{enumerate}

The connection is now established.

\textbf{Important terms:}

\begin{itemize}
    \item SYN: Synchronizes sequence numbers and initiates a connection.
    \item ACK: Acknowledgment flag.
    \item FIN: Initiates an orderly connection termination.
    \item RST: Resets a connection.
\end{itemize}

\subsection{4.3 TCP connection termination}

TCP commonly uses a four-step termination process:

\begin{enumerate}
    \item One endpoint sends FIN.
    \item The other endpoint acknowledges with ACK.
    \item The second endpoint sends FIN when ready to close its side.
    \item The first endpoint sends the final ACK.
\end{enumerate}

TCP supports half-close: one endpoint can stop sending while still receiving data.

\subsection{4.4 TCP flow control and congestion control}

These two concepts are frequently confused in examinations.

\textbf{Flow control}

Flow control prevents a sender from overwhelming the receiver.

TCP uses a \textbf{sliding window mechanism}. The receiver advertises how much data it can currently accept through its receive window.

Example: If a receiver can accept only 4 KB of additional data, the sender must limit the amount of unacknowledged data accordingly.

\textbf{Congestion control}

Congestion control prevents excessive traffic from overwhelming the network.

TCP uses mechanisms such as:

\begin{itemize}
    \item Slow Start.
    \item Congestion Avoidance.
    \item Fast Retransmit.
    \item Fast Recovery (in applicable TCP implementations).
\end{itemize}

The congestion window (cwnd) controls how much data can be outstanding based on the sender's estimate of network capacity.

\begin{table}[htbp]
\centering
\begin{tabularx}{\textwidth}{X X}
\toprule
\textbf{Flow control} & \textbf{Congestion control} \\
\midrule
Protects the receiver & Protects the network \\
Based on receiver capacity & Based on network congestion \\
Uses receiver-advertised window & Uses congestion window \\
Prevents receiver overload & Prevents network overload \\
\bottomrule
\end{tabularx}
\end{table}

\subsection{4.5 User Datagram Protocol (UDP)}

UDP is a \textbf{connectionless transport-layer protocol} that provides lightweight, best-effort delivery of datagrams.

It does not establish a connection before sending data.

Characteristics:

\begin{itemize}
    \item Connectionless.
    \item No delivery acknowledgment at the UDP protocol level.
    \item No retransmission mechanism at the UDP protocol level.
    \item No guarantee of ordered delivery.
    \item No built-in flow control.
    \item No built-in congestion control.
    \item Lower protocol overhead than TCP.
    \item Supports broadcast and multicast applications where appropriate.
\end{itemize}

UDP includes a checksum for error detection. In IPv4, UDP checksum use is optional; in IPv6, it is normally mandatory, with limited exceptions.

\textbf{Example:} A live audio or video application may use UDP because receiving the latest data promptly can be more useful than waiting for an older lost packet to be retransmitted.

\subsection{4.6 TCP vs UDP comparison}

\begin{table}[htbp]
\centering
\begin{tabularx}{\textwidth}{p{4.2cm} X X}
\toprule
\textbf{Feature} & \textbf{TCP} & \textbf{UDP} \\
\midrule
Full form & Transmission Control Protocol & User Datagram Protocol \\
Connection & Connection-oriented & Connectionless \\
Reliability & Reliable delivery mechanism & Best-effort delivery \\
Acknowledgment & Yes & No UDP-level acknowledgment \\
Retransmission & Yes & No UDP-level retransmission \\
Ordering & Preserves byte-stream order & No ordering guarantee \\
Flow control & Yes & No built-in flow control \\
Congestion control & Yes & No built-in congestion control \\
Header size & Minimum 20 bytes & 8 bytes \\
Speed/overhead & More protocol overhead & Less protocol overhead \\
Communication & Unicast & Unicast, multicast and broadcast \\
Data handling & Byte stream & Datagram-based \\
Typical applications & Web, file transfer, email & DNS, voice/video, gaming, streaming \\
\bottomrule
\end{tabularx}
\end{table}

\textbf{TCP header vs UDP header}

\begin{table}[htbp]
\centering
\begin{tabularx}{\textwidth}{p{4.2cm} X X}
\toprule
\textbf{Field} & \textbf{TCP} & \textbf{UDP} \\
\midrule
Source port & Yes & Yes \\
Destination port & Yes & Yes \\
Sequence number & Yes & No \\
Acknowledgment number & Yes & No \\
Header length & Yes & No \\
Flags & Yes & No \\
Window size & Yes & No \\
Checksum & Yes & Yes \\
Length field & No separate UDP-style length field & Yes \\
\bottomrule
\end{tabularx}
\end{table}

\textbf{TCP minimum header:} 20 bytes.

\textbf{TCP maximum header:} 60 bytes.

\textbf{UDP header:} Exactly 8 bytes.

\subsection{4.7 Applications of TCP and UDP}

\begin{table}[htbp]
\centering
\begin{tabularx}{\textwidth}{X X}
\toprule
\textbf{Application} & \textbf{Common transport protocol} \\
\midrule
HTTP/1.1 & TCP \\
HTTP/2 & TCP \\
HTTP/3 & QUIC over UDP \\
HTTPS (traditional HTTP over TLS) & TCP \\
FTP & TCP \\
SMTP & TCP \\
POP3 & TCP \\
IMAP & TCP \\
DNS queries & Usually UDP; TCP also supported \\
DHCP & UDP \\
SSH & TCP \\
Telnet & TCP \\
VoIP & Often UDP \\
Online multiplayer games & Often UDP \\
NTP & UDP \\
\bottomrule
\end{tabularx}
\end{table}

\textbf{Important examination traps:}

\begin{itemize}
    \item HTTP/3 uses QUIC, which operates over UDP.
    \item DNS does not exclusively use UDP; TCP is also used.
    \item UDP is not inherently unreliable at the application level. An application can implement its own acknowledgments and retransmissions.
    \item TCP does not preserve application message boundaries. It provides a continuous byte stream.
    \item UDP preserves datagram boundaries.
\end{itemize}


\newpage

\newpage
\begin{center}
\vspace{0.5em}
{\color{primary}\rule{0.45\textwidth}{1.2pt}\\[4pt]}
{\Large\bfseries\color{primary} Section V: Network Interconnection Devices: Hub, Switch, and Router\\[2pt]}
{\color{primary}\rule{0.45\textwidth}{1.2pt}}
\vspace{0.5em}
\end{center}

\section{Network Interconnection Devices: Hub, Switch, and Router}

\subsection{5.1 Hub}

A hub is a basic Layer 1 networking device that connects multiple devices in a LAN.

When a hub receives a signal on one port, it repeats the signal to all other ports.

Characteristics:

\begin{itemize}
    \item Operates at the Physical Layer.
    \item Does not examine MAC or IP addresses.
    \item Sends incoming traffic to all other ports.
    \item All connected devices share the collision domain.
    \item Usually operates in half-duplex Ethernet environments.
\end{itemize}

Hubs are largely obsolete in modern Ethernet networks.

\subsection{5.2 Switch}

A \textbf{network switch} connects devices within a local area network and forwards Ethernet frames based on MAC addresses.

A traditional Layer 2 switch operates at the Data Link Layer.

Main responsibilities:

\begin{itemize}
    \item MAC address learning.
    \item Frame forwarding.
    \item Frame filtering.
    \item Maintaining a MAC address table.
    \item Separating collision domains by port.
\end{itemize}

\subsection{5.3 How does a switch work?}

A switch maintains a table called a \textbf{MAC address table}, also known as a CAM table.

This table associates MAC addresses with switch ports.

Consider this example:

\begin{table}[htbp]
\centering
\begin{tabularx}{\textwidth}{p{4.3cm} X}
\toprule
\textbf{MAC address} & \textbf{Switch port} \\
\midrule
AA:AA:AA:AA:AA:AA & Port 1 \\
BB:BB:BB:BB:BB:BB & Port 2 \\
CC:CC:CC:CC:CC:CC & Port 3 \\
\bottomrule
\end{tabularx}
\end{table}

Suppose Device A, connected to Port 1, sends a frame to Device B.

The switch performs the following operations:

\textbf{Step 1: Learning}

The switch examines the source MAC address of the incoming frame.

It learns that Device A is reachable through Port 1.

\textbf{Step 2: Destination lookup}

The switch checks its MAC address table for the destination MAC address of Device B.

\textbf{Step 3: Forwarding}

\begin{itemize}
    \item If the destination MAC address is known, the switch forwards the frame through the corresponding port.
    \item If the destination MAC address is unknown, it floods the frame to all eligible ports in the same VLAN except the incoming port.
\end{itemize}

\textbf{Step 4: Filtering}

If the destination device is connected to the same port on which the frame arrived, the switch does not forward that frame to another port.

\subsection{5.4 Switch forwarding methods}

\textbf{Store-and-forward switching}

The switch receives the complete frame, checks it for errors, and then forwards it.

Advantages:

\begin{itemize}
    \item Can discard corrupted frames.
    \item Supports frame buffering.
    \item Useful in networks where error checking is important.
\end{itemize}

\textbf{Cut-through switching}

The switch begins forwarding a frame after reading enough of its header to identify the destination port.

Advantages:

\- Lower forwarding latency.

Disadvantage:

\- May forward frames before their complete integrity has been checked.

\textbf{Fragment-free switching}

A variation of cut-through switching that waits until the first 64 bytes have been received before forwarding.

\subsection{5.5 Router}

A \textbf{router} is a Layer 3 networking device that connects different IP networks and forwards packets based on destination IP addresses.

Main responsibilities:

\begin{itemize}
    \item Connecting different networks.
    \item Selecting appropriate routes.
    \item Forwarding IP packets.
    \item Maintaining routing tables.
    \item Separating broadcast domains.
    \item Supporting routing protocols.
\end{itemize}

Examples:

\begin{itemize}
    \item Connecting a home LAN to the Internet.
    \item Connecting a bank's branch network to its head office.
    \item Connecting multiple institutional networks.
\end{itemize}

\subsection{5.6 How does a router work?}

Suppose a computer has IP address \texttt{192.168.1.10/24} and wants to communicate with a server at \texttt{8.8.8.8}.

The process is:

\begin{enumerate}
    \item The computer determines that \texttt{8.8.8.8} is outside its local subnet.
    \item It identifies its default gateway, for example \texttt{192.168.1.1}.
    \item It sends the packet in an Ethernet frame addressed to the router's local MAC address.
    \item The router receives the frame and examines the destination IP address.
    \item The router checks its routing table.
    \item It selects an appropriate next hop or outgoing interface.
    \item It forwards the packet toward the destination network.
\end{enumerate}

\textbf{Important:} When routing an IPv4 packet, a router normally decreases the TTL (Time To Live) field by at least one. If TTL reaches zero, the packet is discarded and an ICMP Time Exceeded message may be generated.

The Ethernet frame is rewritten for the next link, but the destination IP address generally remains unchanged during ordinary routing. NAT is a separate operation that may modify IP addresses and transport-layer ports.

\subsection{5.7 Routing table}

A routing table contains information used by a router to determine where packets should be forwarded.

A simplified routing table might look like this:

\begin{table}[htbp]
\centering
\begin{tabularx}{\textwidth}{p{4.2cm} X X}
\toprule
\textbf{Destination network} & \textbf{Next hop} & \textbf{Interface} \\
\midrule
192.168.1.0/24 & Directly connected & eth0 \\
192.168.2.0/24 & 10.0.0.2 & eth1 \\
0.0.0.0/0 & 10.0.0.1 & eth1 \\
\bottomrule
\end{tabularx}
\end{table}

The entry \texttt{0.0.0.0/0} represents the default IPv4 route.

\textbf{Longest Prefix Match (LPM):} When multiple routing entries match a destination IP address, the router generally selects the route with the longest matching network prefix.

For example, given:

\begin{itemize}
    \item \texttt{192.168.0.0/16}
    \item \texttt{192.168.1.0/24}
    \item \texttt{192.168.1.128/25}
\end{itemize}

A packet destined for \texttt{192.168.1.150} matches all three prefixes. The \texttt{/25} route is selected because it is the most specific match.

\subsection{5.8 Switch vs Router vs Hub}

\begin{table}[htbp]
\centering
\begin{tabularx}{\textwidth}{p{3.0cm} p{3.2cm} X X}
\toprule
\textbf{Feature} & \textbf{Hub} & \textbf{Switch} & \textbf{Router} \\
\midrule
OSI layer & Layer 1 & Usually Layer 2 & Layer 3 \\
Address used & None & MAC address & IP address \\
Main function & Repeats signals & Forwards frames & Routes packets \\
Network scope & LAN & LAN/VLAN & Different IP networks \\
Collision domains & One shared domain & Separate per port & Separate per interface \\
Broadcast domains & One shared domain & One per VLAN by default & Separated by default \\
Intelligent forwarding & No & Yes & Yes \\
Typical use & Legacy LAN & Modern LAN & Inter-network communication \\
\bottomrule
\end{tabularx}
\end{table}

\textbf{Important exam facts:}

\begin{itemize}
    \item A switch does not normally break up a broadcast domain when all its ports belong to the same VLAN.
    \item A router separates broadcast domains.
    \item Each switch port is a separate collision domain in modern switched Ethernet.
    \item A Layer 3 switch can perform routing, including inter-VLAN routing.
    \item A switch learns MAC addresses; a router uses IP routing information.
\end{itemize}

\subsection{5.9 Other important networking devices}

\begin{table}[htbp]
\centering
\begin{tabularx}{\textwidth}{p{4.2cm} X X}
\toprule
\textbf{Device} & \textbf{OSI layer (typical)} & \textbf{Function} \\
\midrule
Repeater & Layer 1 & Regenerates signals \\
Hub & Layer 1 & Repeats signals to all ports \\
Bridge & Layer 2 & Connects LAN segments using MAC addresses \\
Switch & Layer 2 & Forwards frames using MAC addresses \\
Router & Layer 3 & Routes packets using IP addresses \\
Gateway & Depends on implementation & Connects networks or translates protocols \\
Modem & Depends on technology & Modulates/demodulates signals \\
Wireless Access Point & Primarily Layer 2 & Connects wireless clients to a LAN \\
Firewall & Multiple layers & Controls traffic according to security rules \\
\bottomrule
\end{tabularx}
\end{table}


\newpage

\newpage
\begin{center}
\vspace{0.5em}
{\color{primary}\rule{0.45\textwidth}{1.2pt}\\[4pt]}
{\Large\bfseries\color{primary} Section VI: Internet Protocol Addressing: IPv4 versus IPv6\\[2pt]}
{\color{primary}\rule{0.45\textwidth}{1.2pt}}
\vspace{0.5em}
\end{center}

\section{Internet Protocol Addressing: IPv4 versus IPv6}

\subsection{6.1 What is an IP address?}

An \textbf{Internet Protocol (IP) address} is a logical address assigned to a network interface to identify it within an IP network.

IP addressing enables devices to communicate across interconnected networks.

There are two major versions:

\begin{enumerate}
    \item IPv4 — Internet Protocol version 4.
    \item IPv6 — Internet Protocol version 6.
\end{enumerate}

\subsection{6.2 IPv4 addressing}

IPv4 uses a 32-bit address.

It is commonly represented using four decimal numbers separated by dots. This notation is called \textbf{dotted-decimal notation}.

Example:

\texttt{192.168.1.10}

Each decimal section represents 8 bits, called an octet.

An IPv4 address consists of:


\[
32\text{ bits}=4\times 8\text{ bits}
\]


Each octet can represent a value from 0 to 255.

Example:

\begin{table}[htbp]
\centering
\begin{tabular}{c c}
\toprule
\textbf{Decimal} & \textbf{Binary} \\
\midrule
192 & 11000000 \\
168 & 10101000 \\
1 & 00000001 \\
10 & 00001010 \\
\bottomrule
\end{tabular}
\end{table}

Therefore:

\texttt{192.168.1.10}

In binary:

\texttt{11000000.10101000.00000001.00001010}

\subsection{6.3 IPv4 address structure}

An IPv4 address contains two logical parts:

\begin{itemize}
    \item Network portion: Identifies the network.
    \item Host portion: Identifies a host within that network.
\end{itemize}

The division is determined by the subnet mask or prefix length.

Example:

IP address: \texttt{192.168.1.10}

Subnet mask: \texttt{255.255.255.0}

Prefix: \texttt{/24}

This means:

\begin{itemize}
    \item First 24 bits identify the network.
    \item Remaining 8 bits identify the host portion.
\end{itemize}

Network address: \texttt{192.168.1.0}

Host address: \texttt{192.168.1.10}

Broadcast address: \texttt{192.168.1.255}

Usable host range: \texttt{192.168.1.1} to \texttt{192.168.1.254}.

\subsection{6.4 IPv4 address classes}

Historically, IPv4 addresses were divided into five classes.

\begin{table}[htbp]
\centering
\begin{tabularx}{\textwidth}{p{3.0cm} p{3.2cm} X X}
\toprule
\textbf{Class} & \textbf{First octet range} & \textbf{Default subnet mask} & \textbf{Purpose} \\
\midrule
A & 1--126 & 255.0.0.0 (/8) & Large networks \\
B & 128--191 & 255.255.0.0 (/16) & Medium networks \\
C & 192--223 & 255.255.255.0 (/24) & Small networks \\
D & 224--239 & Not applicable & Multicast \\
E & 240--255 & Not applicable & Reserved/experimental \\
\bottomrule
\end{tabularx}
\end{table}

Special cases:

\begin{itemize}
    \item \texttt{0.x.x.x} is reserved for special purposes.
    \item \texttt{127.x.x.x} is reserved for loopback.
    \item \texttt{255.255.255.255} is the limited broadcast address.
\end{itemize}

\textbf{Important:} Classful addressing is largely obsolete. Modern networks use CIDR (Classless Inter-Domain Routing).

Nevertheless, class-based questions frequently appear in traditional networking examinations.

\subsection{6.5 Private and public IPv4 addresses}

\textbf{Public IP address}

A public IP address is globally routable on the Internet, subject to routing and network policies.

\textbf{Private IP address}

Private IP addresses are intended for internal networks and are not globally routed across the public Internet.

Three IPv4 private address ranges are defined by RFC 1918.

\begin{table}[htbp]
\centering
\begin{tabularx}{\textwidth}{p{2.1cm} X X}
\toprule
\textbf{Class} & \textbf{Private address range} & \textbf{CIDR} \\
\midrule
A & 10.0.0.0 -- 10.255.255.255 & 10.0.0.0/8 \\
B & 172.16.0.0 -- 172.31.255.255 & 172.16.0.0/12 \\
C & 192.168.0.0 -- 192.168.255.255 & 192.168.0.0/16 \\
\bottomrule
\end{tabularx}
\end{table}

\textbf{Exam trap:} Not all addresses beginning with \texttt{172} are private.

Only \texttt{172.16.0.0} through \texttt{172.31.255.255} belong to the RFC 1918 private range.

Other important IPv4 addresses:

\begin{table}[htbp]
\centering
\begin{tabularx}{\textwidth}{p{3.9cm} X}
\toprule
\textbf{Address/range} & \textbf{Purpose} \\
\midrule
127.0.0.1 & Loopback address \\
169.254.0.0/16 & Link-local addressing \\
0.0.0.0 & Unspecified address; also used in default route notation \\
255.255.255.255 & Limited broadcast \\
192.0.2.0/24 & Documentation \\
198.51.100.0/24 & Documentation \\
203.0.113.0/24 & Documentation \\
\bottomrule
\end{tabularx}
\end{table}

\subsection{6.6 IPv6 addressing}

IPv6 was developed to address limitations of IPv4, particularly the exhaustion of available IPv4 addresses.

IPv6 uses a 128-bit address.

It is represented using eight groups of hexadecimal digits separated by colons.

Example:

\texttt{2001:0db8:0000:0000:0000:ff00:0042:8329}

Each group contains 16 bits.


\[
128=8\times16
\]


IPv6 uses hexadecimal notation because it makes 128-bit addresses more compact and readable.

\subsection{6.7 IPv6 address abbreviation rules}

IPv6 provides two important abbreviation rules.

\textbf{Rule 1: Leading zeros can be omitted from a group.}

Original:

\texttt{2001:0db8:0000:0000:0000:ff00:0042:8329}

After removing leading zeros:

\texttt{2001:db8:0:0:0:ff00:42:8329}

\textbf{Rule 2: A continuous sequence of zero groups can be replaced with }\texttt{::}\textbf{.}

The previous address becomes:

\texttt{2001:db8::ff00:42:8329}

Important restrictions:

\begin{itemize}
    \item Leading zeros within a group can be removed.
    \item Trailing zeros cannot be removed merely because they are zeros.
    \item \texttt{::} can be used only once in a single IPv6 address because using it more than once would make expansion ambiguous.
\end{itemize}

\subsection{6.8 IPv6 address types}

\begin{quote}
\textbf{Question 1: Unicast}\\[4pt]
\vspace{4pt}
\textbf{Analysis / Answer:} Identifies a single interface. A packet sent to a unicast address is delivered to that interface. Important unicast categories: | Type           | Prefix/range | Purpose                           | | -------------- | ------------ | --------------------------------- | | Global Unicast | \texttt{2000::/3}   | Globally routable addresses       | | Link-local     | \texttt{FE80::/10}  | Communication within a local link | | Unique Local   | \texttt{FC00::/7}   | Private-like internal addressing  | | Loopback       | \texttt{::1/128}    | Refers to the local device        | | Unspecified    | \texttt{::/128}     | Indicates no address              |
\end{quote}
\begin{quote}
\textbf{Question 2: Multicast}\\[4pt]
\vspace{4pt}
\textbf{Analysis / Answer:} Identifies a group of interfaces. IPv6 multicast addresses begin with: \texttt{FF00::/8} A packet sent to a multicast address is delivered to the members of the relevant multicast group.
\end{quote}
\begin{quote}
\textbf{Question 3: Anycast}\\[4pt]
\vspace{4pt}
\textbf{Analysis / Answer:} The same address is assigned to multiple interfaces. A packet is delivered to one suitable interface, generally selected according to routing distance or policy.
\end{quote}
\textbf{Important:} IPv6 does not use broadcast addressing. Multicast provides many of the functions for which IPv4 uses broadcast.

\subsection{6.9 IPv4 vs IPv6 comparison}

\begin{table}[htbp]
\centering
\begin{tabularx}{\textwidth}{p{4.2cm} X X}
\toprule
\textbf{Feature} & \textbf{IPv4} & \textbf{IPv6} \\
\midrule
Version & 4 & 6 \\
Address size & 32 bits & 128 bits \\
Address notation & Decimal & Hexadecimal \\
Example & 192.168.1.10 & 2001:db8::1 \\
Address space & $2^{32}$ & $2^{128}$ \\
Header size & 20--60 bytes & Fixed 40 bytes \\
Checksum & Header checksum & No IP header checksum \\
Broadcast & Supported & Not supported \\
Multicast & Supported & Supported \\
Anycast & Supported through routing arrangements & Built-in addressing model \\
Address configuration & Manual or DHCP & Manual, SLAAC, DHCPv6 \\
Fragmentation & Sender and routers may fragment & Only source hosts fragment \\
NAT & Commonly used & Generally not required for address conservation \\
Security & IPsec support & IPsec support \\
Address resolution & ARP & Neighbor Discovery Protocol \\
TTL field & TTL & Hop Limit \\
Header options & Options & Extension headers \\
\bottomrule
\end{tabularx}
\end{table}

\textbf{Important correction to a common misconception:} IPv6 was designed with IPsec support, but that does not mean every IPv6 connection is automatically encrypted or authenticated. IPv4 can also use IPsec.

\subsection{6.10 IPv4 address space vs IPv6 address space}

IPv4 contains:


\[
2^{32}=4,294,967,296
\]


possible bit patterns.

IPv6 contains:


\[
2^{128}\approx3.4\times10^{38}
\]


possible bit patterns.

IPv6 has an enormously larger address space.

\subsection{6.11 Important IPv6 protocol changes}

\begin{table}[htbp]
\centering
\begin{tabularx}{\textwidth}{X X}
\toprule
\textbf{IPv4} & \textbf{IPv6 equivalent} \\
\midrule
ARP & Neighbor Discovery Protocol (NDP) \\
ICMP & ICMPv6 \\
TTL & Hop Limit \\
DHCP & DHCPv6 and SLAAC \\
Broadcast & Multicast or other suitable mechanisms \\
IPv4 fragmentation by routers & Fragmentation by source host only \\
\bottomrule
\end{tabularx}
\end{table}

\textbf{SLAAC (Stateless Address Autoconfiguration):} Allows an IPv6 host to configure its address using information advertised by routers, without requiring a DHCPv6 server for address assignment.

\textbf{NDP (Neighbor Discovery Protocol):} Uses ICMPv6 messages to discover neighbouring devices, determine link-layer addresses, find routers, and perform address-related checks.

\newpage

\newpage
\begin{center}
\vspace{0.5em}
{\color{primary}\rule{0.45\textwidth}{1.2pt}\\[4pt]}
{\Large\bfseries\color{primary} Section VII: IP Subnetting Mathematics and Binary Calculations\\[2pt]}
{\color{primary}\rule{0.45\textwidth}{1.2pt}}
\vspace{0.5em}
\end{center}

\section{IP Subnetting Mathematics and Binary Calculations}

\subsection{What is subnetting?}
Subnetting is the process of dividing a larger IP network into smaller logical networks called subnets.
It is performed by borrowing bits from the host portion of an IP address and using them as network bits.
For example, consider:
192.168.1.0/24
This network contains 256 total IPv4 addresses.
It can be divided into four smaller subnets:
192.168.1.0/26
192.168.1.64/26
192.168.1.128/26
192.168.1.192/26
Each subnet contains 64 total addresses.
Why is subnetting necessary?
Efficient IP allocation: Avoids allocating unnecessarily large address blocks.
Network segmentation: Separates departments, buildings or groups of devices.
Reduced broadcast domains: Routers separate broadcast domains between subnets.
Security: Allows network access policies to be applied between segments.
Simplified administration: Makes network management easier.
Improved routing organization: Allows route summarization and hierarchical addressing.
\subsection{Basic IPv4 addressing revision}
An IPv4 address contains 32 bits divided into four octets.
Example:
192.168.1.10
Binary representation:
11000000.10101000.00000001.00001010
Each octet contains 8 bits.
The value of each bit is:
\begin{table}[htbp]
\centering
\caption{Binary Positional Weights for an 8-Bit IPv4 Octet}
\label{tab:ipv4_octet_weights}
\begin{tabular}{ccccccccc}
\toprule
\textbf{Bit Position} & 7 & 6 & 5 & 4 & 3 & 2 & 1 & 0 \\
\midrule
\textbf{Positional Power ($2^p$)} & $2^7$ & $2^6$ & $2^5$ & $2^4$ & $2^3$ & $2^2$ & $2^1$ & $2^0$ \\
\textbf{Decimal Value} & 128 & 64 & 32 & 16 & 8 & 4 & 2 & 1 \\
\bottomrule
\end{tabular}
\end{table}
For example:

\[
192=128+64
\]

Therefore:

\[
192=11000000_2
\]

Similarly:

\[
168=128+32+8
\]

Therefore:

\[
168=10101000_2
\]

Important: The network portion and host portion are determined by the subnet mask, not by the visual appearance of the IP address.
\subsection{Subnet mask}
A subnet mask identifies which bits of an IPv4 address belong to the network portion and which belong to the host portion.
Network bits are represented by binary 1.
Host bits are represented by binary 0.
Example:
IP address:
192.168.1.10
Subnet mask:
255.255.255.0
Binary subnet mask:
11111111.11111111.11111111.00000000
There are 24 network bits and 8 host bits.
Therefore, its CIDR notation is:
/24
\begin{table}[htbp]
\centering
\caption{Comprehensive IPv4 Subnet Mask and Prefix Length Reference}
\label{tab:common_subnet_masks}
\begin{tabularx}{\textwidth}{ccccX}
\toprule
\textbf{CIDR Prefix} & \textbf{Dotted-Decimal Subnet Mask} & \textbf{Network Bits} & \textbf{Host Bits} & \textbf{Conventional Role / Capacity} \\
\midrule
/8  & 255.0.0.0         & 8  & 24 & Class A default ($16,777,214$ usable hosts) \\
/16 & 255.255.0.0       & 16 & 16 & Class B default ($65,534$ usable hosts) \\
/17 & 255.255.128.0     & 17 & 15 & $32,766$ usable hosts \\
/18 & 255.255.192.0     & 18 & 14 & $16,382$ usable hosts \\
/19 & 255.255.224.0     & 19 & 13 & $8,190$ usable hosts \\
/20 & 255.255.240.0     & 20 & 12 & $4,094$ usable hosts \\
/21 & 255.255.248.0     & 21 & 11 & $2,046$ usable hosts \\
/22 & 255.255.252.0     & 22 & 10 & $1,022$ usable hosts \\
/23 & 255.255.254.0     & 23 & 9  & $510$ usable hosts \\
/24 & 255.255.255.0     & 24 & 8  & Class C default ($254$ usable hosts) \\
/25 & 255.255.255.128   & 25 & 7  & $126$ usable hosts \\
/26 & 255.255.255.192   & 26 & 6  & $62$ usable hosts \\
/27 & 255.255.255.224   & 27 & 5  & $30$ usable hosts \\
/28 & 255.255.255.240   & 28 & 4  & $14$ usable hosts \\
/29 & 255.255.255.248   & 29 & 3  & $6$ usable hosts \\
/30 & 255.255.255.252   & 30 & 2  & $2$ usable hosts (traditional P2P link) \\
/31 & 255.255.255.254   & 31 & 1  & $2$ usable hosts (RFC 3021 P2P link) \\
/32 & 255.255.255.255   & 32 & 0  & Single host route / Loopback interface \\
\bottomrule
\end{tabularx}
\end{table}
Exam note: /31 is used for point-to-point IPv4 links under RFC 3021, where both addresses can be used as endpoint addresses. The usual network/broadcast reservation does not apply. /32 represents a single IPv4 address or host route.
\subsection{Important subnetting formulas}
These formulas are essential for solving almost every conventional IPv4 subnetting problem.
Let:
$n$ = number of network bits.
$h$ = number of host bits.
$s$ = number of borrowed bits.
$N$ = number of subnets.
$H$ = number of usable host addresses per conventional subnet.
Since IPv4 contains 32 bits:

\[
n+h=32
\]

Formula 1: Total addresses

\[
\boxed{2^h}
\]

Formula 2: Usable host addresses
For conventional subnets where network and broadcast addresses are reserved:

\[
\boxed{2^h-2}
\]

Exceptions include /31 point-to-point links and /32 host routes.
Formula 3: Number of subnets
If $s$ bits are borrowed from the original host portion:

\[
\boxed{2^s}
\]

Older textbooks sometimes use $2^s-2$, based on obsolete subnet-zero and all-ones subnet restrictions. Modern networks permit both, so use $2^s$.
Formula 4: New prefix length

\[
\boxed{\text{New prefix}=\text{Original prefix}+s}
\]

Formula 5: Block size
For the interesting octet:

\[
\boxed{\text{Block size}=256-\text{Subnet mask value}}
\]

For example, for /26:
Subnet mask:
255.255.255.192
Block size:

\[
256-192=64
\]

The subnet ranges in the last octet therefore increase by 64.
\subsection{Understanding subnetting visually}
Consider 192.168.1.0/24.
The original network has 8 host bits.
If we borrow 2 bits:

\[
2^2=4\text{ subnets}
\]

The new prefix becomes:

\[
24+2=26
\]

Remaining host bits:

\[
32-26=6
\]

Total addresses per subnet:

\[
2^6=64
\]

Usable addresses per conventional subnet:

\[
64-2=62
\]

Visual example: Dividing a /24 into four /26 subnets
Original network
/24
Network
24 bits
Host
8 bits
Borrow 2 host bits
Subnet 1
/26
0
Subnet 2
/26
64
Subnet 3
/26
128
Subnet 4
/26
192
Each subnet contains 64 addresses, with 62 conventionally usable host addresses.
\subsection{How to calculate network address, broadcast address and usable host range}
Consider:
IP address: 192.168.1.70/26
Step 1: Determine the subnet mask.
/26 means 26 network bits and 6 host bits.
Subnet mask:
255.255.255.192
Step 2: Calculate the block size.

\[
256-192=64
\]

Step 3: Identify the subnet range.
The subnet boundaries are:
0
64
128
192
The given host value is 70.
Since:

\[
64\leq70<128
\]

The address belongs to the subnet beginning at 64.
Step 4: Calculate network address.
192.168.1.64
Step 5: Calculate broadcast address.
One less than the next subnet boundary:

\[
128-1=127
\]

Broadcast:
192.168.1.127
Step 6: Calculate usable host range.
First usable:
192.168.1.65
Last usable:
192.168.1.126
Final answer:
Parameter
Result
Given IP
192.168.1.70
CIDR
/26
Subnet mask
255.255.255.192
Network address
192.168.1.64
Broadcast address
192.168.1.127
First usable
192.168.1.65
Last usable
192.168.1.126
Total addresses
64
Usable hosts
62
\subsection{Binary AND operation for finding the network address}
Another important method is binary AND.
The network address is obtained by performing a bitwise AND operation between the IP address and the subnet mask.

\[
\boxed{\text{Network Address}=\text{IP Address}\ AND\ \text{Subnet Mask}}
\]

Consider:
IP: 192.168.1.70
Mask: 255.255.255.192
The final octet:
IP address: 01000110 (70) Subnet mask: 11000000 (192) -------- AND result: 01000000 (64)
Therefore:
192.168.1.64
This confirms the network address calculated using the block-size method.
Remember the AND rules:
A
B
A AND B
0
0
0
0
1
0
1
0
0
1
1
1
\subsection{Solved subnetting problems}
Problem 1: Divide a /24 network into 4 equal subnets
Given:
192.168.10.0/24
Required: 4 subnets.
Step 1: Calculate borrowed bits.

\[
2^s=4
\]


\[
s=2
\]

Step 2: Calculate new prefix.

\[
24+2=26
\]

Step 3: Calculate host bits.

\[
32-26=6
\]

Step 4: Calculate addresses per subnet.

\[
2^6=64
\]

Usable hosts:

\[
64-2=62
\]

\begin{table}[htbp]
\centering
\caption{Four Subnets Created from 192.168.10.0/24}
\label{tab:prob1_subnets}
\begin{tabularx}{\textwidth}{ccccc}
\toprule
\textbf{Subnet \#} & \textbf{Network Address} & \textbf{First Usable Host} & \textbf{Last Usable Host} & \textbf{Broadcast Address} \\
\midrule
1 & 192.168.10.0/26   & 192.168.10.1   & 192.168.10.62  & 192.168.10.63 \\
2 & 192.168.10.64/26  & 192.168.10.65  & 192.168.10.126 & 192.168.10.127 \\
3 & 192.168.10.128/26 & 192.168.10.129 & 192.168.10.190 & 192.168.10.191 \\
4 & 192.168.10.192/26 & 192.168.10.193 & 192.168.10.254 & 192.168.10.255 \\
\bottomrule
\end{tabularx}
\end{table}
Each subnet supports 62 conventionally usable hosts.
Problem 2: How many hosts can a /27 network support?
Given:
192.168.5.0/27
Host bits:

\[
h=32-27=5
\]

Total addresses:

\[
2^5=32
\]

Usable hosts:

\[
32-2=\boxed{30}
\]

Subnet mask:
255.255.255.224
Block size:

\[
256-224=32
\]

Subnet boundaries:
192.168.5.0
192.168.5.32
192.168.5.64
192.168.5.96
192.168.5.128
192.168.5.160
192.168.5.192
192.168.5.224
Problem 3: Find the subnet information for 172.16.35.123/20
Given:
IP: 172.16.35.123
Prefix: /20
Step 1: Determine subnet mask.
/20 means:
11111111.11111111.11110000.00000000
Subnet mask:
255.255.240.0
Step 2: Identify the interesting octet.
The third octet is 240.
Block size:

\[
256-240=16
\]

Step 3: Find the subnet range in the third octet.
Subnet boundaries:
0, 16, 32, 48, 64, ...
The third octet is 35, so it belongs to the range 32--47.
Step 4: Calculate network and broadcast addresses.
Network:
172.16.32.0
Broadcast:
172.16.47.255
Step 5: Calculate usable host range.
First usable:
172.16.32.1
Last usable:
172.16.47.254
Step 6: Calculate host capacity.

\[
h=32-20=12
\]


\[
2^{12}-2=4094
\]

Final result
Parameter
Value
IP address
172.16.35.123
CIDR
/20
Subnet mask
255.255.240.0
Network
172.16.32.0
Broadcast
172.16.47.255
First host
172.16.32.1
Last host
172.16.47.254
Usable hosts
4,094
Problem 4: How many subnets and hosts are created by changing /24 to /28?
Original prefix: /24
New prefix: /28
Borrowed bits:

\[
s=28-24=4
\]

Number of subnets:

\[
2^4=\boxed{16}
\]

Remaining host bits:

\[
h=32-28=4
\]

Usable hosts per subnet:

\[
2^4-2=\boxed{14}
\]

Subnet mask:
255.255.255.240
Block size:

\[
256-240=16
\]

Each subnet contains 16 addresses and 14 conventionally usable host addresses.
Problem 5: What subnet mask is required for at least 100 hosts?
We need the smallest number of host bits satisfying:

\[
2^h-2\geq100
\]

Try $h=6$:

\[
2^6-2=62
\]

Insufficient.
Try $h=7$:

\[
2^7-2=126
\]

Sufficient.
Therefore:

\[
\text{Prefix}=32-7=\boxed{/25}
\]

Subnet mask:
255.255.255.128
Maximum conventional usable hosts: 126.
Problem 6: What is the smallest subnet for 500 hosts?
We need:

\[
2^h-2\geq500
\]

For $h=8$:

\[
2^8-2=254
\]

Insufficient.
For $h=9$:

\[
2^9-2=510
\]

Sufficient.
Therefore:

\[
32-9=\boxed{/23}
\]

Subnet mask:
255.255.254.0
A /23 subnet supports 510 conventionally usable hosts.

\newpage

\newpage
\begin{center}
\vspace{0.5em}
{\color{primary}\rule{0.45\textwidth}{1.2pt}\\[4pt]}
{\Large\bfseries\color{primary} Section VIII: Classless Addressing (CIDR), VLSM, and Supernetting\\[2pt]}
{\color{primary}\rule{0.45\textwidth}{1.2pt}}
\vspace{0.5em}
\end{center}

\section{Classless Addressing (CIDR), VLSM, and Supernetting}

\subsection{What is CIDR?}
CIDR stands for Classless Inter-Domain Routing.
It is an IP addressing and routing method that replaces the older classful addressing system.
CIDR allows networks to use variable-length network prefixes rather than being restricted to Class A, B, or C default masks.
It is represented using slash notation.
Example:
192.168.1.0/24
Here:
192.168.1.0 is the network address.
/24 indicates that the first 24 bits identify the network.
The remaining 8 bits are available for host addressing.
Advantages of CIDR
Efficient IP address allocation.
Reduced wastage of IPv4 addresses.
Supports variable-length prefixes.
Enables route aggregation.
Reduces the size of routing tables.
Supports hierarchical network allocation.
\subsection{Classful addressing vs CIDR}
\begin{table}[htbp]
\centering
\caption{Comparison: Classful Addressing vs. Classless Inter-Domain Routing (CIDR)}
\label{tab:classful_vs_cidr}
\begin{tabularx}{\textwidth}{p{3.8cm} X X}
\toprule
\textbf{Evaluation Feature} & \textbf{Classful Addressing} & \textbf{Classless Routing (CIDR)} \\
\midrule
\textbf{Network Partitioning} & Rigid fixed classes (A, B, C) & Variable-length prefix (/1 to /32) \\
\textbf{Subnet Mask} & Default boundary masks (/8, /16, /24) & Arbitrary custom mask length \\
\textbf{Address Efficiency} & Low (massive address wastage) & High (tailored address blocks) \\
\textbf{Route Aggregation} & Restricted to class boundaries & Fully supported supernetting \\
\textbf{Routing Scalability} & Massive, flat routing tables & Compact, hierarchical routing tables \\
\textbf{Modern Status} & Obsolete legacy system & Ubiquitous global standard \\
\bottomrule
\end{tabularx}
\end{table}
Example:
Under traditional classful addressing, a Class C network has a default /24 prefix.
With CIDR, an organization can receive a /27 or /22 network depending on its requirements.
\subsection{CIDR numerical calculations}
Consider:
192.168.20.0/22
Host bits:

\[
32-22=10
\]

Total addresses:

\[
2^{10}=1024
\]

Conventional usable hosts:

\[
1024-2=\boxed{1022}
\]

Subnet mask:
255.255.252.0
Block size in the third octet:

\[
256-252=4
\]

The network beginning at 192.168.20.0/22 covers:
Network: 192.168.20.0
Broadcast: 192.168.23.255
First host: 192.168.20.1
Last host: 192.168.23.254
\subsection{CIDR Quick Reference Table}
\begin{table}[htbp]
\centering
\caption{CIDR Quick Reference for Subnets /24 to /32}
\label{tab:cidr_quick_ref}
\begin{tabular}{ccc}
\toprule
\textbf{CIDR Prefix} & \textbf{Total IPv4 Addresses} & \textbf{Conventional Usable Hosts} \\
\midrule
/24 & 256 & 254 \\
/25 & 128 & 126 \\
/26 & 64  & 62  \\
/27 & 32  & 30  \\
/28 & 16  & 14  \\
/29 & 8   & 6   \\
/30 & 4   & 2   \\
/31 & 2   & 2 (RFC 3021 Point-to-Point) \\
/32 & 1   & 1 (Host Route) \\
\bottomrule
\end{tabular}
\end{table}
\subsection{Longest Prefix Match in CIDR routing}
CIDR allows multiple network routes to overlap.
A router chooses the most specific matching route, meaning the route with the longest matching prefix.
Example routing table:
Destination
Next hop
10.0.0.0/8
Router A
10.1.0.0/16
Router B
10.1.2.0/24
Router C
0.0.0.0/0
Router D
A packet destined for:
10.1.2.50
Matches all four entries.
The router chooses:
10.1.2.0/24
because /24 is the longest, most specific matching prefix.
Exam point: The default route /0 is the least specific IPv4 route.
3. VLSM (Variable Length Subnet Masking)
\subsection{What is VLSM?}
VLSM stands for Variable Length Subnet Masking.
It is a subnetting technique that allows different subnets within the same allocated network to use different subnet masks.
In conventional fixed-length subnetting, every subnet has the same prefix length.
In VLSM, the prefix length is selected according to the number of hosts required by each subnet.
Why is VLSM necessary?
Consider an organization with three departments:
Administration: 100 hosts.
Accounts: 50 hosts.
IT: 20 hosts.
If we use the same /25 subnet for every department, each department receives 126 usable addresses.
This wastes a substantial number of addresses.
With VLSM, we can allocate:
Administration: /25 $\to$ 126 usable addresses.
Accounts: /26 $\to$ 62 usable addresses.
IT: /27 $\to$ 30 usable addresses.
This reduces address wastage.
\subsection{VLSM calculation procedure}
Follow these steps when solving VLSM numerical problems.
List all subnet requirements.
Sort them in descending order of host requirements.
Determine the smallest suitable subnet for each requirement.
Allocate the largest subnet first.
Assign the next subnet from the next available aligned network address.
Calculate network, usable host range and broadcast addresses.
Verify that no subnet overlaps another.
Important formula
For each subnet, find the smallest $h$ satisfying:

\[
\boxed{2^h-2\geq\text{Required hosts}}
\]

Then:

\[
\boxed{\text{Prefix}=32-h}
\]

\subsection{Solved VLSM problem}
Question:
An organization receives the network:
192.168.10.0/24
It needs to create four subnets:
Department
Required hosts
Administration
100
Accounts
50
IT
25
Management
10
Design the network using VLSM.
Step 1: Sort by descending host requirement
The requirements are already in descending order:
Administration: 100.
Accounts: 50.
IT: 25.
Management: 10.
Step 2: Determine the required prefix for each subnet
Administration: 100 hosts

\[
2^7-2=126
\]

Prefix:

\[
32-7=/25
\]

Accounts: 50 hosts

\[
2^6-2=62
\]

Prefix:

\[
32-6=/26
\]

IT: 25 hosts

\[
2^5-2=30
\]

Prefix:

\[
32-5=/27
\]

Management: 10 hosts

\[
2^4-2=14
\]

Prefix:

\[
32-4=/28
\]

Step 3: Allocate subnets sequentially
Start with:
192.168.10.0/24
Allocate the largest subnet first.
Subnet 1: Administration --- /25
A /25 contains 128 addresses.
Network:
192.168.10.0/25
Broadcast:
192.168.10.127
Usable range:
192.168.10.1 -- 192.168.10.126
Next available address:
192.168.10.128
Subnet 2: Accounts --- /26
A /26 contains 64 addresses.
Network:
192.168.10.128/26
Broadcast:
192.168.10.191
Usable range:
192.168.10.129 -- 192.168.10.190
Next available address:
192.168.10.192
Subnet 3: IT --- /27
A /27 contains 32 addresses.
Network:
192.168.10.192/27
Broadcast:
192.168.10.223
Usable range:
192.168.10.193 -- 192.168.10.222
Next available address:
192.168.10.224
Subnet 4: Management --- /28
A /28 contains 16 addresses.
Network:
192.168.10.224/28
Broadcast:
192.168.10.239
Usable range:
192.168.10.225 -- 192.168.10.238
\subsection{Final VLSM Allocation Table}
\begin{table}[htbp]
\centering
\caption{VLSM Network Allocation for 192.168.10.0/24}
\label{tab:vlsm_allocation_result}
\begin{tabularx}{\textwidth}{p{2.8cm} c p{3.2cm} c c}
\toprule
\textbf{Department} & \textbf{Required} & \textbf{Assigned Subnet} & \textbf{Broadcast} & \textbf{Usable Hosts} \\
\midrule
Administration & 100 & 192.168.10.0/25   & 192.168.10.127 & 126 ($128 - 2$) \\
Accounts       & 50  & 192.168.10.128/26 & 192.168.10.191 & 62 ($64 - 2$) \\
IT             & 25  & 192.168.10.192/27 & 192.168.10.223 & 30 ($32 - 2$) \\
Management     & 10  & 192.168.10.224/28 & 192.168.10.239 & 14 ($16 - 2$) \\
\midrule
\textbf{Unallocated} & -- & 192.168.10.240/28 & 192.168.10.255 & 16 addresses reserved \\
\bottomrule
\end{tabularx}
\end{table}
Step 5: Verify the allocation
Total allocated addresses:

\[
128+64+32+16=240
\]

Original /24 network:

\[
2^8=256\text{ addresses}
\]

Remaining:

\[
256-240=\boxed{16}
\]

All four subnets are correctly aligned, non-overlapping and contained within the original /24 block.
\subsection{FLSM vs. VLSM Comparison}
\begin{table}[htbp]
\centering
\caption{Fixed-Length (FLSM) vs. Variable-Length (VLSM) Subnet Masking}
\label{tab:flsm_vs_vlsm}
\begin{tabularx}{\textwidth}{p{3.8cm} X X}
\toprule
\textbf{Feature} & \textbf{FLSM} & \textbf{VLSM} \\
\midrule
\textbf{Subnet Mask} & Identical prefix across all subnets & Tailored prefix per subnet requirement \\
\textbf{Host Capacity} & Uniform across all subnets & Customized to host counts per segment \\
\textbf{Address Efficiency} & Substantial wastage on small links & High utilization; minimal waste \\
\textbf{Hierarchical Subnetting} & Not supported & Subnetting within subnets \\
\textbf{Design Complexity} & Simple, straightforward calculation & Requires sorted, careful allocation \\
\bottomrule
\end{tabularx}
\end{table}
Exam point: VLSM allows subnetting a subnet, whereas FLSM uses the same mask for all subnets.
4. Supernetting and Route Aggregation
\subsection{What is supernetting?}
Supernetting is the process of combining multiple contiguous smaller networks into one larger network by reducing the network prefix length.
It is the opposite of subnetting.
Example:
Four /24 networks:
192.168.0.0/24
192.168.1.0/24
192.168.2.0/24
192.168.3.0/24
Can be aggregated into:
192.168.0.0/22
The resulting /22 represents all four networks.
Why is supernetting used?
Reduces routing table entries.
Simplifies routing.
Improves routing scalability.
Supports hierarchical IP address allocation.
Reduces the amount of routing information that must be advertised.
\subsection{Supernetting requirements}
For conventional CIDR route aggregation, the component networks should:
Be contiguous.
Have the same prefix length.
Have a count that is a power of two.
Be aligned on the correct aggregate boundary.
For example, four /24 networks require a /22 aggregate because:

\[
24-\log_2 4=22
\]

Their first network must be aligned to a multiple of four in the third octet.
Formula
If $k$ equal-sized networks are combined:

\[
\boxed{\text{Aggregate prefix}=p-\log_2(k)}
\]

where:
$p$ is the original prefix length.
$k$ is the number of networks being combined.
$k$ must be a power of two.
\subsection{Solved supernetting problem}
Question: Aggregate these four networks:
192.168.8.0/24
192.168.9.0/24
192.168.10.0/24
192.168.11.0/24
Step 1: Count the networks.

\[
k=4
\]

Step 2: Calculate the new prefix.

\[
24-\log_2(4)
\]


\[
24-2=\boxed{/22}
\]

Step 3: Determine the aggregate address.
The networks begin at third-octet value 8 and cover 8, 9, 10 and 11.
Since the block size is four, 8 is correctly aligned.
Aggregate network:
192.168.8.0/22
Step 4: Determine the range.
Network: 192.168.8.0
Broadcast: 192.168.11.255
Total addresses: $2^{10}=1024$
Conventional usable hosts: $1022$
\subsection{Subnetting vs. Supernetting Comparison}
\begin{table}[htbp]
\centering
\caption{Comparison: Subnetting versus Supernetting}
\label{tab:subnetting_vs_supernetting}
\begin{tabularx}{\textwidth}{p{3.8cm} X X}
\toprule
\textbf{Parameter} & \textbf{Subnetting} & \textbf{Supernetting (Aggregation)} \\
\midrule
\textbf{Fundamental Operation} & Divides one network into smaller subnets & Aggregates contiguous networks into one \\
\textbf{Prefix Length ($p$)} & Increases ($p_{\text{new}} = p + s$) & Decreases ($p_{\text{new}} = p - \log_2 k$) \\
\textbf{Host Capacity ($2^h$)} & Decreases per subnet & Increases in the aggregate block \\
\textbf{Primary Objective} & Network segmentation, security, broadcast containment & Route table reduction, routing scalability \\
\textbf{Example} & One /24 divided into four /26s & Four /24s summarized into one /22 \\
\bottomrule
\end{tabularx}
\end{table}
Important: Route aggregation must not accidentally include address space that the organization does not control or that should be routed differently. An overly broad aggregate can cause traffic to be sent to the wrong destination.

\newpage

\newpage
\begin{center}
\vspace{0.5em}
{\color{primary}\rule{0.45\textwidth}{1.2pt}\\[4pt]}
{\Large\bfseries\color{primary} Section IX: Analog and Digital Modulation Techniques\\[2pt]}
{\color{primary}\rule{0.45\textwidth}{1.2pt}}
\vspace{0.5em}
\end{center}

\section{Analog and Digital Modulation Techniques}

\subsection{What is modulation?}
Modulation is the process of modifying one or more properties of a high-frequency carrier signal according to a message or information signal.
The carrier signal is used to transport information through a communication medium.
In simple terms, modulation allows information to be transmitted by changing characteristics of a carrier wave.
The three fundamental properties of a sinusoidal carrier signal are:
Amplitude.
Frequency.
Phase.
A general sinusoidal carrier signal can be represented as:

\[
c(t)=A_c\sin(2\pi f_ct+\phi)
\]

Where:
$A_c$ = Carrier amplitude.
$f_c$ = Carrier frequency.
$t$ = Time.
$\phi$ = Carrier phase.
Modulation changes one or more of these parameters according to the information signal.
\subsection{Why is modulation necessary?}
Modulation is used for several reasons:
Long-distance communication: It allows information to be transmitted through communication channels designed for particular frequency ranges.
Efficient transmission: It helps adapt signals to the characteristics of a transmission medium.
Multiplexing: Different carrier frequencies can carry different signals simultaneously.
Antenna practicality: Higher-frequency carrier signals can make practical antenna dimensions more manageable.
Wireless communication: Radio communication systems depend on modulation to transmit information.
Interference management: Frequency allocation and modulation techniques help organize communication channels.
Basic communication system


A simplified communication system consists of:
Information source --- generates the message.
Modulator --- combines the message with a carrier.
Transmitter --- prepares and sends the modulated signal.
Communication channel --- carries the signal.
Receiver --- receives the transmitted signal.
Demodulator --- extracts the original information.
Destination --- receives the recovered message.
Demodulation is the reverse process of modulation: extracting the original information from the modulated carrier.
\subsection{Types of modulation}
Modulation can be broadly classified into:
Analog modulation.
Digital modulation.
Pulse modulation.
A. Analog modulation
Analog modulation is used when the information signal is continuous.
The carrier's amplitude, frequency or phase is varied according to the analog message signal.
Three fundamental techniques are:
Amplitude Modulation (AM).
Frequency Modulation (FM).
Phase Modulation (PM).
1. Amplitude Modulation (AM)
In AM, the amplitude of the carrier signal changes according to the information signal, while carrier frequency and phase remain constant in the idealized model.


Characteristics:
Carrier amplitude varies.
Carrier frequency remains constant.
Relatively simple implementation.
More susceptible to amplitude noise.
Traditional AM radio broadcasting uses this technique.
AM modulation index:

\[
\boxed{m=\frac{A_m}{A_c}}
\]

Where:
$A_m$ = Peak amplitude of the message signal.
$A_c$ = Amplitude of the carrier.
Percentage modulation:

\[
\boxed{\text{Modulation percentage}=m\times100\%}
\]

For conventional AM without overmodulation:

\[
0\leq m\leq1
\]

Example:
Carrier amplitude = 10 V.
Message amplitude = 5 V.

\[
m=\frac{5}{10}=0.5
\]


\[
\text{Modulation percentage}=0.5\times100=50\%
\]

2. Frequency Modulation (FM)
In FM, the frequency of the carrier signal changes according to the information signal, while its amplitude remains constant in the idealized model.


Characteristics:
Carrier frequency varies.
Carrier amplitude remains constant.
Generally more resistant to amplitude noise than conventional AM.
Requires greater bandwidth than conventional narrowband AM in many practical configurations.
Used in FM radio broadcasting.
Frequency deviation:
Frequency deviation is the maximum difference between the instantaneous carrier frequency and its unmodulated carrier frequency.

\[
\Delta f=|f_{\text{instantaneous}}-f_c|_{\max}
\]

FM modulation index:

\[
\boxed{\beta=\frac{\Delta f}{f_m}}
\]

Where:
$\Delta f$ = Peak frequency deviation.
$f_m$ = Message signal frequency.
Example:
Frequency deviation = 75 kHz.
Message frequency = 15 kHz.

\[
\beta=\frac{75}{15}=\boxed{5}
\]

Carson's Rule estimates the bandwidth of an FM signal:

\[
\boxed{BW\approx2(\Delta f+f_m)}
\]

Using the previous values:

\[
BW=2(75+15)
\]


\[
=\boxed{180\text{ kHz}}
\]

3. Phase Modulation (PM)
In PM, the phase of the carrier signal changes according to the information signal.


Characteristics:
Carrier phase varies.
Carrier amplitude remains constant in the idealized model.
Used in digital communication systems as the basis for PSK.
Closely related mathematically to FM.
Important distinction:
AM changes amplitude.
FM changes frequency.
PM changes phase.
\begin{table}[htbp]
\centering
\caption{Comparison of Analog Modulation Techniques}
\label{tab:analog_modulation_comp}
\begin{tabularx}{\textwidth}{p{3.2cm} X X X}
\toprule
\textbf{Feature} & \textbf{Amplitude Modulation (AM)} & \textbf{Frequency Modulation (FM)} & \textbf{Phase Modulation (PM)} \\
\midrule
\textbf{Modulated Parameter} & Carrier amplitude $A_c$ & Carrier frequency $f_c$ & Carrier phase $\phi$ \\
\textbf{Carrier Amplitude} & Varies with message signal & Constant & Constant \\
\textbf{Carrier Frequency} & Constant & Varies with message signal & Instantaneous frequency varies \\
\textbf{Noise Immunity} & Susceptible to electrical noise & High immunity to amplitude noise & High noise immunity \\
\textbf{Bandwidth Requirement} & Narrowband ($2 f_m$) & Wideband ($2(\Delta f + f_m)$ Carson) & Proportional to phase deviation \\
\textbf{Canonical Usage} & Long-wave / medium-wave radio & High-fidelity FM broadcasting & Digital phase modulation basis \\
\bottomrule
\end{tabularx}
\end{table}
\subsection{Digital modulation}
Digital modulation transmits digital information by changing the characteristics of a carrier signal according to binary data or symbols.
The three fundamental digital modulation techniques are:
Amplitude Shift Keying (ASK).
Frequency Shift Keying (FSK).
Phase Shift Keying (PSK).
A fourth important technique is Quadrature Amplitude Modulation (QAM).
A. Amplitude Shift Keying (ASK)
In ASK, the carrier amplitude changes according to the digital information.
For basic binary ASK:
Binary 1: Carrier with a particular amplitude.
Binary 0: Carrier with zero or a different amplitude.


Characteristics:
Simple implementation.
Uses different amplitude levels.
Sensitive to amplitude noise.
Used in some low-complexity communication systems.
A common special case is On-Off Keying (OOK).
B. Frequency Shift Keying (FSK)
In FSK, the carrier frequency changes according to the digital information.
For binary FSK:
Binary 1: Frequency $f_1$.
Binary 0: Frequency $f_0$.


Characteristics:
Uses different carrier frequencies.
Information is represented through frequency changes.
Generally less sensitive to amplitude noise than ASK.
Used in some radio telemetry and low-speed communication systems.
C. Phase Shift Keying (PSK)
In PSK, the carrier phase changes according to the digital information.
In Binary PSK (BPSK):
Binary 0: One carrier phase.
Binary 1: A carrier phase shifted by 180°.


Characteristics:
Information is represented through phase changes.
BPSK uses two phase states.
More robust than basic ASK in many noisy-channel conditions.
Widely used in digital communication systems.
D. Quadrature Amplitude Modulation (QAM)
QAM combines amplitude and phase modulation.
It changes both the amplitude and phase of a carrier to represent digital symbols.
It is widely used in:
Wi-Fi.
Cable modem systems.
Digital television.
Broadband communication.
Examples:
16-QAM.
64-QAM.
256-QAM.
1024-QAM.


Bits per symbol in digital modulation
If a modulation scheme uses $M$ distinct symbols, then:

\[
\boxed{\text{Bits per symbol}=\log_2 M}
\]

\begin{table}[htbp]
\centering
\caption{Digital Modulation Schemes, Symbol Constellations, and Spectral Efficiency}
\label{tab:digital_modulation_symbols}
\begin{tabular}{ccc}
\toprule
\textbf{Modulation Scheme} & \textbf{Number of Constellation Symbols ($M$)} & \textbf{Spectral Efficiency ($\text{Bits/Symbol} = \log_2 M$)} \\
\midrule
BPSK & 2 & 1 bit/symbol \\
QPSK & 4 & 2 bits/symbol \\
8-PSK & 8 & 3 bits/symbol \\
16-QAM & 16 & 4 bits/symbol \\
64-QAM & 64 & 6 bits/symbol \\
256-QAM & 256 & 8 bits/symbol \\
1024-QAM & 1024 & 10 bits/symbol \\
\bottomrule
\end{tabular}
\end{table}
Example:
How many bits per symbol are transmitted by 64-QAM?

\[
\log_2 64=6
\]

Therefore:

\[
\boxed{6\text{ bits/symbol}}
\]

\begin{table}[htbp]
\centering
\caption{Master Comparison of Digital Keying Techniques}
\label{tab:digital_keying_comp}
\begin{tabularx}{\textwidth}{p{2.8cm} p{3.2cm} X}
\toprule
\textbf{Keying Technique} & \textbf{Modulated Characteristic} & \textbf{Primary Operational Features} \\
\midrule
\textbf{ASK} (Amplitude) & Carrier amplitude levels & Simple implementation; highly sensitive to amplitude noise and fading \\
\textbf{FSK} (Frequency) & Discrete frequency tones & Constant envelope; robust against non-linear amplification and noise \\
\textbf{PSK} (Phase) & Distinct phase shifts (e.g. 180$^\circ$) & Constant envelope; superior noise immunity; widely used in modern wireless \\
\textbf{QAM} (Quadrature) & Amplitude and phase combined & High spectral efficiency; standard in Wi-Fi (802.11ac/ax), cable modems, 5G \\
\bottomrule
\end{tabularx}
\end{table}
\subsection{Pulse modulation}
Pulse modulation represents information using pulses.
It can be classified into analog pulse modulation and digital pulse modulation.
A. Pulse Amplitude Modulation (PAM)
The amplitude of equally spaced pulses varies according to the message signal.
B. Pulse Width Modulation (PWM)
The width or duration of pulses changes according to the information signal.
C. Pulse Position Modulation (PPM)
The position of each pulse changes according to the information signal.
D. Pulse Code Modulation (PCM)
PCM converts an analog signal into a digital representation through three fundamental stages:
Sampling.
Quantization.
Encoding.


Sampling: Measures the analog signal at regular intervals.
Quantization: Maps sampled values to discrete amplitude levels.
Encoding: Converts quantized values into binary codes.
Nyquist sampling theorem
The sampling frequency must be at least twice the highest frequency present in a band-limited signal for ideal reconstruction:

\[
\boxed{f_s\geq2f_{\max}}
\]

Where:
$f_s$ = Sampling frequency.
$f_{\max}$ = Highest message frequency.
Example:
An analog signal has a maximum frequency of 4 kHz.
Minimum theoretical sampling frequency:

\[
f_s=2\times4=8\text{ kHz}
\]

Therefore, the minimum Nyquist sampling rate is:

\[
\boxed{8000\text{ samples/second}}
\]

In practical systems, sampling above the theoretical minimum is often used to allow filtering and implementation margins.
PCM bit rate
For PCM:

\[
\boxed{\text{Bit rate}=f_s\times n}
\]

Where:
$f_s$ = Sampling rate.
$n$ = Bits per sample.
Example:
A voice signal is sampled at 8,000 samples/second using 8 bits per sample.

\[
R=8000\times8
\]


\[
=\boxed{64,000\text{ bps}=64\text{ kbps}}
\]

\subsection{Bit rate vs Baud rate}
This is a particularly important numerical topic.
Bit rate: Number of bits transmitted per second.
Baud rate: Number of symbols transmitted per second.
Formula:

\[
\boxed{\text{Bit rate}=\text{Baud rate}\times\text{Bits per symbol}}
\]

For an $M$-ary modulation scheme:

\[
\boxed{R_b=R_s\log_2 M}
\]

Where:
$R_b$ = Bit rate.
$R_s$ = Symbol rate or baud rate.
$M$ = Number of distinct symbols.
Solved example 1
A communication system operates at 2,400 baud using 16-QAM.
Find its bit rate.
Bits per symbol:

\[
\log_2 16=4
\]

Bit rate:

\[
2400\times4
\]


\[
=\boxed{9600\text{ bps}}
\]

Solved example 2
A system transmits at 48,000 bps using 64-QAM.
Find its baud rate.
Bits per symbol:

\[
\log_2 64=6
\]

Baud rate:

\[
\frac{48000}{6}
\]


\[
=\boxed{8000\text{ baud}}
\]

Exam trap: Bit rate and baud rate are not always equal. They are equal for binary modulation with one bit per symbol, such as BPSK.

\newpage

\newpage
\begin{center}
\vspace{0.5em}
{\color{primary}\rule{0.45\textwidth}{1.2pt}\\[4pt]}
{\Large\bfseries\color{primary} Section X: Data Transmission Modes, Methods, and Multiplexing\\[2pt]}
{\color{primary}\rule{0.45\textwidth}{1.2pt}}
\vspace{0.5em}
\end{center}

\section{Data Transmission Modes, Methods, and Multiplexing}

Transmission in computer networks can be classified according to the direction of communication, the way bits are transmitted, and how multiple devices share a communication medium.
These classifications are related but describe different aspects of communication.
\subsection{Transmission modes based on direction}
There are three fundamental communication modes.
\subsubsection{A. Simplex mode}
Simplex communication allows data to travel in only one direction.
One device acts as the sender, while another acts as the receiver.
Examples:
Traditional television broadcasting.
Traditional radio broadcasting.
A keyboard sending input to a computer (as a simplified example).
Characteristics:
One-way communication.
No reverse data transmission through the same communication arrangement.
Simple implementation.
\subsubsection{B. Half-duplex mode}
Half-duplex communication allows data transmission in both directions, but only one direction at a time.
Example: Walkie-talkie communication.
One person speaks while the other listens. They must take turns.
Characteristics:
Bidirectional communication.
Cannot transmit simultaneously in both directions.
Shared communication capacity.
\subsubsection{C. Full-duplex mode}
Full-duplex communication allows simultaneous transmission in both directions.
Example:
Telephone conversation.
Modern switched Ethernet links.
Characteristics:
Bidirectional communication.
Simultaneous sending and receiving.
Better utilization for interactive communication.
\begin{table}[htbp]
\centering
\caption{Comparison of Directional Transmission Modes}
\label{tab:transmission_modes}
\begin{tabularx}{\textwidth}{p{3.2cm} c c X}
\toprule
\textbf{Mode} & \textbf{Direction} & \textbf{Simultaneous?} & \textbf{Canonical Example} \\
\midrule
\textbf{Simplex} & Unidirectional & No & Television / radio broadcast, keyboard to CPU \\
\textbf{Half-Duplex} & Bidirectional & No (alternating) & Walkie-talkie, shared legacy hub Ethernet \\
\textbf{Full-Duplex} & Bidirectional & Yes (simultaneous) & Modern switched Ethernet, telephone call \\
\bottomrule
\end{tabularx}
\end{table}
Important examination point: Traditional shared Ethernet using hubs used half-duplex operation. Modern switched Ethernet links commonly support full-duplex operation.
\subsection{Serial vs Parallel transmission}
Another classification is based on how bits are transmitted.
\subsubsection{A. Serial transmission}
Serial transmission sends data one bit at a time over a communication channel.
Example:
To transmit the binary sequence:
10110010
The bits are transmitted sequentially.
Characteristics:
Uses fewer data lines.
Suitable for longer-distance communication.
Common in modern communication interfaces.
Easier to maintain signal integrity at high speeds over many links.
Examples:
USB.
SATA.
Ethernet.
Serial communication interfaces.
\subsubsection{B. Parallel transmission}
Parallel transmission sends multiple bits simultaneously over multiple communication lines.
For example, an 8-bit data value can be transmitted using eight data lines simultaneously.
Characteristics:
Multiple bits are sent at once.
Requires multiple communication lines.
Can experience timing skew between signals.
Commonly suitable for short-distance internal communication.
Examples:
Traditional parallel printer interfaces.
Internal parallel data buses.
\begin{table}[htbp]
\centering
\caption{Comparison: Serial versus Parallel Data Transmission}
\label{tab:serial_vs_parallel}
\begin{tabularx}{\textwidth}{p{3.8cm} X X}
\toprule
\textbf{Feature} & \textbf{Serial Transmission} & \textbf{Parallel Transmission} \\
\midrule
\textbf{Bit Transfer Rate} & One bit sequentially per clock cycle & Multiple bits ($8, 16, 32$) simultaneously \\
\textbf{Physical Lines} & Single data line (or differential pair) & Multiple parallel physical wires \\
\textbf{Clock Skew} & Negligible across single line & Severe at high frequencies / long distances \\
\textbf{Distance Suitability} & Long distances and high speeds & Short distances (internal CPU buses) \\
\textbf{Modern Standards} & PCIe, SATA, USB, Ethernet & Legacy ATA/IDE, IEEE 1284 parallel port \\
\bottomrule
\end{tabularx}
\end{table}
Exam trap: Serial transmission is not necessarily slower than parallel transmission. Modern high-speed serial interfaces can achieve very high data rates.
\subsection{Synchronous vs Asynchronous transmission}
This classification concerns how the sender and receiver coordinate the timing of data transmission.
\subsubsection{A. Asynchronous transmission}
In asynchronous transmission, data is transmitted in units, commonly characters, with start and stop bits used to identify their boundaries.
Each character can be transmitted independently, without requiring a continuous shared clock for the entire stream.
Example: Traditional UART serial communication.
Characteristics:
Start and stop bits are used.
No continuous shared clock is required.
Suitable for intermittent data transmission.
Has additional framing overhead.
Example of a simplified 8N1 UART frame:
Component
Bits
Start bit
1
Data
8
Parity
0
Stop bit
1
Total
10
The notation 8N1 means:
8 data bits.
No parity bit.
1 stop bit.
Thus, transmitting 8 data bits requires 10 transmitted bits.
Efficiency:

\[
\text{Efficiency}=\frac{8}{10}\times100
\]


\[
=\boxed{80\%}
\]

\subsubsection{B. Synchronous transmission}
In synchronous transmission, the sender and receiver coordinate timing so that data can be transmitted in a continuous stream or in larger blocks.
Characteristics:
Timing synchronization is maintained.
Data can be sent in frames or blocks.
Less per-character framing overhead.
Suitable for high-speed communication.
Examples:
Synchronous Ethernet-related communication.
Synchronous serial protocols.
Synchronous communication links.
\begin{table}[htbp]
\centering
\caption{Comparison: Synchronous versus Asynchronous Transmission}
\label{tab:sync_vs_async}
\begin{tabularx}{\textwidth}{p{3.8cm} X X}
\toprule
\textbf{Feature} & \textbf{Synchronous Transmission} & \textbf{Asynchronous Transmission} \\
\midrule
\textbf{Clocking} & Shared continuous clock signal & Independent clocks; synchronized per unit \\
\textbf{Data Framing} & Continuous byte stream or block frames & Start bit, 7--8 data bits, parity, stop bit \\
\textbf{Protocol Overhead} & Minimal preamble/postamble framing & High overhead ($\ge 20\%$ framing bits) \\
\textbf{Data Flow} & Steady, continuous high-speed data & Bursty, intermittent character transfer \\
\textbf{Example Protocols} & Synchronous Ethernet, SONET/SDH & UART, RS-232 serial communication \\
\bottomrule
\end{tabularx}
\end{table}
\subsection{Transmission methods based on signal representation}
\subsubsection{A. Analog transmission}
Analog transmission uses continuously varying signals to represent information.
Characteristics:
Signal amplitude, frequency or phase varies continuously.
Susceptible to noise and distortion.
Used in traditional analog communication systems.
Examples:
Traditional analog radio.
Traditional analog telephone systems.
\subsubsection{B. Digital transmission}
Digital transmission represents information using discrete signal levels, commonly binary 0 and 1.
Characteristics:
Suitable for computer data.
Supports digital encoding and regeneration.
Common in modern communication systems.
Examples:
Ethernet.
Digital fiber-optic communication.
USB.
Important: Digital data can be transmitted over analog communication channels using modulation techniques. Therefore, analog transmission media and digital data representation are not mutually exclusive concepts.
\subsection{Baseband vs Broadband transmission}
Baseband transmission
Baseband transmission uses the available channel bandwidth to carry a signal directly in its baseband form.
Examples:
Traditional Ethernet LAN transmission.
Digital signals transmitted over a dedicated communication link.
Broadband transmission
Broadband transmission uses modulation and frequency allocation to carry signals over one or more frequency channels.
It can support multiple communication channels over the same medium.
Examples:
Cable television networks.
Cable Internet access.
Broadband radio systems.
Feature
Baseband
Broadband
Signal
Baseband signal
Modulated frequency channels
Channel usage
Typically uses the channel for one baseband transmission
Can carry multiple frequency channels
Multiplexing
Depends on system
Commonly frequency-division based
Example
Traditional Ethernet
Cable broadband
\subsection{Multiplexing}
Multiplexing is a technique used to combine multiple signals or data streams and transmit them through a shared communication medium.
The device performing this operation is called a multiplexer (MUX).
At the receiving end, a demultiplexer (DEMUX) separates the signals.
Basic multiplexing concept
Signal 1
Signal 2
Signal 3
Multiplexer (MUX)
Shared communication channel
Demultiplexer (DEMUX)
Signal 1
Signal 2
Signal 3
Types of multiplexing
1. Frequency Division Multiplexing (FDM)
The available bandwidth is divided into separate frequency bands.
Each signal uses a different frequency range.
Example: Traditional radio broadcasting, where different radio stations operate at different frequencies.
2. Time Division Multiplexing (TDM)
Multiple signals share the same communication channel by using different time slots.
Example: A communication link allocates different time slots to different data streams.
Types:
Synchronous TDM: Fixed time slots are allocated.
Statistical TDM: Time slots are allocated dynamically according to demand.
3. Wavelength Division Multiplexing (WDM)
WDM is used in optical fiber communication.
Different optical wavelengths (colours of light) carry separate data channels through the same fiber.
Types include:
CWDM: Coarse Wavelength Division Multiplexing.
DWDM: Dense Wavelength Division Multiplexing.
4. Code Division Multiplexing (CDM/CDMA)
Multiple users share the same frequency band at the same time using different spreading codes.
Example: Code Division Multiple Access in certain cellular communication systems.
\begin{table}[htbp]
\centering
\caption{Master Comparison of Multiplexing Techniques}
\label{tab:multiplexing_techniques}
\begin{tabularx}{\textwidth}{p{3.2cm} p{3.2cm} X}
\toprule
\textbf{Technique} & \textbf{Separation Domain} & \textbf{Primary Application Domain} \\
\midrule
\textbf{FDM} (Frequency) & Separate frequency channels & Radio broadcast, cable TV, DSL broadband \\
\textbf{TDM} (Time) & Non-overlapping time slots & T1/E1 digital carrier systems, cellular \\
\textbf{WDM} (Wavelength) & Optical laser wavelengths & High-capacity fiber-optic backbones (DWDM) \\
\textbf{CDM / CDMA} (Code) & Orthogonal mathematical codes & 3G cellular, GPS satellite communication \\
\bottomrule
\end{tabularx}
\end{table}
\subsection{Switching methods}
Switching determines how communication data travels through a network.
Three traditional switching methods are:
1. Circuit switching
A dedicated communication path is established before communication begins.
Example: Traditional telephone networks.
2. Packet switching
Data is divided into packets. Packets are forwarded through network devices using routing and forwarding decisions.
Example: The Internet.
Two common packet-switching approaches:
Datagram packet switching: Each packet is routed independently.
Virtual circuit packet switching: A logical path is established, and packets follow that logical connection.
3. Message switching
The complete message is received and stored at an intermediate node before being forwarded to the next node.
This is called store-and-forward communication.
\begin{table}[htbp]
\centering
\caption{Systematic Comparison of Network Switching Methodologies}
\label{tab:switching_methodologies}
\begin{tabularx}{\textwidth}{p{3.5cm} X X X}
\toprule
\textbf{Feature} & \textbf{Circuit Switching} & \textbf{Packet Switching} & \textbf{Message Switching} \\
\midrule
\textbf{Dedicated Path} & Yes (end-to-end setup) & No (independent packets) & No (store-and-forward) \\
\textbf{Transmission Unit} & Continuous byte/bit stream & Discrete packets & Complete message file \\
\textbf{Intermediate Storage} & None (direct circuit) & Temporary packet buffers & Full message hard-drive \\
\textbf{Resource Utilization} & Low (reserved bandwidth) & High (statistical sharing) & Moderate to high \\
\textbf{Latency Profile} & Setup delay, constant transfer & Variable queuing delay & High store-and-forward delay \\
\textbf{Example} & Traditional PSTN telephone & Internet (IPv4/IPv6) & Historical telegraph systems \\
\bottomrule
\end{tabularx}
\end{table}

\newpage

\newpage
\begin{center}
\vspace{0.5em}
{\color{primary}\rule{0.45\textwidth}{1.2pt}\\[4pt]}
{\Large\bfseries\color{primary} Section XI: Network Topologies and Mathematical Link Analysis\\[2pt]}
{\color{primary}\rule{0.45\textwidth}{1.2pt}}
\vspace{0.5em}
\end{center}

\section{Network Topologies and Mathematical Link Analysis}

\subsection{What is network topology?}
A network topology is the physical or logical arrangement of nodes, links and communication paths within a computer network.
Physical topology describes how devices and cables are physically connected.
Logical topology describes how data flows between devices, regardless of the physical arrangement.
For example, a network may physically use a star arrangement while its communication technology determines the logical data-access behaviour.
\subsection{Bus topology}


In bus topology, all devices share a common backbone cable.
Advantages:
Simple installation.
Relatively low cable requirement.
Low initial cost for small networks.
Disadvantages:
Backbone failure can stop communication for all devices.
Fault isolation can be difficult.
Performance can degrade as traffic increases.
Collisions are possible in shared-medium implementations.
Exam point: Traditional Ethernet standards such as 10BASE2 and 10BASE5 used bus-style physical arrangements.
\subsection{Star topology}


In star topology, every device connects to a central device, commonly a switch.
Advantages:
Easy to install and manage.
Failure of one cable usually affects only one device.
Easy fault detection.
Easy to expand.
Good performance with switched Ethernet.
Disadvantages:
Failure of the central switch can disrupt the entire star.
Requires more cable than a traditional bus topology.
Central equipment adds cost.
Exam point: Star topology is one of the most common physical topologies in modern LANs.
\subsection{Ring topology}


In ring topology, each node connects to two neighbouring nodes, forming a closed loop.
Advantages:
Predictable transmission in token-based designs.
Can provide orderly access to the shared communication medium.
Avoids collisions in traditional token-passing implementations.
Disadvantages:
A failure can disrupt the ring unless redundancy is provided.
Adding or removing nodes may affect the network.
Troubleshooting can be difficult.
Example: IBM Token Ring is a historical example.
\subsection{Mesh topology}


In mesh topology, nodes are connected through multiple paths.
There are two major types.
Full mesh: Every node connects directly to every other node.
Partial mesh: Only selected nodes have multiple direct connections.
For a full mesh with $n$ devices:

\[
\boxed{L=\frac{n(n-1)}{2}}
\]

where $L$ is the number of links.
For example, with 6 devices:

\[
L=\frac{6(6-1)}{2}
\]


\[
=\frac{30}{2}=\boxed{15}
\]

Advantages:
High redundancy.
Multiple communication paths.
Better fault tolerance.
Can improve availability.
Disadvantages:
High installation cost.
Complex configuration.
Requires many links and interfaces.
\subsection{Tree topology}


Tree topology is a hierarchical arrangement in which several star networks connect to higher-level devices.
Advantages:
Highly scalable.
Supports hierarchical network management.
Easier segmentation.
Suitable for large organizations.
Disadvantages:
Failure of an upper-level device can affect multiple lower-level networks.
More complex configuration.
Requires careful planning.
Example: A university network with a core switch, building distribution switches and laboratory access switches.
\subsection{Hybrid topology}


Hybrid topology combines two or more different network topologies.
For example:
Star-bus topology.
Star-ring topology.
Star-mesh topology.
Advantages:
Flexible design.
Scalable.
Can combine the strengths of multiple topologies.
Suitable for complex networks.
Disadvantages:
More expensive.
Requires more planning.
Complex troubleshooting and maintenance.
\subsection{Point-to-point topology}


A point-to-point topology establishes a direct communication link between two devices.
Examples:
A dedicated leased line between two offices.
A direct router-to-router link.
A direct communication link between two computers.
Advantages:
Simple architecture.
Dedicated link capacity.
Easier troubleshooting.
Disadvantages:
Connects only two endpoints directly.
Not suitable as a complete architecture for a large network without additional links.
\subsection{Master Topology Comparison}
\begin{table}[htbp]
\centering
\caption{Systematic Comparative Analysis of Network Topologies}
\label{tab:topology_comparisons}
\begin{tabularx}{\textwidth}{p{3.0cm} X X}
\toprule
\textbf{Topology} & \textbf{Primary Advantage} & \textbf{Critical Limitation / Failure Point} \\
\midrule
\textbf{Bus} & Inexpensive, simple cabling & Central backbone cable break halts entire network \\
\textbf{Star} & Fault isolation, easy expansion & Central switch represents single point of failure \\
\textbf{Ring} & Deterministic token-passing access & Cable or node failure breaks loop circulation \\
\textbf{Full Mesh} & Maximum fault tolerance ($L = n(n-1)/2$) & Prohibitive cost and cabling complexity \\
\textbf{Tree} & Scalable hierarchical administration & Failure of root/distribution switch isolates branches \\
\textbf{Hybrid} & Highly adaptable to large campuses & Complex configuration and troubleshooting \\
\textbf{Point-to-Point} & Dedicated, uncontended link bandwidth & Connects only two endpoints directly \\
\bottomrule
\end{tabularx}
\end{table}
\subsection{Numerical questions related to topology}
Question 1: How many links are required in a full mesh network of 8 devices?

\[
L=\frac{n(n-1)}{2}
\]


\[
L=\frac{8(8-1)}{2}
\]


\[
=\frac{56}{2}=\boxed{28}
\]

Question 2: How many links are required in a full mesh network of 10 devices?

\[
L=\frac{10(9)}{2}=\boxed{45}
\]

Question 3: How many additional links are required when a full mesh network grows from 5 to 7 devices?
For 5 devices:

\[
L_1=\frac{5(4)}{2}=10
\]

For 7 devices:

\[
L_2=\frac{7(6)}{2}=21
\]

Additional links:

\[
21-10=\boxed{11}
\]


\newpage

\newpage
\begin{center}
\vspace{0.5em}
{\color{primary}\rule{0.45\textwidth}{1.2pt}\\[4pt]}
{\Large\bfseries\color{primary} Section XII: Communication Latency, Performance Metrics, and BDP\\[2pt]}
{\color{primary}\rule{0.45\textwidth}{1.2pt}}
\vspace{0.5em}
\end{center}

\section{Communication Latency, Performance Metrics, and BDP}

\subsection{What is network latency?}
Network latency is the time taken for data to travel from a source to a destination across a network.
It is commonly measured in milliseconds (ms).
For example, if a computer sends a request to a server and receives a response after 40 milliseconds, the measured round-trip time (RTT) is approximately 40 ms.
Lower latency generally means faster responsiveness.
Applications that are particularly sensitive to latency include:
Online gaming.
Video conferencing.
Voice over IP (VoIP).
Remote desktop systems.
Financial trading systems.
Real-time monitoring and control.
\subsection{Types of network delay}
Network latency is composed of several delay components.
The four fundamental components are:
Processing delay.
Queuing delay.
Transmission delay.
Propagation delay.

\[
\boxed{D_{\text{total}}=D_{\text{processing}}+D_{\text{queuing}}+D_{\text{transmission}}+D_{\text{propagation}}}
\]

For a path containing multiple routers and links, the delays accumulate across the path.
\subsubsection{A. Processing delay}
Processing delay is the time required by a network device to examine and process an incoming packet.
Activities include:
Checking packet headers.
Validating packet information.
Looking up routing table entries.
Performing forwarding decisions.
Checking certain error-detection fields.
Example: A router receives a packet and takes 0.2 ms to examine its header and determine the next hop.
Processing delay:

\[
D_{\text{processing}}=0.2\text{ ms}
\]

\subsubsection{B. Queuing delay}
Queuing delay is the time a packet waits in a buffer before being transmitted.
It depends on:
Network traffic volume.
Link utilization.
Packet arrival rate.
Available transmission capacity.
Queue management policies.
Example: If a packet waits 2 ms in a router's output queue, its queuing delay is 2 ms.
During congestion, queuing delay can increase significantly.
\subsubsection{C. Transmission delay}
Transmission delay is the time required to place all the bits of a packet onto a communication link.
It depends on:
Packet size.
Link bandwidth.
Formula:

\[
\boxed{D_{\text{transmission}}=\frac{\text{Packet size in bits}}{\text{Transmission rate in bits/second}}}
\]

Example:
A packet contains 12,000 bits and the link bandwidth is 6 Mbps.

\[
D=\frac{12000}{6\times10^6}
\]


\[
=0.002\text{ seconds}
\]


\[
=\boxed{2\text{ ms}}
\]

Important: Transmission delay is not the same as propagation delay.
\subsubsection{D. Propagation delay}
Propagation delay is the time required for a signal to travel from one end of a communication link to the other.
It depends on:
Physical distance.
Signal propagation speed in the medium.
Formula:

\[
\boxed{D_{\text{propagation}}=\frac{\text{Distance}}{\text{Propagation speed}}}
\]

Example:
A signal travels through a fiber-optic link of 1,000 km.
Assume its propagation speed is:

\[
2\times10^8\text{ m/s}
\]

Distance:

\[
1000\text{ km}=1,000,000\text{ m}
\]

Propagation delay:

\[
D=\frac{1,000,000}{2\times10^8}
\]


\[
=0.005\text{ seconds}
\]


\[
=\boxed{5\text{ ms}}
\]

\subsection{Transmission Delay vs. Propagation Delay}
\begin{table}[htbp]
\centering
\caption{Comprehensive Comparison: Transmission Delay vs. Propagation Delay}
\label{tab:trans_vs_prop_delay}
\begin{tabularx}{\textwidth}{p{3.8cm} X X}
\toprule
\textbf{Parameter} & \textbf{Transmission Delay ($D_{\text{trans}}$)} & \textbf{Propagation Delay ($D_{\text{prop}}$)} \\
\midrule
\textbf{Physical Meaning} & Time to push all packet bits onto medium & Time for electromagnetic signal to cross link \\
\textbf{Governing Formula} & $D_t = L / R$ & $D_p = d / v$ \\
\textbf{Key Determinants} & Packet size ($L$) and link bandwidth ($R$) & Physical distance ($d$) and propagation speed ($v$) \\
\textbf{Speed of Medium} & Independent of propagation speed & Directly inversely proportional to $v$ ($\approx 2\times 10^8$ m/s) \\
\textbf{Distance Dependency} & Completely independent of distance & Directly proportional to physical distance \\
\textbf{Optimization Path} & Increase link bandwidth or compress packet & Shorten physical path or upgrade medium \\
\bottomrule
\end{tabularx}
\end{table}
Where:
$L$ = packet length in bits.
$R$ = transmission rate in bits per second.
$d$ = physical distance.
$v$ = propagation speed.
\subsection{Solved numerical problem: Total network delay}
Question:
A packet of size 10,000 bits travels through a link with bandwidth 5 Mbps. The link is 500 km long, and the signal propagation speed is $2\times10^8$ m/s.
Assume:
Processing delay = 0.5 ms.
Queuing delay = 1.5 ms.
Calculate the total delay.
Step 1: Calculate transmission delay.

\[
D_t=\frac{L}{R}
\]


\[
=\frac{10000}{5\times10^6}
\]


\[
=0.002\text{ seconds}=2\text{ ms}
\]

Step 2: Calculate propagation delay.
Distance:

\[
500\text{ km}=500000\text{ m}
\]


\[
D_p=\frac{500000}{2\times10^8}
\]


\[
=0.0025\text{ seconds}=2.5\text{ ms}
\]

Step 3: Calculate total delay.

\[
D_{\text{total}}=0.5+1.5+2+2.5
\]


\[
\boxed{D_{\text{total}}=6.5\text{ ms}}
\]

\subsection{Store-and-forward vs cut-through latency}
These switching methods affect how long a switch takes before forwarding a frame.
Store-and-forward:
The switch receives the entire frame before forwarding it.
If a frame has length $L$ bits and the incoming link rate is $R$ bits/s, the complete-frame reception time is:

\[
D=\frac{L}{R}
\]

The switch then performs its forwarding operations.
Cut-through:
The switch begins forwarding after reading enough of the frame header to identify the outgoing port.
It can reduce forwarding latency, especially for large frames.
However, it may forward a frame before receiving the complete frame and checking its integrity.
\subsection{What is Round-Trip Time (RTT)?}
Round-Trip Time is the total time required for a packet to travel from the sender to the receiver and for a corresponding response or acknowledgment to return.

\[
\boxed{RTT=D_{\text{forward}}+D_{\text{return}}}
\]

If the forward and return paths have equal delay:

\[
RTT=2D
\]

Example:
Forward delay = 20 ms.
Return delay = 25 ms.

\[
RTT=20+25=\boxed{45\text{ ms}}
\]

RTT is commonly measured using the ping command.
Example:
ping google.com
Typical output:
64 bytes from ...: icmp\_seq=1 ttl=117 time=32.4 ms
The time value indicates the measured round-trip time for that ICMP echo exchange.
Important: RTT is not necessarily twice the one-way delay because forward and return routes may be different.
\subsection{What is jitter?}
Jitter is the variation in packet delay over time.
Suppose packets arrive at these times after transmission:
Packet
Delay
1
20 ms
2
22 ms
3
45 ms
4
21 ms
5
24 ms
The delays are inconsistent, especially for Packet 3.
This variation is jitter.
Jitter is particularly important in:
Voice calls.
Video conferencing.
Live streaming.
Online gaming.
High jitter can cause:
Audio interruptions.
Distorted voice.
Video playback irregularities.
Unstable real-time communication.
A jitter buffer can temporarily store packets and play them at a more consistent rate, at the cost of some additional latency.
Exam distinction:
Latency measures delay.
Jitter measures variation in delay.
Packet loss measures packets that fail to reach their destination.
\subsection{Bandwidth vs Throughput vs Goodput}
These three terms are often confused.
Bandwidth
Bandwidth represents the maximum data-carrying capacity of a communication link.
It is commonly measured in:
bps (bits per second).
Kbps.
Mbps.
Gbps.
Example: A network connection with 100 Mbps bandwidth.
Throughput
Throughput is the actual rate at which data is successfully transferred over a network during a period.
Example:
A connection has 100 Mbps bandwidth, but the actual measured throughput is 75 Mbps.
Goodput
Goodput is the rate of useful application data delivered successfully, excluding protocol overhead and retransmitted data.
Example:
A network transfers 75 Mbps in total, but 68 Mbps represents useful application payload.
Term
Meaning
Bandwidth
Maximum or configured capacity
Throughput
Actual successful data transfer rate
Goodput
Useful application data transfer rate
Relationship in a typical transfer:

\[
\text{Goodput}\leq\text{Throughput}\leq\text{Available capacity}
\]

This is a simplified relationship; measured throughput and available capacity can differ depending on the measurement method and network conditions.
\subsection{Bandwidth-delay product (BDP)}
The Bandwidth-Delay Product represents the amount of data that can be in transit on a link at a given time.
Formula:

\[
\boxed{BDP=\text{Bandwidth}\times\text{RTT}}
\]

Suppose:
Bandwidth = 100 Mbps.
RTT = 50 ms.
Convert:

\[
50\text{ ms}=0.05\text{ s}
\]

Then:

\[
BDP=100\times10^6\times0.05
\]


\[
=5,000,000\text{ bits}
\]

Convert into bytes:

\[
\frac{5,000,000}{8}=625,000\text{ bytes}
\]

Therefore:

\[
\boxed{BDP=625\text{ KB approximately}}
\]

This indicates the amount of data that can fill the link while waiting for acknowledgments to return.
Practical significance: TCP window sizing and network performance tuning can depend heavily on the bandwidth-delay product.
\subsection{Factors affecting network latency}
Factor
Effect
Physical distance
Longer distance generally increases propagation delay
Bandwidth
Higher bandwidth reduces transmission delay
Network congestion
Increases queuing delay
Router processing
Adds processing delay
Number of network hops
Adds processing and per-link delays
Packet size
Larger packets increase transmission delay
Wireless interference
Can cause retransmissions and variable delays
Server processing
Adds application response time
Routing path
An inefficient route can increase delay
\subsection{Latency-related numerical formulas}
Concept
Formula
Transmission delay
$L/R$
Propagation delay
$d/v$
Total delay
$D_p+D_q+D_t+D_{\text{prop}}$
Round-trip time
Forward delay + Return delay
Bandwidth-delay product
Bandwidth $\times$ RTT
Data transfer time (idealized)
Data size / Transfer rate

\newpage

\newpage
\begin{center}
\vspace{0.5em}
{\color{primary}\rule{0.45\textwidth}{1.2pt}\\[4pt]}
{\Large\bfseries\color{primary} Section XIII: IP Routing Tables and Longest Prefix Matching (LPM)\\[2pt]}
{\color{primary}\rule{0.45\textwidth}{1.2pt}}
\vspace{0.5em}
\end{center}

\section{IP Routing Tables and Longest Prefix Matching (LPM)}

\subsection{What is Longest Prefix Match?}
Longest Prefix Match (LPM) is the routing technique used by IP routers to select the most specific matching route from their routing table.
When a router receives an IP packet, it examines the destination IP address and compares it against available routing table entries.
If multiple routes match the destination address, the router selects the route with the longest matching network prefix.
In other words:
The route with the greatest number of matching leading network bits is selected.
A longer prefix represents a smaller and more specific network.
For example:
/8 is less specific than /16.
/16 is less specific than /24.
/24 is less specific than /28.
\subsection{Why is LPM necessary?}
\begin{table}[htbp]
\centering
\caption{Routing Table Example for Longest Prefix Matching}
\label{tab:lpm_routing_sample}
\begin{tabularx}{0.85\textwidth}{ccX}
\toprule
\textbf{Destination Network Prefix} & \textbf{Prefix Length} & \textbf{Next-Hop Forwarding Action} \\
\midrule
192.168.0.0/16 & /16 (16 bits) & Forward to Router A \\
192.168.1.0/24 & /24 (24 bits) & Forward to Router B \\
192.168.1.128/25 & /25 (25 bits) & Forward to Router C (Most Specific Match) \\
0.0.0.0/0 & /0 (0 bits) & Gateway of Last Resort \\
\bottomrule
\end{tabularx}
\end{table}
Suppose a packet arrives with destination:
192.168.1.150
This destination matches all three routes.
However, the router selects:
192.168.1.128/25
because /25 is the most specific matching prefix.
The packet is therefore forwarded to Router C.
\subsection{How does LPM work?}
A router generally performs the following operations:
Reads the destination IP address from the packet header.
Examines the routing table.
Checks which network prefixes match the destination.
Selects the matching route with the longest prefix.
Uses the selected route's next hop and outgoing interface to forward the packet.
If no more specific route exists, the router may use a default route.
The default IPv4 route is:
0.0.0.0/0
It matches every IPv4 destination, but is less specific than any route with a nonzero prefix length.
\subsection{Solved LPM numerical problems}
\subsubsection*{Problem 1: Select the correct route}
A router has the following routing table:
Network
Next hop
10.0.0.0/8
A
10.1.0.0/16
B
10.1.2.0/24
C
10.1.2.128/25
D
0.0.0.0/0
E
Find the next hop for:
10.1.2.200
Step 1: Check the matching routes.
The destination 10.1.2.200 belongs to:
10.0.0.0/8
10.1.0.0/16
10.1.2.0/24
10.1.2.128/25
0.0.0.0/0
Step 2: Compare prefix lengths.

\[
8,\ 16,\ 24,\ 25,\ 0
\]

Step 3: Select the longest prefix.

\[
\boxed{/25}
\]

Therefore, the selected route is:
10.1.2.128/25
Next hop:

\[
\boxed{\text{Router D}}
\]

\subsubsection*{Problem 2: Find the matching route for 172.16.35.10}
Routing table:
Network
Next hop
172.16.0.0/16
A
172.16.32.0/20
B
172.16.35.0/24
C
172.16.35.8/29
D
0.0.0.0/0
E
Destination:
172.16.35.10
The destination belongs to:
172.16.0.0/16
172.16.32.0/20
172.16.35.0/24
172.16.35.8/29
The most specific matching route is:
172.16.35.8/29
Therefore:

\[
\boxed{\text{Next hop D}}
\]

\subsubsection*{Problem 3: What happens if no specific route matches?}
Routing table:
Network
Next hop
192.168.1.0/24
A
172.16.0.0/16
B
0.0.0.0/0
C
Destination:
8.8.8.8
The destination does not match either of the first two routes.
It matches the default route:
0.0.0.0/0
Therefore, the packet is forwarded to:

\[
\boxed{\text{Router C}}
\]

If there is no matching route and no applicable default route, the router cannot forward the packet and will typically discard it.
\subsection{LPM and subnet masks}
LPM can also be understood using a bitwise operation.
A destination matches a route if:

\[
\boxed{(D\ AND\ M)=N}
\]

Where:
$D$ = Destination IP address.
$M$ = Route subnet mask.
$N$ = Network address.
For example:
Destination:
192.168.1.150
Route:
192.168.1.128/25
Mask:
255.255.255.128
Last octet:
Destination: 10010110 (150) Mask: 10000000 (128) -------- AND: 10000000 (128)
The result is 128, matching the route's network address.
Therefore, the destination matches this route.
\subsection{LPM vs. Administrative Distance vs. Metric}
\begin{table}[htbp]
\centering
\caption{Decision Hierarchy: Longest Prefix Match vs. Administrative Distance vs. Metric}
\label{tab:lpm_vs_ad_metric}
\begin{tabularx}{\textwidth}{p{3.8cm} X X}
\toprule
\textbf{Routing Concept} & \textbf{Functional Scope / Decision Stage} & \textbf{Arbitration Criteria} \\
\midrule
\textbf{Longest Prefix Match} & Packet forwarding stage & Selects the most specific route prefix length in forwarding table \\
\textbf{Administrative Distance} & Routing table installation stage & Compares trustworthiness between different routing sources for same prefix \\
\textbf{Routing Metric} & Protocol-specific route calculation & Compares paths within the same protocol (cost, hops, bandwidth, delay) \\
\bottomrule
\end{tabularx}
\end{table}

\newpage

\newpage
\begin{center}
\vspace{0.5em}
{\color{primary}\rule{0.45\textwidth}{1.2pt}\\[4pt]}
{\Large\bfseries\color{primary} Section XIV: Dynamic Routing Protocols and Loop Prevention: STP, OSPF, and BGP\\[2pt]}
{\color{primary}\rule{0.45\textwidth}{1.2pt}}
\vspace{0.5em}
\end{center}

\section{Dynamic Routing Protocols and Loop Prevention: STP, OSPF, and BGP}



These are three important networking protocols, but they solve different problems.


\begin{table}[htbp]
\centering
\caption{Protocol Classification: STP vs. OSPF vs. BGP}
\label{tab:stp_ospf_bgp_summary}
\begin{tabularx}{\textwidth}{p{2.2cm} p{3.8cm} X p{3.2cm}}
\toprule
\textbf{Protocol} & \textbf{Full Designation} & \textbf{Primary Operational Purpose} & \textbf{OSI Layer / Area} \\
\midrule
\textbf{STP} & Spanning Tree Protocol & Prevents Layer 2 Ethernet switching loops & Data Link Layer (L2) \\
\textbf{OSPF} & Open Shortest Path First & Interior routing within an Autonomous System & Network Layer (L3 IGP) \\
\textbf{BGP} & Border Gateway Protocol & Exterior routing between Autonomous Systems & Network Layer (L3 EGP) \\
\bottomrule
\end{tabularx}
\end{table}

\subsection{Spanning Tree Protocol (STP)}

\subsubsection{What is STP?}

Spanning Tree Protocol (STP) is a Layer 2 network protocol that prevents switching loops in Ethernet networks containing redundant links.

Redundant links are useful because they provide alternative paths if a connection fails. However, they can create switching loops if all links forward traffic simultaneously.

STP solves this problem by logically disabling selected redundant ports while retaining them as backup paths.

\textbf{Standard:} IEEE 802.1D.

\subsubsection{Why are switching loops dangerous?}

Consider three switches connected in a triangle:

\begin{itemize}
\item Switch A connects to Switch B.
\item Switch B connects to Switch C.
\item Switch C connects back to Switch A.
\end{itemize}

If all links forward frames, an Ethernet broadcast frame can circulate repeatedly.

This can cause:

\begin{enumerate}
\item \textbf{Broadcast storm:} Broadcast frames continuously circulate and consume bandwidth.
\item \textbf{MAC address table instability:} Switches may learn the same source MAC address on different ports repeatedly.
\item \textbf{Duplicate frames:} Devices may receive multiple copies of the same frame.
\item \textbf{Network congestion:} Excessive traffic can make legitimate communication extremely slow or unavailable.
\end{enumerate}

Unlike IP packets, traditional Ethernet frames do not have a TTL field that automatically limits their circulation.

\subsubsection{How STP works}

STP creates a loop-free logical topology by selecting one switch as the root bridge and calculating the best path from other switches to it.

Its basic operation involves four steps.

\textbf{Step 1: Elect the Root Bridge}

Every switch has a Bridge ID (BID), consisting of:

\begin{itemize}
\item Bridge priority.
\item MAC address.
\end{itemize}

The switch with the lowest Bridge ID becomes the root bridge.

The default priority in traditional STP is 32768, with the MAC address breaking ties.

\textbf{Step 2: Select Root Ports}

Every non-root switch selects one Root Port.

The Root Port is the port providing the lowest-cost path toward the root bridge.

The root bridge itself has no Root Port.

\textbf{Step 3: Select Designated Ports}

Each network segment has one Designated Port.

The Designated Port is responsible for forwarding traffic from that segment toward the root bridge.

On the root bridge, all operational ports connected to network segments are normally designated ports.

\textbf{Step 4: Block Redundant Paths}

Ports that are not selected as Root Ports or Designated Ports are prevented from forwarding ordinary data traffic in the stable topology.

This breaks Layer 2 loops while retaining redundant physical connections.

\subsubsection{STP Port Roles}
\begin{table}[htbp]
\centering
\caption{Spanning Tree Protocol (STP) Port Roles}
\label{tab:stp_port_roles}
\begin{tabularx}{\textwidth}{p{3.5cm} X}
\toprule
\textbf{Port Role} & \textbf{Functional Responsibilities and Selection Rule} \\
\midrule
\textbf{Root Port (RP)} & Lowest-cost path from a non-root switch toward the elected Root Bridge; exactly one per non-root bridge \\
\textbf{Designated Port (DP)} & Lowest-cost port forwarding traffic from a local network segment toward the Root Bridge; one per segment \\
\textbf{Alternate Port (AP)} & Backup port providing an alternative path to the Root Bridge (RSTP / 802.1w); blocks data traffic \\
\textbf{Backup Port (BP)} & Redundant backup port for a shared collision domain segment; blocks ordinary data traffic \\
\bottomrule
\end{tabularx}
\end{table}


\subsubsection{STP Port States}
\begin{table}[htbp]
\centering
\caption{IEEE 802.1D STP Port Operational States}
\label{tab:stp_port_states}
\begin{tabularx}{\textwidth}{p{2.8cm} X}
\toprule
\textbf{Port State} & \textbf{Operational Traffic and MAC Learning Behavior} \\
\midrule
\textbf{Disabled} & Administratively shut down; does not participate in STP, forward frames, or learn MACs \\
\textbf{Blocking} & Discards ordinary data frames; processes received BPDUs; does not populate MAC table \\
\textbf{Listening} & Processes and transmits BPDUs; determines port roles; does not learn MACs or forward data \\
\textbf{Learning} & Populates MAC address table from incoming frames; does not yet forward ordinary data frames \\
\textbf{Forwarding} & Fully operational; sends/receives data frames and actively updates MAC address table \\
\bottomrule
\end{tabularx}
\end{table}

Traditional 802.1D STP uses five port states:


A port generally transitions through:

`Blocking $\to$ Listening $\to$ Learning $\to$ Forwarding`

The Disabled state is not part of the normal transition sequence.

\subsubsection{What are BPDUs?}

Bridge Protocol Data Units (BPDUs) are control messages exchanged between switches to share spanning-tree information.

They help switches:

\begin{itemize}
\item Elect the root bridge.
\item Calculate path costs.
\item Determine port roles.
\item Detect topology changes.
\end{itemize}

\subsubsection{STP Variants and Standards}
\begin{table}[htbp]
\centering
\caption{Spanning Tree Protocol Standards and Variants}
\label{tab:stp_variants}
\begin{tabularx}{\textwidth}{p{3.0cm} p{2.8cm} X}
\toprule
\textbf{Protocol Variant} & \textbf{Standard / RFC} & \textbf{Key Characteristics and Capabilities} \\
\midrule
\textbf{STP} & IEEE 802.1D & Legacy standard; timer-based convergence (30--50 seconds); single instance \\
\textbf{RSTP} & IEEE 802.1w & Rapid convergence ($<1$ second) via proposal-agreement handshakes; 3 port states \\
\textbf{MSTP} & IEEE 802.1s & Maps multiple VLANs into distinct spanning-tree instances; highly scalable \\
\textbf{PVST+} & Cisco Proprietary & Maintains an independent spanning-tree instance for each active VLAN \\
\bottomrule
\end{tabularx}
\end{table}


\textbf{Important examination distinction:} RSTP uses three principal port states: Discarding, Learning, and Forwarding. It also introduces alternate and backup port roles.

\subsubsection{STP numerical example}

Suppose a switch has three possible paths to the root bridge:


The switch selects Path B as its Root Port path because it has the lowest total path cost.

If two paths have equal cost, additional tie-breaking rules are used, including upstream Bridge ID and Port ID.

\subsubsection{STP versus Rapid STP (RSTP)}
\begin{table}[htbp]
\centering
\caption{Comparative Analysis: Classic STP (802.1D) versus Rapid STP (802.1w)}
\label{tab:stp_vs_rstp}
\begin{tabularx}{\textwidth}{p{3.5cm} X X}
\toprule
\textbf{Evaluation Dimension} & \textbf{Classic STP (IEEE 802.1D)} & \textbf{Rapid STP (IEEE 802.1w)} \\
\midrule
\textbf{Convergence Mechanism} & Timer-driven (Listening/Learning timers: 30--50 s) & Handshake-driven (Explicit proposal/agreement: sub-second) \\
\textbf{Number of Port States} & Five (Disabled, Blocking, Listening, Learning, Forwarding) & Three (Discarding, Learning, Forwarding) \\
\textbf{Port Roles} & Root, Designated, Blocked & Root, Designated, Alternate, Backup \\
\textbf{Network Suitability} & Legacy networks & Modern switched campus/enterprise LANs \\
\bottomrule
\end{tabularx}
\end{table}


\subsubsection{Important STP commands (Cisco IOS)}

\begin{lstlisting}
show spanning-tree
show spanning-tree vlan 10
show spanning-tree root
show spanning-tree blockedports
\end{lstlisting}

To configure a switch as the preferred root bridge:

\begin{lstlisting}
configure terminal
spanning-tree vlan 10 priority 4096
\end{lstlisting}

Alternatively:

\begin{lstlisting}
spanning-tree vlan 10 root primary
\end{lstlisting}

The exact behavior of the second command depends on the switch's current configuration and platform.

\subsection{OSPF --- Open Shortest Path First}

\subsubsection{What is OSPF?}

OSPF is a link-state Interior Gateway Protocol (IGP) used to exchange routing information within a single Autonomous System (AS).

An autonomous system is a collection of IP networks and routers administered under a common routing policy.

OSPF calculates the shortest paths using Dijkstra's Shortest Path First (SPF) algorithm.

\textbf{Standard:} OSPFv2 is specified in RFC 2328 for IPv4. OSPFv3 supports IPv6 and is specified in RFC 5340.

\subsubsection{Why is OSPF needed?}

Imagine an organization has multiple routers connecting different departments and network segments.

Static routing requires administrators to manually configure routes. When the topology changes, static routes may need manual updates.

OSPF allows routers to dynamically:

\begin{itemize}
\item Discover neighboring routers.
\item Exchange network topology information.
\item Calculate optimal routes.
\item Adapt when links fail or recover.
\end{itemize}

\subsubsection{How OSPF works}

OSPF operation can be understood in six stages.

\begin{enumerate}
\item \textbf{Neighbor discovery:} Routers send Hello packets to discover neighboring OSPF routers.
\item \textbf{Adjacency formation:} Eligible neighbors establish an adjacency and synchronize their link-state information.
\item \textbf{LSA exchange:} Routers exchange Link-State Advertisements describing links and network topology.
\item \textbf{LSDB construction:} Each router builds a Link-State Database representing the topology of its OSPF area.
\item \textbf{SPF calculation:} Dijkstra's algorithm calculates the shortest paths.
\item \textbf{Routing table installation:} The best routes are installed in the routing table, subject to route-selection rules.
\end{enumerate}

\subsubsection{OSPF Packet Types}
\begin{table}[htbp]
\centering
\caption{OSPF Packet Types and Protocol Functions}
\label{tab:ospf_packet_types}
\begin{tabularx}{\textwidth}{p{4.2cm} p{2.2cm} X}
\toprule
\textbf{Packet Type Name} & \textbf{Type ID} & \textbf{Primary Operational Purpose} \\
\midrule
\textbf{Hello} & Type 1 & Discovers neighbors, establishes adjacencies, and acts as keepalive \\
\textbf{Database Description (DBD)} & Type 2 & Summarizes Link-State Database contents for database synchronization \\
\textbf{Link-State Request (LSR)} & Type 3 & Requests detailed LSA information missing from the local database \\
\textbf{Link-State Update (LSU)} & Type 4 & Carries one or more Link-State Advertisements (LSAs) across routers \\
\textbf{Link-State Ack (LSAck)} & Type 5 & Provides explicit Layer 3 acknowledgment for received LSUs \\
\bottomrule
\end{tabularx}
\end{table}

OSPF uses five principal packet types:


\subsubsection{OSPF areas}

OSPF divides a large routing domain into areas to improve scalability and reduce unnecessary topology information.

\textbf{Area 0 --- Backbone Area}

\begin{itemize}
\item The central area of a multi-area OSPF network.
\item Other OSPF areas connect to the backbone, either directly or through a virtual link.
\item Inter-area routing traffic passes through the backbone architecture.
\end{itemize}

Other areas can be used to organize networks by geography, department, or topology.

\subsubsection{OSPF Hierarchical Router Classifications}
\begin{table}[htbp]
\centering
\caption{OSPF Hierarchical Router Classifications}
\label{tab:ospf_router_types}
\begin{tabularx}{\textwidth}{p{4.2cm} X}
\toprule
\textbf{Router Classification} & \textbf{Topological Role and Architectural Placement} \\
\midrule
\textbf{Internal Router (IR)} & All active interfaces reside completely within a single non-backbone OSPF area \\
\textbf{Backbone Router (BR)} & Has at least one active interface assigned to the backbone area (Area 0) \\
\textbf{Area Border Router (ABR)} & Interfaces attach to two or more areas; connects standard areas to Area 0 \\
\textbf{AS Boundary Router (ASBR)} & Redistributes external routes from foreign routing protocols (BGP, RIP, static) into OSPF \\
\bottomrule
\end{tabularx}
\end{table}


A router can belong to more than one of these categories.

\subsubsection{OSPF cost}

OSPF uses a cost metric to determine the preferred path.

A commonly used Cisco default reference-bandwidth formula is:



\[
\text{Cost}=\frac{\text{Reference Bandwidth}}{\text{Interface Bandwidth}}
\]



For example, with a reference bandwidth of 100 Mbps:


Modern networks should configure an appropriate reference bandwidth so that higher-speed interfaces are distinguishable.

\textbf{Important:} OSPF implementations use integer interface costs. The formula and default reference bandwidth can vary by vendor and configuration.

\subsubsection{OSPF numerical example}

Consider the following paths between Router A and Router D:

\begin{itemize}
\item A $\to$ B $\to$ D: costs 10 + 20.
\item A $\to$ C $\to$ D: costs 5 + 10.
\item A $\to$ E $\to$ D: costs 15 + 15.
\end{itemize}

Calculate the total cost:



\[
C_{ABD}=10+20=30
\]





\[
C_{ACD}=5+10=15
\]





\[
C_{AED}=15+15=30
\]



Therefore, OSPF selects:



\[
\boxed{A\rightarrow C\rightarrow D}
\]



because its total cost is 15, the lowest among the available paths.

\subsubsection{OSPF characteristics}

\begin{itemize}
\item Dynamic routing protocol.
\item Link-state routing algorithm.
\item Uses Dijkstra's SPF algorithm.
\item Supports classless routing and CIDR.
\item Supports VLSM.
\item Supports authentication mechanisms.
\item Supports hierarchical routing through areas.
\item Uses IP protocol number 89, not TCP or UDP.
\item Supports equal-cost multipath routing (ECMP).
\end{itemize}

\subsubsection{OSPF administrative distance}

A commonly used Cisco administrative distance is:



\[
\boxed{110}
\]



Administrative distance is a Cisco route-source preference value, not an OSPF path metric.

\subsubsection{OSPF commands (Cisco IOS)}

\begin{lstlisting}
show ip route
show ip ospf
show ip ospf neighbor
show ip ospf database
show ip protocols
\end{lstlisting}

Example configuration:

\begin{lstlisting}
router ospf 1
 router-id 1.1.1.1
 network 192.168.1.0 0.0.0.255 area 0
 network 10.0.0.0 0.0.0.3 area 0
\end{lstlisting}

Here:

\begin{itemize}
\item \texttt{1} is the locally significant OSPF process ID.
\item \texttt{1.1.1.1} is the router ID.
\item \texttt{0.0.0.255} is the wildcard mask corresponding to a \texttt{/24} network mask.
\item \texttt{area 0} places the matching interfaces in the backbone area.
\end{itemize}

\subsection{BGP --- Border Gateway Protocol}

\subsubsection{What is BGP?}

Border Gateway Protocol (BGP) is a path-vector routing protocol used to exchange routing information between Autonomous Systems.

It is the principal inter-domain routing protocol of the Internet.

For example:

\begin{itemize}
\item An ISP operates one autonomous system.
\item A university network may operate another autonomous system.
\item A cloud provider may operate another autonomous system.
\end{itemize}

BGP allows these independently administered networks to exchange reachability information.

\textbf{Standard:} BGP-4, RFC 4271.

\subsubsection{Types of BGP: eBGP versus iBGP}
\begin{table}[htbp]
\centering
\caption{Comparison: External BGP (eBGP) versus Internal BGP (iBGP)}
\label{tab:ebgp_vs_ibgp}
\begin{tabularx}{\textwidth}{p{3.8cm} X X}
\toprule
\textbf{Parameter} & \textbf{External BGP (eBGP)} & \textbf{Internal BGP (iBGP)} \\
\midrule
\textbf{Autonomous System} & Peers reside in \textbf{different} AS numbers & Peers reside in the \textbf{same} AS number \\
\textbf{Primary Role} & Inter-domain routing across Internet edges & Distributing exterior routing information within an AS \\
\textbf{IP TTL Default} & TTL = 1 (peers directly connected by default) & TTL = 255 (multi-hop peering standard) \\
\textbf{Administrative Distance} & 20 (Cisco IOS default; highly trusted) & 200 (Cisco IOS default; lower precedence than IGPs) \\
\textbf{Loop Prevention} & Rejects routes containing local AS in \texttt{AS\_PATH} & Does not advertise routes learned from iBGP to other iBGP peers \\
\bottomrule
\end{tabularx}
\end{table}

There are two principal forms:

\textbf{1. eBGP --- External BGP}

\begin{itemize}
\item Exchanges routes between different autonomous systems.
\item Commonly used between ISPs and between organizations and ISPs.
\item Neighboring routers belong to different AS numbers.
\end{itemize}

\textbf{2. iBGP --- Internal BGP}

\begin{itemize}
\item Exchanges BGP routes between routers within the same autonomous system.
\item Helps distribute externally learned routes across an organization's internal BGP network.
\end{itemize}


The administrative distances shown are common Cisco defaults.

\subsubsection{How BGP works}

\begin{enumerate}
\item Two routers establish a TCP connection.
\item They establish a BGP session and exchange OPEN messages.
\item They exchange routing information through UPDATE messages.
\item Each router evaluates available paths according to BGP path attributes and routing policy.
\item The selected routes are advertised to eligible neighbors.
\end{enumerate}

\subsubsection{BGP Message Types}
\begin{table}[htbp]
\centering
\caption{Border Gateway Protocol (BGP-4) Message Types}
\label{tab:bgp_message_types}
\begin{tabularx}{\textwidth}{p{3.5cm} X}
\toprule
\textbf{Message Type} & \textbf{Primary Operational Purpose} \\
\midrule
\textbf{OPEN} & Initiates BGP peering session, negotiates BGP version, AS number, Hold Time, and BGP Router ID \\
\textbf{UPDATE} & Advertises new feasible network prefixes and path attributes, or withdraws invalid routes \\
\textbf{KEEPALIVE} & Periodic heartbeat verifying link liveness; sent at 1/3 of the negotiated Hold Time interval \\
\textbf{NOTIFICATION} & Signals error conditions or connection closure; immediately terminates the TCP peering session \\
\textbf{ROUTE-REFRESH} & Dynamically requests full re-advertisement of routes from a peer without session tear-down \\
\bottomrule
\end{tabularx}
\end{table}


\subsubsection{BGP Path Attributes and Selection Rules}
\begin{table}[htbp]
\centering
\caption{Essential BGP Path Attributes and Selection Precedence}
\label{tab:bgp_path_attributes}
\begin{tabularx}{\textwidth}{p{3.2cm} p{2.8cm} X}
\toprule
\textbf{Attribute} & \textbf{Category} & \textbf{Selection Rule and Interpretation} \\
\midrule
\textbf{WEIGHT} & Cisco proprietary & Locally significant; highest value preferred (default = 0, local = 32768) \\
\textbf{LOCAL\_PREF} & Well-known discretionary & AS-wide outbound exit preference; highest value preferred (default = 100) \\
\textbf{AS\_PATH} & Well-known mandatory & Sequence of traversed AS numbers; shortest AS path length preferred \\
\textbf{ORIGIN} & Well-known mandatory & Route origin type: IGP ($I$) $<$ EGP ($e$) $<$ Incomplete ($?$) \\
\textbf{MED} & Optional non-transitive & Inbound entry suggestion advertised to neighboring AS; lowest value preferred \\
\textbf{NEXT\_HOP} & Well-known mandatory & IP address of the next router used to reach the destination prefix \\
\textbf{COMMUNITY} & Optional transitive & Tags prefixes for group policy enforcement (e.g. \texttt{no-export}, \texttt{no-advertise}) \\
\bottomrule
\end{tabularx}
\end{table}

BGP does not simply choose the route with the lowest number of hops or the shortest geographical distance. It uses path attributes and routing policies.


A simplified BGP selection example:

Suppose a router receives two routes to the same destination:


If the routes are otherwise comparable, Route A may be preferred because its LOCAL\_PREF is higher. This illustrates that a shorter AS\_PATH does not automatically win over a higher LOCAL\_PREF.

Actual BGP best-path selection depends on vendor implementation, configuration, and the complete comparison sequence.

\subsubsection{AS\_PATH and loop prevention}

The AS\_PATH attribute records the autonomous systems a route has traversed.

When a BGP router receives a route containing its own AS number in the AS\_PATH, it normally rejects that route to prevent an inter-AS routing loop.

\subsubsection{BGP characteristics}

\begin{itemize}
\item Path-vector routing protocol.
\item Used for inter-domain routing.
\item Supports policy-based routing decisions.
\item Uses TCP port 179.
\item Supports CIDR and route aggregation.
\item Handles large-scale Internet routing.
\item Uses AS\_PATH for inter-AS loop prevention.
\end{itemize}

\subsubsection{BGP administrative distance}

Common Cisco defaults:

\begin{itemize}
\item eBGP: 20
\item iBGP: 200
\end{itemize}

\subsubsection{BGP commands (Cisco IOS)}

\begin{lstlisting}
show ip bgp
show ip bgp summary
show ip bgp neighbors
show ip route bgp
show running-config
\end{lstlisting}

Example configuration:

\begin{lstlisting}
router bgp 65001
 neighbor 192.0.2.2 remote-as 65002
 network 203.0.113.0 mask 255.255.255.0
\end{lstlisting}

This configures an eBGP neighbor in AS 65002 from a router in AS 65001.

The \texttt{network} statement advertises a route only when the corresponding prefix is present in the local routing table.

\subsubsection{STP vs OSPF vs BGP --- Important comparison}


\textbf{Exam memory points:}

\begin{itemize}
\item STP prevents switching loops.
\item OSPF uses Dijkstra's algorithm.
\item BGP is a path-vector protocol.
\item OSPF uses IP protocol number 89.
\item BGP uses TCP port 179.
\item OSPF is an IGP; BGP is an EGP.
\item STP elects a Root Bridge; OSPF elects a DR and BDR on eligible multi-access networks.
\end{itemize}


\newpage
\begin{center}
\vspace{0.5em}
{\color{primary}\rule{0.45\textwidth}{1.2pt}\\[4pt]}
{\Large\bfseries\color{primary} Section XV: Domain Name System (DNS) Architecture and Troubleshooting\\[2pt]}
{\color{primary}\rule{0.45\textwidth}{1.2pt}}
\vspace{0.5em}
\end{center}

\section{Domain Name System (DNS) Architecture and Troubleshooting}

\subsection{What is DNS?}
DNS stands for Domain Name System.
DNS is a hierarchical and distributed naming system that translates human-readable domain names into information needed by network applications, most commonly IP addresses.
For example:
www.example.com
may resolve to an IPv4 address such as:
93.184.216.34
Humans find domain names easier to remember than numerical IP addresses.
DNS allows users to access network resources using names instead of remembering their IP addresses.
Why is DNS necessary?
Without DNS, users would need to remember the IP address of every website or network service they want to access.
DNS provides:
Name-to-IP address resolution.
IP-to-name resolution through reverse DNS.
Service discovery through service records.
Distributed and scalable domain management.
Load distribution through multiple DNS records.
Flexible changes to server IP addresses without requiring users to remember new addresses.
Important: DNS is not limited to translating domain names into IP addresses. It also stores other information, including mail server records, name server records and service information.
\subsection{DNS hierarchy}
DNS follows a hierarchical structure.
The main components are:
Root DNS servers.
Top-Level Domain (TLD) DNS servers.
Authoritative DNS servers.
A DNS resolver performs the queries required to obtain the requested information.
A. Root DNS servers
Root DNS servers occupy the top level of the DNS hierarchy.
They direct resolvers toward the appropriate TLD name servers.
There are 13 named logical root server identities, represented by letters A through M. Each identity is implemented through many distributed server instances using anycast.
Root servers do not normally provide the final IP address of a website.
B. Top-Level Domain (TLD) servers
TLD servers manage information about domain extensions.
Examples:
.com
.org
.net
.edu
.gov
.bd
For example, a .com TLD server can direct a resolver toward the authoritative name servers responsible for example.com.
C. Authoritative DNS servers
An authoritative DNS server stores or serves the official DNS records for a domain.
For example, it may provide the A record for:
www.example.com
The authoritative server supplies the answer from the domain's DNS data.
D. Recursive DNS resolver
A recursive resolver receives DNS queries from client devices and obtains the answer on their behalf.
It may:
Return a cached answer.
Query root servers.
Query TLD servers.
Query authoritative servers.
Cache the result according to its TTL.
Common public recursive DNS resolver services include:
Google Public DNS: 8.8.8.8 and 8.8.4.4.
Cloudflare DNS: 1.1.1.1 and 1.0.0.1.
Quad9: 9.9.9.9.
These are examples of recursive resolver addresses, not authoritative servers for every domain.
\subsection{DNS resolution process}
Consider a user accessing:
www.example.com
The following sequence describes a simplified DNS lookup when the answer is not cached.
DNS resolution sequence
Client device
The browser or operating system requests the IP address of www.example.com.
Recursive DNS resolver
The resolver receives the query and checks its cache.
Root DNS server
If necessary, the resolver asks the root server where to find the .com TLD servers.
TLD DNS server
The resolver asks the .com server for the authoritative name servers of example.com.
Authoritative DNS server
The resolver queries the authoritative server for the required record.
Recursive resolver
The resolver receives the answer and caches it according to its TTL.
Client device
The resolver returns the IP address to the client, which can then establish a connection to the destination.
Important: The client usually sends its query to a recursive resolver, not directly to root, TLD and authoritative servers. The recursive resolver performs the iterative lookup process.
Recursive vs iterative DNS queries
Recursive query:
The client asks the resolver to obtain a complete answer or report that it cannot resolve the name.
The resolver takes responsibility for finding the answer.
Iterative query:
The queried DNS server returns the best information it has, which may be:
The requested answer.
A referral to another DNS server.
An error.
The resolver then continues the lookup using the referral.
\subsection{Recursive vs. Iterative DNS Queries}
\begin{table}[htbp]
\centering
\caption{Comparison: Recursive versus Iterative DNS Queries}
\label{tab:dns_recursive_vs_iterative}
\begin{tabularx}{\textwidth}{p{3.8cm} X X}
\toprule
\textbf{Evaluation Dimension} & \textbf{Recursive Query} & \textbf{Iterative Query} \\
\midrule
\textbf{Resolver Obligation} & Queried server MUST find answer or return failure & Queried server returns best referral it knows \\
\textbf{Operating Stage} & Client host $\to$ Local Recursive Resolver & Recursive Resolver $\to$ Root / TLD / Authoritative \\
\textbf{Resolution Workload} & Heavy workload on the queried recursive server & Workload distributed; asking resolver walks tree \\
\textbf{Returned Content} & Complete final IP address or NXDOMAIN & Delegation referral to next-level name server \\
\bottomrule
\end{tabularx}
\end{table}
\subsection{Types of DNS records}
DNS records store different types of information about domain names.
These are particularly important for examinations and practical DNS troubleshooting.
\subsection{DNS Resource Record Summary}
\begin{table}[htbp]
\centering
\caption{Primary Domain Name System (DNS) Resource Record Types}
\label{tab:dns_resource_records}
\begin{tabularx}{\textwidth}{p{2.2cm} p{4.2cm} X}
\toprule
\textbf{Record} & \textbf{Primary Purpose} & \textbf{Canonical Example / Syntax} \\
\midrule
\textbf{A} & Hostname to 32-bit IPv4 address & \texttt{example.com. IN A 93.184.216.34} \\
\textbf{AAAA} & Hostname to 128-bit IPv6 address & \texttt{example.com. IN AAAA 2606:2800:220:1:248::} \\
\textbf{CNAME} & Canonical alias to another domain name & \texttt{www.example.com. IN CNAME example.com.} \\
\textbf{MX} & Mail exchange server with priority & \texttt{example.com. IN MX 10 mail.example.com.} \\
\textbf{NS} & Authoritative name server for zone & \texttt{example.com. IN NS ns1.example.net.} \\
\textbf{PTR} & Pointer for reverse DNS resolution & \texttt{34.216.184.93.in-addr.arpa. IN PTR host.net.} \\
\textbf{SOA} & Zone authority, serial number, timers & Primary server, admin email, refresh, expire, negative TTL \\
\textbf{TXT} & Arbitrary text data (SPF, DKIM, DMARC) & \texttt{"v=spf1 include:\_spf.google.com \textasciitilde{}all"} \\
\textbf{SRV} & Service location, protocol, priority, port & \texttt{\_sip.\_tcp.example.com. IN SRV 10 60 5060 sip.com.} \\
\textbf{CAA} & Certification Authority Authorization & Specifies permitted SSL/TLS certificate issuers \\
\bottomrule
\end{tabularx}
\end{table}
A record
An A record maps a hostname to an IPv4 address.
Example:
www.example.com. 300 IN A 93.184.216.34
AAAA record
An AAAA record maps a hostname to an IPv6 address.
CNAME record
A CNAME record makes one domain name an alias of another canonical domain name.
Example:
www.example.com. IN CNAME example.com.
A CNAME record points to another name, not directly to an IP address.
A CNAME generally cannot coexist with other ordinary data records at the same owner name. This is why a zone apex commonly uses A/AAAA records rather than a conventional CNAME.
MX record
An MX record identifies the mail servers responsible for receiving email for a domain.
Example:
example.com. IN MX 10 mail.example.com.
The number 10 represents the preference value.
A lower preference number indicates a higher priority.
If multiple MX records exist, mail servers normally try the lowest-preference-value server first.
NS record
An NS record identifies the name servers responsible for a DNS zone.
Example:
example.com. IN NS ns1.example.net.
PTR record
A PTR record is used for reverse DNS lookup.
It maps an IP address back to a domain name.
For IPv4, reverse DNS uses the special domain:
in-addr.arpa
For IPv6, reverse DNS uses:
ip6.arpa
SOA record
SOA stands for Start of Authority.
It contains important information about a DNS zone.
Common fields include:
Primary name server.
Responsible party's email address, represented in DNS format.
Serial number.
Refresh interval.
Retry interval.
Expire interval.
Negative caching TTL information.
The SOA serial number helps secondary DNS servers determine whether the zone data has changed.
TXT record
TXT records store text information associated with a domain.
Common uses:
SPF email authentication information.
Domain ownership verification.
DKIM public key information.
DMARC policy information.
SRV record
SRV records specify the hostname and port for a particular service.
They commonly contain:
Priority.
Weight.
Port.
Target hostname.
\subsection{DNS port numbers and transport protocols}
DNS commonly uses:
Protocol
Port
Purpose
UDP
53
Most conventional DNS queries and responses
TCP
53
Zone transfers and DNS responses requiring TCP
Why does DNS use UDP?
Traditional DNS queries are usually small and benefit from UDP's low overhead.
UDP does not establish a connection before sending data.
However, DNS also uses TCP.
When does DNS use TCP?
Important situations include:
Zone transfers between DNS servers.
DNS responses that cannot be handled appropriately over the original UDP exchange.
Certain DNS configurations and resolver behaviours.
Exam point: DNS uses both UDP and TCP on port 53. The answer "DNS always uses UDP" is incorrect.
\subsection{DNS caching and TTL}
TTL stands for Time To Live.
A DNS record's TTL specifies how long a resolver or cache may retain that record before it should be considered expired and refreshed.
Example:
example.com. 3600 IN A 192.0.2.10
TTL:

\[
3600\text{ seconds}=1\text{ hour}
\]

If a resolver caches this record, it can normally use the cached answer for up to one hour, subject to DNS caching rules.
Why is DNS caching important?
Reduces repeated DNS queries.
Improves response time.
Reduces load on authoritative DNS servers.
Reduces network traffic.
What happens when a DNS record changes?
Suppose a domain changes its IP address.
If a resolver still has the old record cached, users may continue receiving the old IP address until the cached record expires or is otherwise refreshed.
This is commonly called DNS propagation delay, although the underlying issue is often distributed caching and different refresh times.
Important: Lowering a record's TTL shortly before changing it does not immediately clear copies already cached with the previous TTL.
4. DNS Troubleshooting
DNS problems can make a network connection appear broken even when the actual network connectivity is working.
For example:
A user cannot open www.example.com.
The same user can successfully reach the server using its IP address.
This suggests a possible DNS resolution problem.
However, DNS is not the only possible cause; browser, routing, firewall, TLS and application problems must also be considered.
\subsection{Common DNS Diagnostic Problems and Errors}
\begin{table}[htbp]
\centering
\caption{Common DNS Errors, Root Causes, and Diagnostics}
\label{tab:dns_errors_diagnostics}
\begin{tabularx}{\textwidth}{p{2.8cm} p{4.5cm} X}
\toprule
\textbf{Error Response} & \textbf{Diagnostic Meaning} & \textbf{Common Underlying Causes} \\
\midrule
\textbf{NXDOMAIN} & Non-Existent Domain & Typo in domain name, expired domain, missing zone record \\
\textbf{SERVFAIL} & Server Failure & DNSSEC validation failure, unreachable authoritative nameservers \\
\textbf{REFUSED} & Query Refused & Resolver access control list (ACL) block, open-resolver protection \\
\textbf{Timeout} & No response received & UDP port 53 blocked by firewall, network link failure, down resolver \\
\textbf{Wrong IP} & Incorrect address returned & Stale local resolver cache, poisoned cache, split-horizon DNS mismatch \\
\bottomrule
\end{tabularx}
\end{table}
\subsection{Step-by-step DNS troubleshooting}
A systematic approach is important when troubleshooting DNS issues.
Step 1: Check basic connectivity
First, determine whether the computer has network connectivity.
Linux:
ip addr ip route ping -c 4 1.1.1.1
Windows:
ipconfig ping 1.1.1.1
If communication with a known reachable IP address works but domain names fail, DNS becomes a stronger suspect.
However, ICMP may be blocked, so a failed ping alone does not prove that the network is down.
Step 2: Check DNS resolution
Use:
nslookup example.com
Or:
dig example.com
These commands query DNS and display the returned information.
Example:
\begin{lstlisting}
\$ nslookup example.com
Server: 192.168.1.1
Address: 192.168.1.1\#53

Non-authoritative answer:
Name: example.com
Address: 93.184.216.34
\end{lstlisting}
This indicates that the DNS resolver returned an IPv4 address.
Step 3: Query a different DNS resolver
Try querying a public DNS resolver directly:
nslookup example.com 8.8.8.8
Or:
dig @1.1.1.1 example.com
If the public resolver succeeds while the configured resolver fails, the problem may be related to the configured DNS server or its upstream resolution.
Step 4: Check DNS records
Use:
dig example.com A
For IPv6:
dig example.com AAAA
For mail servers:
dig example.com MX
For authoritative name servers:
dig example.com NS
For the SOA record:
dig example.com SOA
Step 5: Trace the DNS resolution process
Use:
dig +trace example.com
This shows the iterative lookup path through the DNS hierarchy.
It can help identify where the resolution process is failing.
Step 6: Check the configured DNS server
Windows:
ipconfig /all
Linux systems using systemd-resolved:
resolvectl status
Traditional Linux configuration may also be inspected with:
cat /etc/resolv.conf
The DNS server may be configured through:
DHCP.
Static network settings.
NetworkManager.
systemd-resolved.
VPN configuration.
Step 7: Clear DNS cache if stale information is suspected
Windows:
ipconfig /flushdns
Linux with systemd-resolved:
sudo resolvectl flush-caches
On other Linux distributions, cache-clearing commands depend on the DNS caching service being used.
Examples include nscd or dnsmasq.
Step 8: Check the hosts file
The hosts file can override normal DNS resolution for particular names.
Windows:
\texttt{C:\textbackslash{}Windows\textbackslash{}System32\textbackslash{}drivers\textbackslash{}etc\textbackslash{}hosts}
Linux:
/etc/hosts
A manually configured entry may cause a domain name to resolve to an unexpected IP address.
Step 9: Check firewall and DNS ports
DNS traffic commonly uses:
UDP port 53.
TCP port 53.
A firewall blocking these connections can prevent DNS resolution.
Testing TCP connectivity to a DNS server on port 53:
nc -vz 8.8.8.8 53
Note that a successful TCP test does not establish that UDP DNS traffic is working.
Step 10: Check authoritative DNS and delegation
If only a particular domain fails, investigate:
Whether its NS records are correct.
Whether its authoritative servers are reachable.
Whether the requested record exists.
Whether the zone has valid SOA information.
Whether DNSSEC validation is failing.
Whether parent-zone delegation is correct.
\subsection{Essential DNS Troubleshooting Commands}
\begin{table}[htbp]
\centering
\caption{Essential Diagnostic CLI Commands for DNS Troubleshooting}
\label{tab:dns_cli_commands}
\begin{tabularx}{\textwidth}{p{4.5cm} X}
\toprule
\textbf{CLI Diagnostic Command} & \textbf{Operational Diagnostic Purpose} \\
\midrule
\texttt{dig example.com +trace} & Performs full iterative traversal from root servers down to authoritative \\
\texttt{dig @1.1.1.1 example.com A} & Bypasses local resolver to test resolution directly against Cloudflare DNS \\
\texttt{dig example.com MX} & Queries Mail Exchange records with preference levels \\
\texttt{nslookup example.com 8.8.8.8} & Standard cross-platform query against Google Public DNS \\
\texttt{ipconfig /flushdns} & Flushes local Windows DNS client resolver cache \\
\texttt{sudo resolvectl flush-caches} & Clears systemd-resolved local cache on modern Linux systems \\
\texttt{nc -vz 8.8.8.8 53} & Verifies TCP port 53 network transport reachability through firewalls \\
\bottomrule
\end{tabularx}
\end{table}
\subsection{Scenario-based DNS troubleshooting questions}
\subsubsection*{Scenario 1: A website opens using its IP address but not its domain name. What is the likely issue?}
Possible DNS resolution problem.
Check:
nslookup example.com
Then verify the configured DNS server.
\subsubsection*{Scenario 2: One DNS server resolves a domain, but another returns NXDOMAIN. What might be the reason?}
Possible reasons:
Different DNS records or zone data.
Inconsistent authoritative servers.
Stale or incorrectly configured DNS information.
Different split-horizon DNS views.
\subsubsection*{Scenario 3: A website's IP address was changed, but some users still access the old server. Why?}
Likely DNS caching.
The old record may still exist in resolver caches until its TTL expires.
\subsubsection*{Scenario 4: A DNS query returns SERVFAIL. What should be investigated?}
Check:
DNS server availability.
Authoritative server health.
DNSSEC validation.
Zone configuration.
Delegation and upstream resolution.
\subsubsection*{Scenario 5: DNS works when connected directly to the office network but fails through VPN. Why?}
Possible causes:
Incorrect DNS server pushed by VPN.
Split-DNS configuration errors.
DNS traffic blocked by VPN routing or firewall rules.
Internal DNS zones unavailable through the VPN.

\newpage

\newpage
\begin{center}
\vspace{0.5em}
{\color{primary}\rule{0.45\textwidth}{1.2pt}\\[4pt]}
{\Large\bfseries\color{primary} Section XVI: Dynamic Host Configuration Protocol (DHCP) and Broadcast Management\\[2pt]}
{\color{primary}\rule{0.45\textwidth}{1.2pt}}
\vspace{0.5em}
\end{center}

\section{Dynamic Host Configuration Protocol (DHCP) and Broadcast Management}

\subsection{Excessive Broadcast Traffic}

\subsection{What is broadcast traffic?}

Broadcast traffic refers to network packets or Ethernet frames sent to all devices within a particular broadcast domain.

In IPv4 Ethernet networks, the Layer 2 broadcast MAC address is:

\begin{lstlisting}
FF:FF:FF:FF:FF:FF
\end{lstlisting}

The IPv4 limited broadcast address is:

\begin{lstlisting}
255.255.255.255
\end{lstlisting}

A broadcast is intended for all devices in the relevant broadcast domain, rather than a single destination.

\subsubsection{Types of IPv4 Communication}
\begin{table}[htbp]
\centering
\caption{IPv4 Transmission Addressing Paradigms}
\label{tab:ipv4_transmission_paradigms}
\begin{tabularx}{\textwidth}{p{2.8cm} p{4.2cm} X}
\toprule
\textbf{Traffic Type} & \textbf{Delivery Scope} & \textbf{Canonical Protocol Example} \\
\midrule
\textbf{Unicast} & Point-to-point (one sender to one host) & HTTP/HTTPS web traffic, SSH, file transfer \\
\textbf{Broadcast} & One sender to all hosts in broadcast domain & ARP requests, DHCP Discover, NetBIOS \\
\textbf{Multicast} & One sender to subscribed host group & OSPF updates (224.0.0.5), IPTV streaming \\
\bottomrule
\end{tabularx}
\end{table}


IPv6 does not use broadcast. It uses multicast for comparable group communication functions.

\subsection{What is excessive broadcast traffic?}

Excessive broadcast traffic occurs when a network receives an unusually large volume of broadcast frames, consuming bandwidth and processing resources.

It may result in:

\begin{itemize}
\item Increased network congestion.
\item Higher CPU utilization on network devices.
\item Increased packet-processing overhead on hosts.
\item Slower application performance.
\item Packet loss and communication failures.
\end{itemize}

\subsection{Common causes}

\subsubsection{1. Large broadcast domains}

A large Layer 2 network contains many devices in one broadcast domain.

Every broadcast frame is forwarded to the relevant ports within that domain, subject to switching and VLAN rules.

\subsubsection{2. Switching loops}

A Layer 2 switching loop can cause broadcast frames to circulate repeatedly.

This can develop into a broadcast storm.

\subsubsection{3. ARP traffic}

Address Resolution Protocol (ARP) requests are broadcasts in IPv4 Ethernet networks.

When a device needs to discover the MAC address corresponding to a local IPv4 address, it sends an ARP request.

For example:

\begin{lstlisting}
Who has 192.168.1.20?
Tell 192.168.1.10
\end{lstlisting}

The device owning \texttt{192.168.1.20} replies with its MAC address.

Large numbers of devices or excessive ARP requests can increase broadcast traffic.

\subsubsection{4. Faulty network devices}

A malfunctioning device or misconfigured network interface may generate excessive broadcasts.

\subsubsection{5. Malware or misconfigured applications}

Some malicious or incorrectly configured software may generate unusually high volumes of network traffic.

\subsection{Broadcast storm}

A broadcast storm is a condition in which excessive broadcast frames overwhelm network resources.

A common cause is a Layer 2 switching loop.

For example:

\begin{enumerate}
\item A host sends a broadcast Ethernet frame.
\item A switch forwards it through multiple paths.
\item Another switch receives and forwards copies.
\item Frames continue circulating through the loop.
\item The network becomes congested.
\end{enumerate}

A broadcast storm can make the network unusable.

\subsection{How to identify excessive broadcast traffic}

Network administrators can use:

\begin{itemize}
\item Switch interface statistics.
\item Network monitoring systems.
\item Packet analyzers such as Wireshark.
\item SNMP monitoring.
\item Interface utilization reports.
\item CPU and memory monitoring.
\end{itemize}

Useful Cisco commands:

\begin{lstlisting}
show interfaces
show interfaces counters
show spanning-tree
show mac address-table
show interfaces status
\end{lstlisting}

The exact available commands depend on the switch model and operating system.

In Wireshark, useful display filters include:

\begin{lstlisting}
eth.dst == ff:ff:ff:ff:ff:ff
\end{lstlisting}

This displays Ethernet broadcast frames.

For ARP traffic:

\begin{lstlisting}
arp
\end{lstlisting}

\subsection{Methods to Control Excessive Broadcast Traffic}
\begin{table}[htbp]
\centering
\caption{Mitigation Strategies for Excessive Broadcast Traffic}
\label{tab:broadcast_mitigation_strategies}
\begin{tabularx}{\textwidth}{p{3.8cm} X}
\toprule
\textbf{Mitigation Strategy} & \textbf{Technical Implementation and Impact} \\
\midrule
\textbf{VLAN Segmentation} & Breaks flat Layer 2 networks into smaller logical broadcast domains \\
\textbf{STP / RSTP Enforcement} & Eliminates switching loops and prevents broadcast storm formation \\
\textbf{Storm Control} & Suppresses ingress broadcast/multicast traffic when defined bandwidth threshold is crossed \\
\textbf{DHCP Snooping} & Drops invalid DHCP messages and intercepts rogue DHCP server broadcasts \\
\textbf{Dynamic ARP Inspection} & Intercepts and validates ARP requests against DHCP binding database \\
\textbf{Layer 3 Boundary Routing} & Terminates broadcast propagation at router and multilayer switch interfaces \\
\bottomrule
\end{tabularx}
\end{table}


\subsubsection{VLAN example}

Suppose an organization has 120 computers.

Without VLAN segmentation, all 120 devices may belong to the same broadcast domain.

With VLAN segmentation:

\begin{itemize}
\item VLAN 10: 40 computers.
\item VLAN 20: 40 computers.
\item VLAN 30: 40 computers.
\end{itemize}

Now, an ordinary Layer 2 broadcast in VLAN 10 is not forwarded into VLANs 20 and 30.

Inter-VLAN communication requires Layer 3 routing.

\textbf{Important:} VLANs divide broadcast domains, but they do not automatically prevent all network congestion or eliminate broadcasts within each VLAN.

\subsection{Dynamic Host Configuration Protocol (DHCP) Process}

\subsection{What is DHCP?}

Dynamic Host Configuration Protocol (DHCP) is a network management protocol that automatically assigns IP configuration information to devices.

It eliminates the need to manually configure every device's network settings.

DHCP can provide:

\begin{itemize}
\item IPv4 address.
\item Subnet mask.
\item Default gateway.
\item DNS server addresses.
\item Lease duration.
\item Other network configuration parameters.
\end{itemize}

\textbf{Standard:} DHCPv4 is specified in RFC 2131.

DHCPv4 uses:

\begin{itemize}
\item UDP port 67: Server.
\item UDP port 68: Client.
\end{itemize}

\subsection{Why is DHCP important?}

Consider an organization with 500 computers.

Manually configuring the IP address, subnet mask, gateway, and DNS settings for every computer would be time-consuming and error-prone.

DHCP automates this process.

Benefits include:

\begin{enumerate}
\item Centralized IP address management.
\item Reduced configuration errors.
\item Efficient reuse of IP addresses.
\item Easier network administration.
\item Support for mobile devices and changing network environments.
\end{enumerate}

\subsection{DORA process}

DHCP commonly uses a four-step process known as DORA:

High-priority examination topic

\begin{itemize}
\item Discover
\item Offer
\item Request
\item Acknowledge
\end{itemize}

\subsubsection{Step 1: DHCP Discover}

When a client joins a network and needs an IP address, it broadcasts a DHCP Discover message.

The client may not yet have an IP address, so it commonly uses:

\begin{lstlisting}
Source IP:      0.0.0.0
Destination IP: 255.255.255.255
\end{lstlisting}

The message is sent using UDP source port 68 and destination port 67.

Purpose: Discover available DHCP servers.

\subsubsection{Step 2: DHCP Offer}

A DHCP server receives the Discover message and responds with a DHCP Offer.

The offer may contain:

\begin{itemize}
\item Proposed IP address.
\item Subnet mask.
\item Default gateway.
\item DNS server.
\item Lease duration.
\end{itemize}

Example:

\begin{lstlisting}
Offered IP:     192.168.1.100
Subnet Mask:    255.255.255.0
Default Gateway: 192.168.1.1
DNS Server:     8.8.8.8
Lease Time:     24 hours
\end{lstlisting}

\subsubsection{Step 3: DHCP Request}

The client broadcasts a DHCP Request indicating which offered address and server it has selected.

Broadcasting this message also allows other DHCP servers to know that their offers were not selected.

\subsubsection{Step 4: DHCP Acknowledgment}

The selected DHCP server sends a DHCP ACK message confirming the lease.

The client can then configure its network interface with the assigned parameters.

\subsubsection{DORA Process Summary}
\begin{table}[htbp]
\centering
\caption{The DHCPv4 DORA Exchange Sequence}
\label{tab:dhcp_dora_summary}
\begin{tabularx}{\textwidth}{p{1.2cm} p{2.4cm} p{1.8cm} p{1.8cm} X}
\toprule
\textbf{Step} & \textbf{Message} & \textbf{Source IP} & \textbf{Dest IP} & \textbf{Operational Function} \\
\midrule
\textbf{1} & \textbf{Discover} & \texttt{0.0.0.0} & \texttt{255.255.255.255} & Broadcast from client searching for available DHCP servers \\
\textbf{2} & \textbf{Offer} & Server IP & Broadcast / Unicast & Unicast/broadcast from server offering IP lease configuration \\
\textbf{3} & \textbf{Request} & \texttt{0.0.0.0} & \texttt{255.255.255.255} & Client broadcast selecting one server's offer and declining others \\
\textbf{4} & \textbf{ACK} & Server IP & Broadcast / Unicast & Server confirms lease binding and assigns subnet, gateway, and DNS \\
\bottomrule
\end{tabularx}
\end{table}


\subsection{DHCP lease}

A DHCP lease is the period during which a client is permitted to use an assigned IP address.

For example:

\begin{lstlisting}
IP Address: 192.168.1.100
Lease Duration: 24 hours
\end{lstlisting}

The client can use this address during the lease period, subject to the lease and renewal process.

\subsubsection{Lease renewal}

DHCP clients do not normally wait until the lease expires before attempting renewal.

For a typical DHCPv4 lease:

\begin{itemize}
\item At T1, commonly 50\% of the lease duration, the client attempts renewal with the original server.
\item At T2, commonly 87.5\% of the lease duration, the client attempts rebinding, potentially contacting any available DHCP server.
\item If the lease expires without successful renewal, the client must stop using the leased address.
\end{itemize}

These are default timing rules, and actual client behavior can vary.

\subsection{DHCPv4 Protocol Message Types}
\begin{table}[htbp]
\centering
\caption{Complete DHCPv4 Protocol Message Specifications}
\label{tab:dhcp_message_specifications}
\begin{tabularx}{\textwidth}{p{3.2cm} X}
\toprule
\textbf{Message Type} & \textbf{Protocol Function} \\
\midrule
\textbf{DHCPDISCOVER} & Client broadcast to locate active DHCP servers on the subnet \\
\textbf{DHCPOFFER} & Server response offering specific IP lease, mask, gateway, and DNS parameters \\
\textbf{DHCPREQUEST} & Client broadcast confirming chosen offer, requesting lease renewal, or verifying lease \\
\textbf{DHCPACK} & Server acknowledgment containing committed network configuration parameters \\
\textbf{DHCPNAK} & Server negative acknowledgment rejecting client request (e.g. invalid subnet or expired lease) \\
\textbf{DHCPDECLINE} & Client message indicating offered IP address is already in use (detected via gratuitous ARP) \\
\textbf{DHCPRELEASE} & Client message gracefully relinquishing its assigned IP lease before expiration \\
\textbf{DHCPINFORM} & Client query requesting local configuration parameters without requesting an IP address \\
\bottomrule
\end{tabularx}
\end{table}


\subsection{DHCP relay agent}

A DHCP Discover message is normally broadcast within its local broadcast domain.

Routers do not ordinarily forward broadcasts between networks.

So how can a DHCP server in a central data center provide addresses to clients on multiple VLANs?

The solution is a DHCP relay agent.

A relay agent receives a DHCP broadcast and forwards the request to a DHCP server on another network.

Cisco configuration example:

\begin{lstlisting}
interface vlan 10
 ip address 192.168.10.1 255.255.255.0
 ip helper-address 192.168.1.10
\end{lstlisting}

Here:

\begin{itemize}
\item \texttt{192.168.10.1} is the Layer 3 interface for VLAN 10.
\item \texttt{192.168.1.10} is the DHCP server.
\item \texttt{ip helper-address} forwards eligible DHCP requests to that server.
\end{itemize}

\subsubsection{DHCP server scope}

A DHCP scope defines the range of IP addresses and related configuration parameters available for a particular subnet.

Example:

\begin{lstlisting}
Network:        192.168.1.0/24
Address range:  192.168.1.100--192.168.1.200
Gateway:        192.168.1.1
DNS:            8.8.8.8
Lease:          24 hours
\end{lstlisting}

Addresses outside the dynamic pool can be reserved for servers, printers, and infrastructure devices, provided they are managed to avoid address conflicts.

\subsection{DHCP Reservation versus Static IP Assignment}
\begin{table}[htbp]
\centering
\caption{Comparison: DHCP Reservation versus Static IP Assignment}
\label{tab:dhcp_res_vs_static}
\begin{tabularx}{\textwidth}{p{3.5cm} X X}
\toprule
\textbf{Parameter} & \textbf{DHCP Reservation} & \textbf{Static IP Configuration} \\
\midrule
\textbf{Configuration Point} & Central DHCP server scope & Local host network interface settings \\
\textbf{Binding Mechanism} & Client MAC address mapped to dedicated IP & Manually assigned IP address, mask, and gateway \\
\textbf{Parameter Changes} & Centrally pushed on lease renewal & Requires manual update on every device \\
\textbf{Primary Use Case} & Printers, VoIP phones, networked cameras & Core servers, router interfaces, firewalls \\
\bottomrule
\end{tabularx}
\end{table}


A DHCP reservation is often preferable when an administrator wants a device to receive a consistent IP address while retaining centralized management.

\subsection{DHCP troubleshooting}

Common problems:

\begin{itemize}
\item DHCP server unavailable.
\item DHCP scope exhausted.
\item Incorrect VLAN configuration.
\item DHCP relay misconfigured.
\item Firewall blocking DHCP traffic.
\item Rogue DHCP server.
\item Duplicate IP addresses.
\end{itemize}

Useful commands:

\textbf{Windows:}

\begin{lstlisting}
ipconfig /all
ipconfig /release
ipconfig /renew
ipconfig /flushdns
\end{lstlisting}

\textbf{Linux:}

\begin{lstlisting}
ip addr
ip route
resolvectl status
\end{lstlisting}

The DHCP client renewal command depends on the Linux distribution and network manager.

\subsubsection{APIPA}

If a Windows client is configured for DHCP but cannot obtain an address, it may automatically assign itself an IPv4 link-local address.

This is known as Automatic Private IP Addressing (APIPA).

Range:



\[
\boxed{169.254.0.0/16}
\]



Example:

\begin{lstlisting}
169.254.20.15
\end{lstlisting}

An APIPA address usually indicates that the client could not obtain a DHCP lease. It does not provide ordinary communication with other subnets through a default gateway.


\newpage
\begin{center}
\vspace{0.5em}
{\color{primary}\rule{0.45\textwidth}{1.2pt}\\[4pt]}
{\Large\bfseries\color{primary} Section XVII: Network Address Translation: Static NAT, Dynamic NAT, and PAT\\[2pt]}
{\color{primary}\rule{0.45\textwidth}{1.2pt}}
\vspace{0.5em}
\end{center}

\section{Network Address Translation: Static NAT, Dynamic NAT, and PAT}

\subsection{What is NAT?}
NAT stands for Network Address Translation.
NAT is a technique used by routers and firewalls to translate IP addresses in packet headers as traffic passes between networks.
It is commonly used to allow devices with private IPv4 addresses to communicate with external networks using public IPv4 addresses.
For example, consider a home network:
Device
Private IP
Computer
192.168.1.10
Smartphone
192.168.1.11
Laptop
192.168.1.12
Router's LAN interface
192.168.1.1
The router's WAN interface may have a public IPv4 address:
203.0.113.10
When these devices access the Internet, the router translates their private source addresses into an externally usable address.
This allows multiple internal devices to share a public IPv4 address when PAT is used.
Why is NAT necessary?
IPv4 address conservation: Multiple private devices can share a smaller number of public IPv4 addresses.
Private network connectivity: Allows private IPv4 networks to communicate with public networks.
Address flexibility: Internal IP addressing can change without necessarily changing externally visible addresses.
Network organization: Internal addressing can be independent of the public addressing scheme.
Inbound connection mapping: Static NAT and port forwarding can make internal services reachable from external networks.
Important: NAT is not a replacement for a firewall. Although it can make direct inbound connections less straightforward, NAT alone does not provide comprehensive network security.
\subsection{Private IPv4 address ranges}
Private IPv4 addresses are reserved for internal networks and are not globally routable on the public Internet.
According to RFC 1918, there are three private IPv4 address blocks.
Classful category
Private IP range
CIDR
Total addresses
Class A range
10.0.0.0 -- 10.255.255.255
/8
16,777,216
Class B range
172.16.0.0 -- 172.31.255.255
/12
1,048,576
Class C range
192.168.0.0 -- 192.168.255.255
/16
65,536
These are historical class labels for the ranges; modern IP networking uses CIDR.
Important distinction
Not every address that is not globally routable is an RFC 1918 private address.
For example:
127.0.0.1 is a loopback address.
169.254.0.0/16 is IPv4 link-local addressing.
100.64.0.0/10 is shared address space commonly used by carrier-grade NAT providers.
6. Types of NAT
The three major NAT techniques commonly discussed in networking examinations are:
Static NAT.
Dynamic NAT.
Port Address Translation (PAT).
\subsection{Static NAT}
Static NAT establishes a permanent one-to-one mapping between a private IP address and a public IP address.
For example:
Private IP
Public IP
192.168.1.10
203.0.113.10
192.168.1.11
203.0.113.11
192.168.1.12
203.0.113.12
The mapping remains fixed.
How does Static NAT work?
Suppose an internal web server has:
192.168.1.10
Its public NAT address is:
203.0.113.10
When an external client accesses the public address, the NAT device translates the destination address to the internal server's private address, subject to its firewall and forwarding policies.
For outgoing traffic, the reverse translation occurs.
Advantages
Permanent public-to-private IP mapping.
Suitable for servers that need a consistent public address.
Easier to identify specific internal devices externally.
Supports inbound connections when appropriate firewall policies are configured.
Disadvantages
Requires one public IP address for each permanently mapped private address.
Can consume many public IPv4 addresses.
Requires manual configuration and maintenance.
Example
A company operates an internal web server:
Private IP: 192.168.1.100 Public IP: 203.0.113.100
Static NAT establishes:
192.168.1.100 <-> 203.0.113.100
The public address remains associated with that internal device.
\subsection{Dynamic NAT}
Dynamic NAT translates private IP addresses into public IP addresses using a pool of available public IP addresses.
Unlike Static NAT, mappings are assigned dynamically when devices initiate communication.
Example:
Public IP pool:
203.0.113.10
203.0.113.11
203.0.113.12
Private devices:
192.168.1.10
192.168.1.11
192.168.1.12
Possible temporary mappings:
Private IP
Assigned public IP
192.168.1.10
203.0.113.11
192.168.1.11
203.0.113.10
192.168.1.12
203.0.113.12
The assignments are dynamic and may change when translation entries expire or are released.
Characteristics
Uses a pool of public IP addresses.
Mapping is temporary.
Public IPs are assigned as needed.
A public address generally cannot be assigned to another active mapping simultaneously in the same translation context.
Advantages
Reduces the need for permanent one-to-one public IP assignments.
Public addresses can be reused when mappings expire.
Useful for networks with a pool of public addresses.
Disadvantages
Requires a sufficient public IP pool for concurrent mappings.
Does not inherently allow several internal devices to share one public address simultaneously.
Mapping and troubleshooting are more complex than simple static mappings.
Important: If all available public IP addresses in the pool are occupied, additional devices may be unable to establish new translations until an address becomes available.
\subsection{Port Address Translation (PAT)}
PAT stands for Port Address Translation.
It is also commonly called:
NAT Overload.
Network Address Port Translation (NAPT).
Many-to-one NAT.
PAT allows multiple private devices to share a single public IPv4 address by using different transport-layer port numbers.
It is the most common form of NAT used by home broadband routers.
How does PAT work?
Consider three internal devices:
Device
Private IP
PC
192.168.1.10
Laptop
192.168.1.11
Smartphone
192.168.1.12
All devices share the public IP:
203.0.113.10
The router creates different translation entries.
Internal source
Translated source
192.168.1.10:5000
203.0.113.10:40001
192.168.1.11:5000
203.0.113.10:40002
192.168.1.12:5000
203.0.113.10:40003
Although all three devices use the same public IP address, their translated source ports distinguish the communication sessions.
The router maintains a translation table to identify which internal device should receive returning traffic.
PAT translation table
Protocol
Internal address
Public translated address
TCP
192.168.1.10:5000
203.0.113.10:40001
TCP
192.168.1.11:5000
203.0.113.10:40002
TCP
192.168.1.12:5000
203.0.113.10:40003
The port numbers shown are illustrative. Actual port selection depends on the NAT implementation and available mappings.
Advantages
Conserves public IPv4 addresses.
Allows many internal devices to share one public IP.
Commonly supported by consumer and enterprise routers.
Reduces the need for large public IPv4 address pools.
Supports simultaneous communication sessions.
Disadvantages
Requires translation-table maintenance.
Can complicate inbound connections.
Can cause problems for certain protocols that embed addressing information in payloads.
May complicate peer-to-peer communication.
Translation state can be lost when the NAT device restarts or mappings expire.
\subsection{Master Comparison: Static NAT vs. Dynamic NAT vs. PAT}
\begin{table}[htbp]
\centering
\caption{Master Comparative Analysis: Static NAT vs. Dynamic NAT vs. PAT}
\label{tab:nat_master_comparison}
\begin{tabularx}{\textwidth}{p{3.2cm} X X X}
\toprule
\textbf{Feature} & \textbf{Static NAT} & \textbf{Dynamic NAT} & \textbf{PAT (NAT Overload)} \\
\midrule
\textbf{Mapping Relation} & Strict 1-to-1 & Temporary 1-to-1 & Many-to-1 (or many-to-few) \\
\textbf{Public IP Pool} & 1 dedicated IP per host & Pool of public IPs & Single shared public IP \\
\textbf{Port Translation} & No (IP header only) & No (IP header only) & Yes (translates source ports) \\
\textbf{Inbound Access} & Seamless (predictable IP) & Complex / temporary & Requires explicit Port Forwarding \\
\textbf{IPv4 Conservation} & None ($1:1$ ratio) & Low (bounded by pool) & Extremely high ($>60,000$ sessions) \\
\textbf{Primary Application} & Web/mail enterprise servers & Temporary pool access & Home/office broadband Internet \\
\bottomrule
\end{tabularx}
\end{table}
Remember the three NAT techniques
Static NAT
Private IP
Public IP
One private address permanently maps to one public address.
Dynamic NAT
Private IPs
Public IP pool
Private addresses receive temporary public IP mappings from a pool.
PAT
Many private IPs
One public IP
Port numbers distinguish the different communication sessions.
8. How NAT Works: Packet Translation
Understanding source and destination address translation is important for practical network troubleshooting.
Consider a computer:
192.168.1.10
It wants to access a web server:
198.51.100.20
The router has a public address:
203.0.113.10
Outgoing packet
Before NAT:
Field
Value
Source IP
192.168.1.10
Source port
50000
Destination IP
198.51.100.20
Destination port
443
After PAT:
Field
Value
Source IP
203.0.113.10
Source port
40001
Destination IP
198.51.100.20
Destination port
443
The NAT device changes the source IP address and, in this example, the source port.
It also updates relevant packet checksums.
Returning packet
The web server sends a response to:
203.0.113.10:40001
The NAT router checks its translation table and determines that this response belongs to:
192.168.1.10:50000
It translates the destination address and port back to the internal values and forwards the packet to the original computer.
Important: Basic NAT translates IP addresses. PAT additionally translates transport-layer port numbers, commonly TCP or UDP ports.
9. Source NAT vs Destination NAT
Another important distinction is based on which IP address is translated.
\subsection{Source NAT (SNAT)}
Source NAT changes the source IP address of a packet.
It is commonly used when internal devices access external networks.
Example:
Before: Source: 192.168.1.10 Destination: 198.51.100.20 After: Source: 203.0.113.10 Destination: 198.51.100.20
The source address changes from private to public.
\subsection{Destination NAT (DNAT)}
Destination NAT changes the destination IP address of a packet.
It is commonly used for port forwarding and publishing internal services.
Example:
Before: Source: 198.51.100.20 Destination: 203.0.113.10 After: Source: 198.51.100.20 Destination: 192.168.1.100
An external request to a public address is forwarded to an internal server.
\subsection{Comparison: Source NAT (SNAT) vs. Destination NAT (DNAT)}
\begin{table}[htbp]
\centering
\caption{Comparison: Source NAT (SNAT) versus Destination NAT (DNAT)}
\label{tab:snat_vs_dnat}
\begin{tabularx}{\textwidth}{p{3.8cm} X X}
\toprule
\textbf{Parameter} & \textbf{Source NAT (SNAT)} & \textbf{Destination NAT (DNAT)} \\
\midrule
\textbf{Translated Header} & Source IP address (and source port in PAT) & Destination IP address (and destination port) \\
\textbf{Traffic Direction} & Outbound traffic initiated by internal host & Inbound traffic initiated by external client \\
\textbf{Primary Objective} & Hiding private IP space; Internet access & Publishing internal servers; Port Forwarding \\
\textbf{Cisco Terminology} & \texttt{ip nat inside source list ...} & \texttt{ip nat inside source static <local> <global>} \\
\bottomrule
\end{tabularx}
\end{table}
10. NAT vs PAT vs Port Forwarding
These terms are related but should not be treated as interchangeable.
NAT is the general technique of translating IP addresses.
PAT additionally uses port numbers to distinguish multiple communication sessions.
Port forwarding is a configuration that directs incoming traffic on a particular public IP address and port to a specified internal IP address and port.
Example: Hosting an internal web server
Suppose:
Public IP: 203.0.113.10
Internal web server: 192.168.1.100
Public service port: 8080
Internal service port: 80
A port-forwarding rule can be configured:
203.0.113.10:8080 | v 192.168.1.100:80
An external client connecting to the public address on port 8080 is forwarded to the internal server on port 80.
This is a form of destination translation, commonly called DNAT.
A firewall policy must also permit the incoming connection.
11. Double NAT and Carrier-Grade NAT (CGNAT)
\subsection{Double NAT}
Double NAT occurs when traffic passes through two devices that both perform NAT.
Example:
Computer | v Home Router (NAT) | v ISP Router (NAT) | v Internet
Potential problems:
Complicated port forwarding.
Difficulty hosting services.
Problems with certain VPNs.
Complicated troubleshooting.
Some peer-to-peer connectivity issues.
Double NAT is not necessarily a problem for ordinary outbound web browsing, but it can complicate inbound connectivity.
\subsection{Carrier-Grade NAT (CGNAT)}
Carrier-Grade NAT is a technique used by Internet Service Providers to allow multiple subscribers to share public IPv4 addresses.
An ISP may assign a customer an address from:
100.64.0.0/10
This is the shared address range reserved for carrier-grade NAT environments.
The customer's router may perform NAT, and the ISP performs another translation before traffic reaches the public Internet.
This is often called large-scale NAT.
\subsection{Comparison: Conventional NAT vs. Carrier-Grade NAT (CGNAT)}
\begin{table}[htbp]
\centering
\caption{Comparison: Enterprise / SOHO NAT versus Carrier-Grade NAT (CGNAT)}
\label{tab:cgnat_comparison}
\begin{tabularx}{\textwidth}{p{3.8cm} X X}
\toprule
\textbf{Feature} & \textbf{Conventional Customer-Premises NAT} & \textbf{Carrier-Grade NAT (CGNAT / LSN)} \\
\midrule
\textbf{Operator} & Home user or enterprise IT department & Internet Service Provider (ISP) core network \\
\textbf{Address Allocation} & RFC 1918 (e.g. 192.168.1.0/24, 10.0.0.0/8) & Shared address block \textbf{100.64.0.0/10} (RFC 6598) \\
\textbf{Number of NAT Layers} & Single NAT translation stage & Two stages (Double NAT: Customer CPE + ISP CGNAT) \\
\textbf{Inbound Port Forwarding} & Fully controllable by end-user router & Blocked at ISP level; requires public IP addon \\
\bottomrule
\end{tabularx}
\end{table}
12. NAT Configuration and Troubleshooting Commands
For Cisco-style network equipment, these commands are useful for understanding and inspecting NAT configuration.
Display NAT translations:
show ip nat translations
Display NAT statistics:
show ip nat statistics
Display running configuration:
show running-config
Clear dynamic NAT translations:
clear ip nat translation *
The exact behaviour and syntax can depend on the router's operating system and configuration.
Basic Cisco IOS Static NAT example
configure terminal interface gigabitEthernet 0/0 ip nat inside exit interface gigabitEthernet 0/1 ip nat outside exit ip nat inside source static 192.168.1.100 203.0.113.100 end
Here:
192.168.1.100 is the internal private address.
203.0.113.100 is the public address.
ip nat inside identifies the internal-facing interface.
ip nat outside identifies the external-facing interface.
Basic Cisco IOS PAT example
A typical PAT configuration uses an access control list to identify the internal addresses and overloads the router's outside interface address.
configure terminal access-list 1 permit 192.168.1.0 0.0.0.255 interface gigabitEthernet 0/0 ip nat inside exit interface gigabitEthernet 0/1 ip nat outside exit ip nat inside source list 1 interface gigabitEthernet 0/1 overload end
The keyword overload enables PAT.
The router translates multiple internal devices to the outside interface's IP address using port mappings.
13. High-Value Numerical and Conceptual Questions
\subsubsection*{Problem 1: How many public IP addresses are needed for Static NAT for 50 internal devices?}
Static NAT uses one public IP for each permanently mapped internal IP.
Therefore:

\[
\boxed{50\text{ public IP addresses}}
\]

\subsubsection*{Problem 2: An organization has 100 private devices but only 5 public IP addresses. Can Dynamic NAT support all 100 devices simultaneously?}
Not with a basic one-to-one Dynamic NAT pool of only five addresses.
At most five simultaneous one-to-one mappings can exist using that pool.
PAT can allow many devices to share fewer public IP addresses by distinguishing sessions using port numbers.
\subsubsection*{Problem 3: How many internal devices can use one public IP through PAT?}
There is no simple fixed number based solely on the number of devices.
It depends on:
Available translated port numbers.
NAT implementation.
Protocol.
Concurrent sessions.
Resource limits.
Translation and timeout policies.
PAT supports many devices sharing a public IP address, subject to these limits.
\subsubsection*{Problem 4: Which NAT technique is most commonly used in home routers?}

\[
\boxed{\text{PAT (NAT Overload)}}
\]

It allows multiple private devices to share one public IPv4 address.
\subsubsection*{Problem 5: Which NAT technique provides a permanent one-to-one mapping?}

\[
\boxed{\text{Static NAT}}
\]

\subsubsection*{Problem 6: Which NAT technique uses a pool of public IP addresses?}

\[
\boxed{\text{Dynamic NAT}}
\]

\subsubsection*{Problem 7: What is the difference between NAT and PAT?}
NAT translates IP addresses.
PAT translates IP addresses and transport-layer port numbers, allowing multiple internal devices to share a public IP address.

\newpage

\newpage
\begin{center}
\vspace{0.5em}
{\color{primary}\rule{0.45\textwidth}{1.2pt}\\[4pt]}
{\Large\bfseries\color{primary} Section XVIII: Network Security Systems: Firewall, IDS, IPS, and DMZ Architecture\\[2pt]}
{\color{primary}\rule{0.45\textwidth}{1.2pt}}
\vspace{0.5em}
\end{center}

\section{Network Security Systems: Firewall, IDS, IPS, and DMZ Architecture}



These technologies help protect networks by controlling traffic, detecting suspicious activity, and preventing attacks.

\subsection{Firewall}

\subsubsection{What is a firewall?}

A firewall is a network security system that monitors and controls incoming and outgoing network traffic according to predefined security rules.

It acts as a security boundary between networks or between a device and a network.

For example, an organization's firewall may:

\begin{itemize}
\item Allow employees to access approved websites.
\item Allow HTTPS traffic to web servers.
\item Block unauthorized incoming connections.
\item Restrict access to internal databases.
\item Allow remote access only through approved VPN connections.
\end{itemize}

\subsubsection{How does a firewall work?}

A firewall examines network traffic and compares it against configured rules.

A simplified rule might look like this:


Actual firewall rules also depend on direction, interface, connection state, address translation, and rule order.

\subsubsection{Types of firewalls}

\textbf{1. Packet-filtering firewall}

\begin{itemize}
\item Examines packet header information.
\item Uses source IP, destination IP, protocol, and port numbers.
\item Does not normally inspect application payloads.
\item Generally operates at the Network and Transport layers.
\end{itemize}

Example:

Allow TCP traffic from \texttt{192.168.1.0/24} to destination port 443.

\textbf{2. Stateful inspection firewall}

\begin{itemize}
\item Tracks the state of active network connections.
\item Maintains a state table.
\item Evaluates packets in the context of existing connections.
\item Can recognize whether a packet belongs to an established session.
\end{itemize}

For example, when an internal computer initiates an HTTPS connection, the firewall can permit legitimate return traffic without requiring a separate broad inbound rule.

\textbf{3. Application-layer firewall}

\begin{itemize}
\item Examines application-layer traffic.
\item Can inspect protocols such as HTTP.
\item Can apply rules based on application-level information.
\item May block malicious requests that ordinary packet filtering cannot identify.
\end{itemize}

\textbf{4. Next-Generation Firewall (NGFW)}

A Next-Generation Firewall combines traditional firewall capabilities with additional security features, which may include:

\begin{itemize}
\item Deep Packet Inspection (DPI).
\item Application awareness and control.
\item Intrusion prevention.
\item User identity-based policies.
\item Threat intelligence integration.
\item URL filtering.
\item Malware detection.
\end{itemize}

The exact feature set varies by vendor and product.

\textbf{5. Web Application Firewall (WAF)}

A WAF protects web applications by inspecting HTTP and HTTPS requests.

It can help detect and block attacks such as:

\begin{itemize}
\item SQL injection.
\item Cross-site scripting (XSS).
\item Malicious HTTP requests.
\item Certain forms of application-layer abuse.
\end{itemize}

A WAF protects web applications; it does not replace a general network firewall.

\subsubsection{Firewall Deployment Architectures}
\begin{table}[htbp]
\centering
\caption{Network Firewall Architectural Deployment Models}
\label{tab:firewall_architectures}
\begin{tabularx}{\textwidth}{p{3.8cm} X}
\toprule
\textbf{Deployment Model} & \textbf{Architectural Scope and Protection Focus} \\
\midrule
\textbf{Perimeter Firewall} & Sits at network boundary separating internal trusted LAN from untrusted Internet \\
\textbf{Host-Based Firewall} & Software agent running directly on endpoints/servers (e.g. \texttt{iptables}, Windows Defender) \\
\textbf{DMZ Architecture} & Multi-homed firewall isolating publicly accessible services from the internal LAN \\
\textbf{Internal Segmentation} & Controls lateral east-west traffic between departmental subnets or data center tiers \\
\bottomrule
\end{tabularx}
\end{table}

Common firewall deployment models include:


\subsubsection{DMZ --- Demilitarized Zone}

A DMZ is a separate network segment used to host services that must be accessible from external networks.

Examples include:

\begin{itemize}
\item Public web servers.
\item Public DNS servers.
\item Mail gateways.
\item Reverse proxies.
\end{itemize}

A simplified logical architecture:

\begin{lstlisting}
Internet
   |
Firewall
   |
   +---- DMZ
   |      |
   |      +---- Public Web Server
   |      +---- Mail Gateway
   |
   +---- Internal Network
          |
          +---- Employee Computers
          +---- Database Servers
          +---- File Servers
\end{lstlisting}

The DMZ limits direct exposure of the internal network.

For example, external users may be allowed to access a public web server in the DMZ, while direct access from the Internet to the internal database is denied.

\subsubsection{Advantages of firewalls}

\begin{itemize}
\item Prevent unauthorized network access.
\item Enforce security policies.
\item Control inbound and outbound traffic.
\item Restrict access to sensitive services.
\item Support network segmentation.
\item Provide logging and monitoring information.
\end{itemize}

\subsubsection{Limitations of firewalls}

\begin{itemize}
\item Cannot guarantee protection against every attack.
\item Misconfigured rules may create security weaknesses.
\item Encrypted traffic can limit inspection unless suitable inspection capabilities are configured.
\item Cannot independently prevent every insider threat.
\item Cannot replace endpoint protection, secure application development, or intrusion detection.
\end{itemize}

\subsection{IDS --- Intrusion Detection System}

\subsubsection{What is IDS?}

An Intrusion Detection System (IDS) monitors network traffic or system activities to identify suspicious behavior, security violations, or potential attacks.

Its primary function is to \textbf{detect and alert}.

An IDS generally does not automatically block the suspicious traffic.

\subsubsection{How does IDS work?}

An IDS collects and analyzes information from network traffic or host activities.

When it identifies suspicious behavior, it generates an alert for the administrator or security monitoring system.

For example:

\begin{enumerate}
\item An IDS monitors traffic entering a server network.
\item It identifies repeated connection attempts matching a known attack pattern.
\item It generates an alert.
\item A security administrator investigates the event.
\end{enumerate}

\subsubsection{Types of IDS}

\textbf{1. NIDS --- Network-based Intrusion Detection System}

\begin{itemize}
\item Monitors network traffic.
\item Usually deployed at strategic network locations.
\item Can identify suspicious traffic patterns and known attack signatures.
\end{itemize}

Example: Monitoring traffic entering a data center.

\textbf{2. HIDS --- Host-based Intrusion Detection System}

\begin{itemize}
\item Runs on an individual host.
\item Monitors system activities, logs, files, and other host-level events.
\item Can detect unauthorized file changes and suspicious processes, depending on its capabilities.
\end{itemize}

Example: Monitoring a database server for unauthorized modifications to critical system files.

\subsubsection{IDS detection techniques}

\textbf{1. Signature-based detection}

Compares observed activity with known attack signatures.

Advantages:

\begin{itemize}
\item Effective for known threats with reliable signatures.
\item Can produce relatively precise alerts when signatures are well maintained.
\end{itemize}

Limitations:

\begin{itemize}
\item May fail to detect new or previously unknown attacks.
\item Requires updated signatures.
\end{itemize}

\textbf{2. Anomaly-based detection}

Establishes a baseline of normal behavior and identifies deviations.

Example:

A server normally receives 100 requests per minute but suddenly receives 50,000 requests per minute.

This unusual pattern may trigger an alert.

Advantages:

\- Can identify previously unseen or unusual behavior.

Limitations:

\begin{itemize}
\item May generate false positives.
\item Requires suitable baselines and tuning.
\end{itemize}

\textbf{3. Stateful protocol analysis}

Compares observed protocol behavior with expected protocol standards and normal protocol states.

It can detect unusual or invalid protocol behavior.

\subsubsection{IDS limitations}

\begin{itemize}
\item Generates alerts rather than necessarily blocking traffic.
\item May produce false positives.
\item Can miss sophisticated or unknown attacks.
\item Requires ongoing monitoring and tuning.
\item Encrypted traffic may limit network-level visibility.
\end{itemize}

\subsection{IPS --- Intrusion Prevention System}

\subsubsection{What is IPS?}

An Intrusion Prevention System (IPS) monitors network or system activities to detect malicious behavior and automatically take preventive action.

Its primary function is to \textbf{detect and prevent}.

Depending on its configuration, an IPS may:

\begin{itemize}
\item Drop malicious packets.
\item Block suspicious connections.
\item Reset network sessions.
\item Temporarily block an offending source.
\item Generate security alerts.
\end{itemize}

\subsubsection{Types of IPS}

\textbf{1. NIPS --- Network-based Intrusion Prevention System}

Monitors network traffic and can block malicious traffic before it reaches protected systems.

\textbf{2. HIPS --- Host-based Intrusion Prevention System}

Runs on an individual host and can prevent suspicious or unauthorized activities on that system.

\subsubsection{How an IPS works}

\begin{enumerate}
\item Traffic passes through or is monitored by the IPS.
\item The IPS inspects the traffic.
\item It identifies activity matching its detection policies.
\item It applies a configured preventive action.
\item It records the event for security monitoring.
\end{enumerate}

\subsubsection{Intrusion Detection (IDS) vs. Intrusion Prevention (IPS)}
\begin{table}[htbp]
\centering
\caption{Comparative Analysis: Intrusion Detection (IDS) vs. Intrusion Prevention (IPS)}
\label{tab:ids_vs_ips}
\begin{tabularx}{\textwidth}{p{3.2cm} X X}
\toprule
\textbf{Parameter} & \textbf{Intrusion Detection System (IDS)} & \textbf{Intrusion Prevention System (IPS)} \\
\midrule
\textbf{Core Function} & Passive detection and real-time alerting & Active detection and automated threat mitigation \\
\textbf{Inline Placement} & Out-of-band via SPAN / TAP mirroring & In-line directly in the packet forwarding path \\
\textbf{Response Action} & Generates alerts, logs events, notifies SIEM & Drops packets, resets TCP sessions (\texttt{TCP RST}), blocks IPs \\
\textbf{Traffic Latency Impact} & Zero impact on network throughput & Can introduce latency; potential single point of failure \\
\textbf{False Positive Risk} & Generates administrative alert fatigue & Can disrupt legitimate business traffic \\
\bottomrule
\end{tabularx}
\end{table}


\subsubsection{Important distinction}

An IDS can be deployed out-of-band, receiving a copy of network traffic for analysis.

An IPS is commonly deployed inline, meaning traffic passes through it and can be blocked.

However, modern products may combine IDS and IPS functions.

\subsubsection{Master Security Comparison: Firewall vs. IDS vs. IPS}
\begin{table}[htbp]
\centering
\caption{Master Security Tool Comparison: Firewall vs. IDS vs. IPS}
\label{tab:firewall_vs_ids_vs_ips}
\begin{tabularx}{\textwidth}{p{2.8cm} X X X}
\toprule
\textbf{Feature} & \textbf{Firewall} & \textbf{IDS} & \textbf{IPS} \\
\midrule
\textbf{Primary Function} & Access control \& session state & Threat detection \& monitoring & Threat detection \& blocking \\
\textbf{Decision Basis} & Port, IP, protocol, session state & Attack signatures \& anomalies & Attack signatures \& anomalies \\
\textbf{Placement} & In-line boundary gateway & Out-of-band (SPAN/TAP) & In-line network path \\
\textbf{Action on Match} & Allow or Deny (Drop/Reject) & Generate alert to admin/SIEM & Drop packet, reset session \\
\textbf{Inspection Depth} & Headers \& state (NGFW: DPI) & Full packet payload \& stream & Full packet payload \& stream \\
\bottomrule
\end{tabularx}
\end{table}


\textbf{Memory aid:}

\begin{itemize}
\item Firewall: Controls access.
\item IDS: Detects suspicious activity.
\item IPS: Detects and prevents suspicious activity.
\end{itemize}


\newpage
\begin{center}
\vspace{0.5em}
{\color{primary}\rule{0.45\textwidth}{1.2pt}\\[4pt]}
{\Large\bfseries\color{primary} Section XIX: Secure Network Architecture, Defense in Depth, and Layer 2 Security\\[2pt]}
{\color{primary}\rule{0.45\textwidth}{1.2pt}}
\vspace{0.5em}
\end{center}

\section{Secure Network Architecture, Defense in Depth, and Layer 2 Security}



\subsection{What is secure network architecture?}

Secure network architecture is the design and organization of network infrastructure, devices, services, and security controls to protect information systems from unauthorized access, attacks, data loss, and service disruption.

It involves integrating security into the network design rather than adding security controls only after deployment.

A secure architecture should provide:

\begin{itemize}
\item Confidentiality.
\item Integrity.
\item Availability.
\item Authentication.
\item Authorization.
\item Accountability.
\end{itemize}

\subsection{The CIA Triad Information Security Model}
\begin{table}[htbp]
\centering
\caption{The CIA Triad Information Security Model}
\label{tab:cia_triad}
\begin{tabularx}{\textwidth}{p{3.2cm} p{4.2cm} X}
\toprule
\textbf{Pillar} & \textbf{Primary Security Objective} & \textbf{Technical Implementation Mechanisms} \\
\midrule
\textbf{Confidentiality} & Preventing unauthorized information disclosure & AES-256 encryption, TLS 1.3, IPsec, RBAC, MFA \\
\textbf{Integrity} & Guaranteeing data accuracy and non-tampering & Cryptographic hashes (SHA-256), HMAC, digital signatures \\
\textbf{Availability} & Ensuring authorized users have reliable access & Redundant links, link aggregation, HA failover, BCDR \\
\bottomrule
\end{tabularx}
\end{table}

The CIA Triad is a fundamental information security model.

\textbf{1. Confidentiality}

Ensures that information is accessible only to authorized individuals or systems.

Examples:

\begin{itemize}
\item Encryption.
\item Access control.
\item Authentication.
\item Data classification.
\end{itemize}

\textbf{2. Integrity}

Ensures that information remains accurate and is not modified without authorization.

Examples:

\begin{itemize}
\item Cryptographic hashes.
\item Digital signatures.
\item File integrity monitoring.
\item Access controls.
\end{itemize}

\textbf{3. Availability}

Ensures that systems and information remain accessible when required.

Examples:

\begin{itemize}
\item Redundant network links.
\item Backup systems.
\item Load balancing.
\item Disaster recovery.
\item DDoS mitigation.
\end{itemize}


\subsection{Core principles of secure network architecture}

\subsubsection{1. Defense in Depth}

Defense in Depth means implementing multiple independent layers of security.

If one security mechanism fails, other mechanisms can still provide protection.

Example:

An organization may use:

\begin{enumerate}
\item Perimeter firewall.
\item Network segmentation.
\item IDS/IPS.
\item Endpoint protection.
\item Multi-factor authentication.
\item Encryption.
\item Centralized security monitoring.
\end{enumerate}

No single security control should be assumed to provide complete protection.

\subsubsection{2. Least Privilege}

The principle of least privilege means giving users, applications, and systems only the permissions necessary to perform their assigned tasks.

Example:

A regular employee should not have administrator access to the organization's database servers.

Benefits:

\begin{itemize}
\item Reduces unauthorized access.
\item Limits damage from compromised accounts.
\item Reduces accidental modification.
\item Limits the impact of malware.
\end{itemize}

\subsubsection{3. Network Segmentation}

Network segmentation divides a network into smaller logical or physical sections.

It can be implemented using:

\begin{itemize}
\item VLANs.
\item Subnets.
\item Routers.
\item Firewalls.
\item Access Control Lists (ACLs).
\item Software-defined networking policies.
\end{itemize}

Example:

An organization may separate:


Segmentation reduces unnecessary communication between devices and can restrict lateral movement after a system is compromised.

\subsubsection{4. Zero Trust Architecture}

Zero Trust is a security approach based on the principle:

\textbf{Never implicitly trust; continuously verify.}

Traditional perimeter-based security often assumes that devices inside the corporate network are trustworthy.

Zero Trust does not automatically trust a user or device simply because it is connected to an internal network.

Its major principles include:

\begin{itemize}
\item Verify identity and device security.
\item Apply least-privilege access.
\item Authenticate and authorize access to individual resources.
\item Monitor access continuously.
\item Assume that a network may already be compromised.
\end{itemize}

Example:

An employee accessing an internal application may need to provide:

\begin{itemize}
\item Valid credentials.
\item Multi-factor authentication.
\item Device compliance information.
\item Authorization for that specific application.
\end{itemize}

Being connected to the office Wi-Fi alone does not grant unrestricted access.

\subsubsection{5. Secure Remote Access}

Remote access should be protected using suitable authentication and encryption mechanisms.

Common technologies:

\begin{itemize}
\item IPsec VPN.
\item SSL/TLS-based VPN.
\item Multi-factor authentication.
\item Secure administrative access through SSH.
\item Role-based access control.
\end{itemize}

Avoid exposing administrative interfaces directly to untrusted networks.

\subsubsection{6. Encryption}

Encryption protects data from unauthorized disclosure.

Two important categories are:

\textbf{Encryption in transit}

Protects information while it travels across a network.

Examples:

\begin{itemize}
\item HTTPS using TLS.
\item SSH.
\item IPsec VPN.
\end{itemize}

\textbf{Encryption at rest}

Protects stored information.

Examples:

\begin{itemize}
\item Full-disk encryption.
\item Database encryption.
\item Encrypted backups.
\end{itemize}

Encryption does not automatically guarantee data integrity, authentication, or secure key management. Appropriate cryptographic modes and protocols are required.

\subsubsection{7. Secure DNS and DHCP}

DNS and DHCP are important network services and should be protected.

Security measures include:

\begin{itemize}
\item Restricting access to DNS administration.
\item Monitoring suspicious DNS queries.
\item Using DNSSEC where appropriate.
\item Preventing unauthorized DHCP servers.
\item Enabling DHCP snooping on supported switches.
\item Using appropriate network segmentation.
\end{itemize}

\subsubsection{8. Access Control Lists (ACLs)}

An ACL is a set of rules used to permit or deny network traffic.

Example:

\begin{lstlisting}
access-list 10 permit 192.168.1.0 0.0.0.255
access-list 10 deny any
\end{lstlisting}

This standard Cisco IPv4 ACL permits the specified source network and denies other sources when applied to an interface or feature.

ACLs can control access based on parameters such as:

\begin{itemize}
\item Source IP address.
\item Destination IP address.
\item Protocol.
\item TCP/UDP port.
\end{itemize}

\textbf{Standard ACLs} generally filter by source IPv4 address.

\textbf{Extended ACLs} can filter by source, destination, protocol, and port information.

\subsubsection{9. Security Monitoring and Logging}

A secure architecture needs continuous visibility into network activities.

Useful tools and systems include:

\begin{itemize}
\item SIEM (Security Information and Event Management).
\item IDS/IPS.
\item Syslog servers.
\item Network monitoring systems.
\item NetFlow or IPFIX.
\item Endpoint Detection and Response (EDR).
\end{itemize}

SIEM platforms collect and correlate security events from multiple sources to support detection, investigation, and response.

\subsubsection{10. High Availability and Redundancy}

Security architecture must also protect availability.

Common approaches:

\begin{itemize}
\item Redundant routers.
\item Redundant switches.
\item Firewall high-availability pairs.
\item Multiple ISP connections.
\item Load balancers.
\item Backup power.
\item Regular backups and disaster recovery testing.
\end{itemize}

Redundancy reduces single points of failure, although it must be designed carefully to avoid introducing switching loops or configuration inconsistencies.

\subsection{Common Network Attacks and Architectural Defense Controls}
\begin{table}[htbp]
\centering
\caption{Major Network Attacks and Architectural Defense Controls}
\label{tab:network_attacks_defenses}
\begin{tabularx}{\textwidth}{p{3.2cm} X X}
\toprule
\textbf{Attack Paradigm} & \textbf{Threat Mechanism} & \textbf{Defensive Network Countermeasures} \\
\midrule
\textbf{ARP Spoofing} & Forges ARP replies to hijack gateway IP & Dynamic ARP Inspection (DAI), DHCP Snooping \\
\textbf{MAC Flooding} & Floods switch with forged MACs to force hub behavior & Port Security (limit MACs per port, shutdown violation) \\
\textbf{VLAN Hopping} & Exploits DTP auto-trunking or double tagging & Disable DTP, change native VLAN, disable unused ports \\
\textbf{Rogue DHCP Server} & Distributes malicious gateway/DNS parameters & DHCP Snooping (designate trusted vs. untrusted ports) \\
\textbf{IP Spoofing} & Forges source IP headers to bypass filters & Unicast Reverse Path Forwarding (uRPF), egress ACLs \\
\textbf{DDoS Volumetric} & Overwhelms bandwidth with botnet traffic & Anycast BGP, Cloud scrubbing, upstream rate limiting \\
\textbf{Man-in-the-Middle} & Intercepts and alters unencrypted traffic & End-to-end TLS 1.3, Mutual TLS (mTLS), IPsec ESP \\
\bottomrule
\end{tabularx}
\end{table}


\subsubsection{DHCP snooping and Dynamic ARP Inspection}

These are particularly important Layer 2 security mechanisms.

\textbf{DHCP Snooping}

A Layer 2 security feature that distinguishes trusted and untrusted switch ports for DHCP traffic.

It helps prevent unauthorized DHCP servers from distributing incorrect network configuration.

It can also build a DHCP snooping binding database containing information such as:

\begin{itemize}
\item Client MAC address.
\item Assigned IP address.
\item VLAN.
\item Interface.
\item Lease information.
\end{itemize}

\textbf{Dynamic ARP Inspection (DAI)}

DAI validates ARP packets against trusted IP-to-MAC bindings, commonly obtained from the DHCP snooping database.

It helps prevent ARP spoofing and ARP poisoning.

\subsubsection{Comparison: DHCP Snooping vs. Dynamic ARP Inspection}
\begin{table}[htbp]
\centering
\caption{Comparison: DHCP Snooping versus Dynamic ARP Inspection (DAI)}
\label{tab:dhcp_snoop_vs_dai}
\begin{tabularx}{\textwidth}{p{3.8cm} X X}
\toprule
\textbf{Feature} & \textbf{DHCP Snooping} & \textbf{Dynamic ARP Inspection (DAI)} \\
\midrule
\textbf{Target Threat} & Rogue DHCP servers, IP starvation & ARP poisoning, ARP spoofing, MitM attacks \\
\textbf{Inspected Traffic} & DHCP protocol messages (UDP 67/68) & ARP request and reply frames (EtherType 0x0806) \\
\textbf{Port Trust Model} & Trusted (uplink to DHCP) / Untrusted (hosts) & Trusted (switch uplinks) / Untrusted (access ports) \\
\textbf{Binding Database} & \textbf{Builds} the DHCP Snooping Binding Table & \textbf{Leverages} the DHCP Snooping Binding Table \\
\bottomrule
\end{tabularx}
\end{table}


\subsection{Example: Secure enterprise network design}

Consider a medium-sized organization with employees, servers, guests, and public web services.

A reasonable architecture would include:

\begin{itemize}
\item A perimeter firewall between the Internet and the organization.
\item A DMZ for public-facing services.
\item VLAN segmentation for departments.
\item Internal firewalls or ACLs to control inter-VLAN traffic.
\item IDS/IPS monitoring at suitable network boundaries.
\item MFA for remote access.
\item Centralized logging and security monitoring.
\item Encrypted communication for sensitive services.
\item Redundant network infrastructure.
\end{itemize}

The objective is to reduce the attack surface, limit unauthorized access, and contain the impact of a security incident.


\newpage
\begin{center}
\vspace{0.5em}
{\color{primary}\rule{0.45\textwidth}{1.2pt}\\[4pt]}
{\Large\bfseries\color{primary} Section XX: Flow Control, Data Link Reliability, and Stop-and-Wait ARQ\\[2pt]}
{\color{primary}\rule{0.45\textwidth}{1.2pt}}
\vspace{0.5em}
\end{center}

\section{Flow Control, Data Link Reliability, and Stop-and-Wait ARQ}



\subsection{What is Stop-and-Wait?}

Stop-and-Wait is a simple flow-control and reliable data transmission protocol in which the sender transmits one frame and waits for an acknowledgment (ACK) before transmitting the next frame.

It ensures that the sender does not overwhelm the receiver with frames.

The sender proceeds only after receiving an acknowledgment for the transmitted frame or after handling a timeout.

It is a fundamental concept in data communication and is commonly studied alongside Automatic Repeat reQuest (ARQ) protocols.

\subsection{Why is Stop-and-Wait needed?}

Consider a sender transmitting data to a receiver over a network.

If the sender continuously transmits frames without waiting, the receiver may not be able to process them quickly enough.

Frames may be lost, and the receiver may not have sufficient buffer space.

Stop-and-Wait addresses this by allowing only one outstanding frame at a time.

Its basic operation is:

\begin{enumerate}
\item Sender transmits a frame.
\item Sender starts a timer.
\item Receiver receives the frame.
\item Receiver checks whether the frame is valid.
\item Receiver sends an ACK.
\item Sender receives the ACK and transmits the next frame.
\end{enumerate}

\subsection{Stop-and-Wait operation}

Suppose a sender needs to transmit four frames:

\begin{lstlisting}
Sender                         Receiver
   |                               |
   |---------- Frame 1 ----------->|
   |                               |
   |<------------- ACK ------------|
   |                               |
   |---------- Frame 2 ----------->|
   |                               |
   |<------------- ACK ------------|
   |                               |
   |---------- Frame 3 ----------->|
   |                               |
   |<------------- ACK ------------|
   |                               |
   |---------- Frame 4 ----------->|
   |                               |
   |<------------- ACK ------------|
   |                               |
\end{lstlisting}

The sender waits for an acknowledgment after each frame.

\subsection{Stop-and-Wait ARQ}

Basic Stop-and-Wait provides flow control, but reliable transmission also requires handling lost or corrupted frames.

Stop-and-Wait ARQ adds:

\begin{itemize}
\item Acknowledgments.
\item Timers.
\item Retransmissions.
\item Sequence numbers.
\item Error detection.
\end{itemize}

ARQ stands for Automatic Repeat reQuest.

\subsubsection{How Stop-and-Wait ARQ works}

\textbf{Case 1: Successful transmission}

\begin{enumerate}
\item Sender transmits Frame 0.
\item Receiver receives Frame 0 correctly.
\item Receiver sends ACK.
\item Sender receives ACK and sends Frame 1.
\end{enumerate}

\textbf{Case 2: Data frame is lost}

\begin{enumerate}
\item Sender transmits a frame.
\item The frame is lost before reaching the receiver.
\item Receiver cannot acknowledge it.
\item Sender's timer expires.
\item Sender retransmits the frame.
\end{enumerate}

\textbf{Case 3: ACK is lost}

\begin{enumerate}
\item Sender transmits Frame 0.
\item Receiver receives Frame 0 and sends ACK.
\item The ACK is lost.
\item Sender does not receive the ACK before timeout.
\item Sender retransmits Frame 0.
\item Receiver recognizes the retransmission as a duplicate and does not deliver the data a second time.
\item Receiver sends the acknowledgment again.
\end{enumerate}

This prevents duplicate delivery to the upper layer.

\subsection{Sequence numbers}

Stop-and-Wait ARQ commonly uses a one-bit sequence number:



\[
\boxed{0,\ 1,\ 0,\ 1,\ldots}
\]



Why only one bit?

Because the sender can have only one outstanding frame.

The receiver needs to distinguish:

\begin{itemize}
\item A new frame.
\item A retransmission of the previous frame.
\end{itemize}

For example:


Sequence numbers are essential because an ACK may be lost, causing the sender to retransmit a frame already received.

\subsection{Stop-and-Wait performance}

Stop-and-Wait is simple but can be inefficient on high-bandwidth or long-distance links because the sender remains idle while waiting for acknowledgments.

\subsubsection{Important formulas}

Let:

\begin{itemize}
\item $L$ = Frame size in bits.
\item $R$ = Link transmission rate in bits per second.
\item $T_t$ = Frame transmission time.
\item $T_p$ = One-way propagation delay.
\item $T_a$ = ACK transmission time.
\item $T_{proc}$ = Relevant processing time.
\end{itemize}

Frame transmission time:



\[
\boxed{T_t=\frac{L}{R}}
\]



Propagation time:



\[
\boxed{T_p=\frac{d}{v}}
\]



where $d$ is distance and $v$ is propagation speed.

Ignoring ACK transmission and processing time, the cycle time is:



\[
T_{\text{cycle}}=T_t+2T_p
\]



The sender transmits one frame during this cycle.

Therefore, the approximate utilization is:



\[
\boxed{U=\frac{T_t}{T_t+2T_p}}
\]



Substituting $T_t=L/R$:



\[
\boxed{U=\frac{1}{1+2a}}
\]



where:



\[
a=\frac{T_p}{T_t}
\]



This is the standard idealized Stop-and-Wait utilization formula.

If ACK transmission time is significant, a more complete cycle-time model is:



\[
T_{\text{cycle}}=T_t+2T_p+T_a+T_{proc}
\]



The exact processing terms depend on the system being modeled.

\subsubsection{Numerical problem 1: Calculate utilization}

A sender transmits frames over a link with:

\begin{itemize}
\item Frame size = 1000 bits.
\item Bandwidth = 1 Mbps.
\item One-way propagation delay = 10 ms.
\end{itemize}

Calculate Stop-and-Wait utilization.

\textbf{Step 1: Calculate transmission time}



\[
T_t=\frac{L}{R}
\]





\[
T_t=\frac{1000}{1,000,000}
\]





\[
T_t=0.001s=1ms
\]



\textbf{Step 2: Calculate the cycle time}



\[
T_{\text{cycle}}=T_t+2T_p
\]





\[
=1+2(10)
\]





\[
=21ms
\]



\textbf{Step 3: Calculate utilization}



\[
U=\frac{1}{21}
\]





\[
U=0.04762
\]





\[
\boxed{U=4.76\%}
\]



The sender uses only approximately 4.76\% of the link's transmission capacity for frame transmission under this idealized model.

\subsubsection{Numerical problem 2: Calculate throughput}

Given:

\begin{itemize}
\item Frame size = 2000 bits.
\item Bandwidth = 2 Mbps.
\item One-way propagation delay = 5 ms.
\item Ignore ACK transmission and processing time.
\end{itemize}

First, transmission time:



\[
T_t=\frac{2000}{2,000,000}
\]





\[
T_t=0.001s=1ms
\]



Cycle time:



\[
T_{\text{cycle}}=1+2(5)=11ms
\]



Throughput:



\[
\text{Throughput}=\frac{\text{Frame size}}{\text{Cycle time}}
\]





\[
=\frac{2000}{0.011}
\]





\[
\boxed{181,818.18\text{ bits/s}}
\]



This is approximately 181.82 kbps.

\subsubsection{\textbf{Numerical problem 3: Calculate utilization using }$a$}

Suppose:



\[
a=4
\]



Then:



\[
U=\frac{1}{1+2a}
\]





\[
U=\frac{1}{1+8}
\]





\[
\boxed{U=11.11\%}
\]



\subsection{Advantages and disadvantages}


\subsection{Stop-and-Wait versus Sliding Window Protocols}
\begin{table}[htbp]
\centering
\caption{Comparison: Stop-and-Wait versus Sliding Window Protocols}
\label{tab:stop_and_wait_vs_sliding_window}
\begin{tabularx}{\textwidth}{p{3.8cm} X X}
\toprule
\textbf{Parameter} & \textbf{Stop-and-Wait ARQ} & \textbf{Sliding Window (Go-Back-N / Selective Repeat)} \\
\midrule
\textbf{Sender Window ($W_s$)} & Exactly $1$ frame outstanding & $W_s > 1$ (pipelined multiple frames in flight) \\
\textbf{Sequence Numbers} & 1 bit ($0$ and $1$) & $k$ bits (e.g. $0$ to $2^k-1$) \\
\textbf{Link Utilization} & Very low on high-BDP links ($U = \frac{1}{1+2a}$) & High link utilization ($U \approx 1$ when $W_s \ge 1+2a$) \\
\textbf{Buffer Overhead} & Minimal (single frame buffer) & Larger (buffers for outstanding / out-of-order frames) \\
\textbf{Protocol Complexity} & Very simple logic and timing & More complex acknowledgment and retransmission logic \\
\bottomrule
\end{tabularx}
\end{table}

Sliding Window Protocols improve performance by allowing multiple frames to be transmitted before receiving acknowledgments for all of them.


\subsubsection{Go-Back-N ARQ}

\begin{itemize}
\item Sender can transmit multiple frames within a window.
\item Uses cumulative acknowledgments.
\item If a frame is lost or corrupted, the sender retransmits that frame and subsequent outstanding frames.
\end{itemize}

\subsubsection{Selective Repeat ARQ}

\begin{itemize}
\item Sender can transmit multiple frames within a window.
\item Receiver can buffer correctly received out-of-order frames.
\item Only lost or corrupted frames need retransmission.
\end{itemize}

\textbf{Important examination point:} Stop-and-Wait ARQ is a special case of the sliding-window approach with a sender window size of one.

\subsection{Exam-Oriented Flow Control Revision and Formulas}

\subsection{Important facts to memorize}


\subsection{Important numerical formulas}



\newpage
\begin{center}
\vspace{0.5em}
{\color{primary}\rule{0.45\textwidth}{1.2pt}\\[4pt]}
{\Large\bfseries\color{primary} Section XXI: High-Value Examination Questions, Scenarios, and Practice Problems\\[2pt]}
{\color{primary}\rule{0.45\textwidth}{1.2pt}}
\vspace{0.5em}
\end{center}

\section{High-Value Examination Questions, Scenarios, and Practice Problems}

\subsection{Subnetting, Topology, and Latency Numerical Problems}

\begin{quote}
\textbf{Numerical Problem 1: A /26 IPv4 subnet has how many conventional usable host addresses?}\\[4pt]
\quad A. 30\\[2pt]
\quad B. 62\\[2pt]
\quad C. 64\\[2pt]
\quad D. 126\\[2pt]
\vspace{4pt}
\textbf{Answer: B (62).} Host bits $h = 32 - 26 = 6$. Total addresses $= 2^6 = 64$. Conventional usable hosts $= 2^6 - 2 = 62$.
\end{quote}

\begin{quote}
\textbf{Numerical Problem 2: Borrowing 3 host bits creates how many subnets?}\\[4pt]
\quad A. 6\\[2pt]
\quad B. 8\\[2pt]
\quad C. 14\\[2pt]
\quad D. 16\\[2pt]
\vspace{4pt}
\textbf{Answer: B (8).} Borrowed bits $s = 3$. Number of subnets $= 2^s = 2^3 = 8$.
\end{quote}

\begin{quote}
\textbf{Numerical Problem 3: What is the subnet mask for a /28 prefix?}\\[4pt]
\quad A. 255.255.255.192\\[2pt]
\quad B. 255.255.255.224\\[2pt]
\quad C. 255.255.255.240\\[2pt]
\quad D. 255.255.255.248\\[2pt]
\vspace{4pt}
\textbf{Answer: C (255.255.255.240).} Prefix /28 has 4 host bits in the fourth octet. Block size $= 2^4 = 16$. Mask $= 256 - 16 = 240$.
\end{quote}

\begin{quote}
\textbf{Numerical Problem 4: Which subnet contains the host IP 192.168.1.150/25?}\\[4pt]
\quad A. 192.168.1.0/25\\[2pt]
\quad B. 192.168.1.64/25\\[2pt]
\quad C. 192.168.1.128/25\\[2pt]
\quad D. 192.168.1.192/25\\[2pt]
\vspace{4pt}
\textbf{Answer: C (192.168.1.128/25).} Prefix /25 block size $= 128$. Subnets: \texttt{192.168.1.0/25} (0--127) and \texttt{192.168.1.128/25} (128--255). 150 falls within 128--255.
\end{quote}

\begin{quote}
\textbf{Numerical Problem 5: What minimum prefix length supports 200 usable hosts?}\\[4pt]
\quad A. /24\\[2pt]
\quad B. /25\\[2pt]
\quad C. /26\\[2pt]
\quad D. /27\\[2pt]
\vspace{4pt}
\textbf{Answer: A (/24).} Usable host requirement $2^h - 2 \ge 200 \implies h = 8$ ($2^8 - 2 = 254$). Prefix $= 32 - 8 = /24$. (/25 supports only 126 hosts).
\end{quote}

\begin{quote}
\textbf{Numerical Problem 6: Combining four contiguous /24 networks results in what prefix length?}\\[4pt]
\quad A. /21\\[2pt]
\quad B. /22\\[2pt]
\quad C. /23\\[2pt]
\quad D. /26\\[2pt]
\vspace{4pt}
\textbf{Answer: B (/22).} Number of combined networks $k = 4$. New prefix $= 24 - \log_2(4) = 24 - 2 = /22$.
\end{quote}

\begin{quote}
\textbf{Numerical Problem 7: In VLSM design, which requirement must be assigned first?}\\[4pt]
\quad A. Smallest host requirement\\[2pt]
\quad B. Largest host requirement\\[2pt]
\quad C. Random order\\[2pt]
\quad D. Last alphabetically\\[2pt]
\vspace{4pt}
\textbf{Answer: B (Largest host requirement).} Allocating the largest subnet first maintains proper binary power-of-two boundary alignment and prevents subnet fragmentation.
\end{quote}

\begin{quote}
\textbf{Numerical Problem 8: A full mesh topology with 9 devices requires how many physical duplex links?}\\[4pt]
\quad A. 18\\[2pt]
\quad B. 27\\[2pt]
\quad C. 36\\[2pt]
\quad D. 72\\[2pt]
\vspace{4pt}
\textbf{Answer: C (36).} Formula $L = n(n-1)/2 = 9(8)/2 = 36$ links.
\end{quote}

\begin{quote}
\textbf{Numerical Problem 9: A 20,000-bit packet is transmitted over a 10 Mbps link. What is the transmission delay?}\\[4pt]
\quad A. 0.2 ms\\[2pt]
\quad B. 2 ms\\[2pt]
\quad C. 20 ms\\[2pt]
\quad D. 200 ms\\[2pt]
\vspace{4pt}
\textbf{Answer: B (2 ms).} Transmission delay $D_t = L / R = 20,000 / (10 \times 10^6) = 0.002\text{ s} = 2\text{ ms}$.
\end{quote}

\begin{quote}
\textbf{Numerical Problem 10: A signal travels 600 km at $2 \times 10^8$ m/s. What is the propagation delay?}\\[4pt]
\quad A. 0.3 ms\\[2pt]
\quad B. 3 ms\\[2pt]
\quad C. 30 ms\\[2pt]
\quad D. 300 ms\\[2pt]
\vspace{4pt}
\textbf{Answer: B (3 ms).} Propagation delay $D_p = d / v = 600,000\text{ m} / (2 \times 10^8\text{ m/s}) = 0.003\text{ s} = 3\text{ ms}$.
\end{quote}

\begin{quote}
\textbf{Numerical Problem 11: A network link has 200 Mbps bandwidth and 40 ms RTT. What is the Bandwidth-Delay Product (BDP)?}\\[4pt]
\quad A. 800 kb\\[2pt]
\quad B. 1 Mb\\[2pt]
\quad C. 5 Mb\\[2pt]
\quad D. 8 Mb (1 MB)\\[2pt]
\vspace{4pt}
\textbf{Answer: D (8 Mb or 1 MB).} $\text{BDP} = R \times \text{RTT} = (200 \times 10^6\text{ bps}) \times 0.04\text{ s} = 8,000,000\text{ bits} = 1,000,000\text{ bytes} = 1\text{ MB}$.
\end{quote}

\subsection{Modulation, LPM, DNS, and NAT Practice Questions}
\begin{quote}
\textbf{Practice Question 1: Which modulation technique changes the amplitude of a carrier signal?}\\[4pt]
\quad A. FM\\[2pt]
\quad B. PM\\[2pt]
\quad C. AM\\[2pt]
\quad D. FSK\\[2pt]
\vspace{4pt}
\textbf{Answer: C (AM).} Amplitude Modulation modifies the carrier amplitude $A_c$ proportionally to the instantaneous amplitude of the information signal while maintaining frequency and phase constant.
\end{quote}
\begin{quote}
\textbf{Practice Question 2: Which digital modulation technique changes both amplitude and phase?}\\[4pt]
\quad A. ASK\\[2pt]
\quad B. FSK\\[2pt]
\quad C. BPSK\\[2pt]
\quad D. QAM\\[2pt]
\vspace{4pt}
\textbf{Answer: D (QAM).} Quadrature Amplitude Modulation alters both carrier amplitude and phase simultaneously, achieving high spectral efficiency.
\end{quote}
\begin{quote}
\textbf{Practice Question 3: How many bits per symbol are transmitted using 256-QAM?}\\[4pt]
\quad A. 4\\[2pt]
\quad B. 6\\[2pt]
\quad C. 8\\[2pt]
\quad D. 16\\[2pt]
\vspace{4pt}
\textbf{Answer: C (8 bits/symbol).} Number of bits per symbol $= \log_2 M = \log_2 256 = 8\text{ bits/symbol}$.
\end{quote}
\begin{quote}
\textbf{Practice Question 4: A 16-QAM system operates at 2,000 baud. What is its ideal bit rate?}\\[4pt]
\quad A. 2 kbps\\[2pt]
\quad B. 4 kbps\\[2pt]
\quad C. 8 kbps\\[2pt]
\quad D. 16 kbps\\[2pt]
\vspace{4pt}
\textbf{Answer: C (8 kbps).} In 16-QAM, bits/symbol $= \log_2 16 = 4$. Ideal bit rate $= \text{Baud rate} \times 4 = 2000 \times 4 = 8000\text{ bps} = 8\text{ kbps}$.
\end{quote}
\begin{quote}
\textbf{Practice Question 5: What is the minimum Nyquist sampling rate for a 5 kHz band-limited signal?}\\[4pt]
\quad A. 2.5 kHz\\[2pt]
\quad B. 5 kHz\\[2pt]
\quad C. 10 kHz\\[2pt]
\quad D. 20 kHz\\[2pt]
\vspace{4pt}
\textbf{Answer: C (10 kHz).} Nyquist sampling theorem requires $f_s \ge 2 f_{\max} = 2 \times 5\text{ kHz} = 10\text{ kHz}$ (or $10,000\text{ samples/second}$).
\end{quote}
\begin{quote}
\textbf{Practice Question 6: Which route is selected by Longest Prefix Match?}\\[4pt]
\quad A. The route with the shortest prefix\\[2pt]
\quad B. The route with the longest matching prefix\\[2pt]
\quad C. The route with the lowest IP address\\[2pt]
\quad D. Always the default route\\[2pt]
\vspace{4pt}
\textbf{Answer: B (The route with the longest matching prefix).} IP routers select the route with the most specific (longest) matching network prefix.
\end{quote}
\begin{quote}
\textbf{Practice Question 7: Which DNS record maps a hostname to an IPv6 address?}\\[4pt]
\quad A. A\\[2pt]
\quad B. AAAA\\[2pt]
\quad C. MX\\[2pt]
\quad D. PTR\\[2pt]
\vspace{4pt}
\textbf{Answer: B (AAAA).} An AAAA record maps a domain hostname to a 128-bit IPv6 address (compared to an A record for 32-bit IPv4).
\end{quote}
\begin{quote}
\textbf{Practice Question 8: Which DNS record identifies mail servers?}\\[4pt]
\quad A. NS\\[2pt]
\quad B. SOA\\[2pt]
\quad C. MX\\[2pt]
\quad D. CNAME\\[2pt]
\vspace{4pt}
\textbf{Answer: C (MX).} Mail Exchange (MX) records specify mail servers responsible for accepting email on behalf of the domain, ordered by priority preference.
\end{quote}
\begin{quote}
\textbf{Practice Question 9: Which port is conventionally used by DNS?}\\[4pt]
\quad A. 25\\[2pt]
\quad B. 53\\[2pt]
\quad C. 67\\[2pt]
\quad D. 110\\[2pt]
\vspace{4pt}
\textbf{Answer: B (Port 53).} DNS standardly operates on UDP port 53 for queries and TCP port 53 for zone transfers and responses exceeding UDP limits.
\end{quote}
\begin{quote}
\textbf{Practice Question 10: What does NXDOMAIN indicate?}\\[4pt]
\quad A. The DNS server is unreachable\\[2pt]
\quad B. The domain name does not exist according to DNS\\[2pt]
\quad C. The server refused the query\\[2pt]
\quad D. The domain has no IPv6 address\\[2pt]
\vspace{4pt}
\textbf{Answer: B (The domain name does not exist according to DNS).} NXDOMAIN (Non-Existent Domain) is an authoritative response stating that no record exists for the queried domain.
\end{quote}
\begin{quote}
\textbf{Practice Question 11: Which command traces the DNS lookup hierarchy?}\\[4pt]
\quad A. ipconfig /all\\[2pt]
\quad B. dig +trace example.com\\[2pt]
\quad C. ping example.com\\[2pt]
\quad D. netstat -r\\[2pt]
\vspace{4pt}
\textbf{Answer: B (dig +trace example.com).} The \texttt{+trace} option in \texttt{dig} initiates an iterative root-to-authoritative tree walk to trace the entire resolution chain.
\end{quote}
\begin{quote}
\textbf{Practice Question 12: Which NAT technique creates a permanent one-to-one address mapping?}\\[4pt]
\quad A. PAT\\[2pt]
\quad B. Dynamic NAT\\[2pt]
\quad C. Static NAT\\[2pt]
\quad D. CGNAT\\[2pt]
\vspace{4pt}
\textbf{Answer: C (Static NAT).} Static NAT maps one private IP address permanently to one dedicated public IP address in a $1:1$ relationship.
\end{quote}
\begin{quote}
\textbf{Practice Question 13: Which NAT technique uses a pool of public IP addresses?}\\[4pt]
\quad A. Static NAT\\[2pt]
\quad B. Dynamic NAT\\[2pt]
\quad C. PAT only\\[2pt]
\quad D. DNS\\[2pt]
\vspace{4pt}
\textbf{Answer: B (Dynamic NAT).} Dynamic NAT assigns temporary $1:1$ mappings from a configured pool of public IP addresses as outbound sessions are initiated.
\end{quote}
\begin{quote}
\textbf{Practice Question 14: Which technique allows many internal devices to share one public IPv4 address?}\\[4pt]
\quad A. Static NAT\\[2pt]
\quad B. Dynamic NAT\\[2pt]
\quad C. PAT\\[2pt]
\quad D. ARP\\[2pt]
\vspace{4pt}
\textbf{Answer: C (PAT).} Port Address Translation (NAT Overload) allows thousands of internal hosts to share one single public IP address using distinct transport source ports.
\end{quote}
\begin{quote}
\textbf{Practice Question 15: What does SNAT modify?}\\[4pt]
\quad A. Destination IP\\[2pt]
\quad B. Source IP\\[2pt]
\quad C. DNS record\\[2pt]
\quad D. Subnet mask\\[2pt]
\vspace{4pt}
\textbf{Answer: B (Source IP).} Source NAT (SNAT) translates the source IP address of outbound packets leaving private networks.
\end{quote}
\begin{quote}
\textbf{Practice Question 16: What does DNAT modify?}\\[4pt]
\quad A. Source IP\\[2pt]
\quad B. Source port only\\[2pt]
\quad C. Destination IP\\[2pt]
\quad D. MAC address\\[2pt]
\vspace{4pt}
\textbf{Answer: C (Destination IP).} Destination NAT (DNAT) translates the destination IP address of incoming packets (standard in port forwarding).
\end{quote}
\begin{quote}
\textbf{Practice Question 17: Which address range is reserved for carrier-grade NAT?}\\[4pt]
\quad A. 10.0.0.0/8\\[2pt]
\quad B. 127.0.0.0/8\\[2pt]
\quad C. 100.64.0.0/10\\[2pt]
\quad D. 192.168.0.0/16\\[2pt]
\vspace{4pt}
\textbf{Answer: C (100.64.0.0/10).} RFC 6598 reserves the shared address block \texttt{100.64.0.0/10} exclusively for Carrier-Grade NAT (CGNAT) deployments.
\end{quote}

\newpage

\subsection{Routing, Security, and Protocol Practice Questions}
\begin{quote}
\textbf{Practice Question 1: What is the primary purpose of STP?}\\[4pt]
\vspace{4pt}
\textbf{Answer: Prevent Layer 2 switching loops.} Spanning Tree Protocol (IEEE 802.1D/w) logically blocks redundant switch ports to prevent broadcast storms, frame duplication, and MAC table corruption.
\end{quote}
\begin{quote}
\textbf{Practice Question 2: Which switch becomes the STP root bridge?}\\[4pt]
\vspace{4pt}
\textbf{Answer: Switch with the lowest Bridge ID.} The Root Bridge election chooses the switch with the lowest Bridge ID (Priority + MAC address).
\end{quote}
\begin{quote}
\textbf{Practice Question 3: Which algorithm does OSPF use?}\\[4pt]
\vspace{4pt}
\textbf{Answer: Dijkstra.} OSPF utilizes Dijkstra's Shortest Path First (SPF) algorithm to calculate loop-free shortest paths from its Link-State Database.
\end{quote}
\begin{quote}
\textbf{Practice Question 4: Which IP protocol number does OSPF use?}\\[4pt]
\vspace{4pt}
\textbf{Answer: 89.} OSPF packets are encapsulated directly into IP datagrams using IP protocol number 89 (not TCP or UDP).
\end{quote}
\begin{quote}
\textbf{Practice Question 5: Which protocol is primarily used for routing between autonomous systems?}\\[4pt]
\vspace{4pt}
\textbf{Answer: BGP.} Border Gateway Protocol (BGP-4) is the standardized Exterior Gateway Protocol (EGP) connecting independent Autonomous Systems.
\end{quote}
\begin{quote}
\textbf{Practice Question 6: Which BGP attribute records the sequence of autonomous systems traversed?}\\[4pt]
\vspace{4pt}
\textbf{Answer: AS\_PATH.} The \texttt{AS\_PATH} attribute lists autonomous systems traversed by a route prefix and acts as the primary BGP inter-AS loop prevention mechanism.
\end{quote}
\begin{quote}
\textbf{Practice Question 7: What is a common cause of a broadcast storm?}\\[4pt]
\vspace{4pt}
\textbf{Answer: A Layer 2 switching loop.} Ethernet frames lack a Time-to-Live (TTL) field; redundant unblocked paths cause broadcast frames to circulate endlessly.
\end{quote}
\begin{quote}
\textbf{Practice Question 8: Which technology divides a switched LAN into separate broadcast domains?}\\[4pt]
\vspace{4pt}
\textbf{Answer: VLAN.} Virtual Local Area Networks (VLANs) segment a physical switch into distinct logical Layer 2 broadcast domains.
\end{quote}
\begin{quote}
\textbf{Practice Question 9: What does DORA stand for in DHCP?}\\[4pt]
\vspace{4pt}
\textbf{Answer: Discover, Offer, Request, Acknowledge.} The DORA sequence is the standardized 4-step DHCPv4 client-server lease negotiation exchange.
\end{quote}
\begin{quote}
\textbf{Practice Question 10: Which UDP port does a DHCP server use?}\\[4pt]
\vspace{4pt}
\textbf{Answer: 67.} DHCP servers listen on UDP port 67, while clients send/receive on UDP port 68.
\end{quote}
\begin{quote}
\textbf{Practice Question 11: A client receives an IPv4 address beginning with 169.254. What does this commonly indicate?}\\[4pt]
\vspace{4pt}
\textbf{Answer: APIPA/link-local self-assignment.} When DHCP allocation fails, Windows hosts self-assign an IPv4 address from \texttt{169.254.0.0/16} (APIPA).
\end{quote}
\begin{quote}
\textbf{Practice Question 12: Which device or function forwards DHCP broadcasts to a server on another subnet?}\\[4pt]
\vspace{4pt}
\textbf{Answer: DHCP relay agent.} A DHCP relay agent (configured via Cisco \texttt{ip helper-address}) forwards local client broadcast requests to unicast remote servers.
\end{quote}
\begin{quote}
\textbf{Practice Question 13: What is the primary function of an IDS?}\\[4pt]
\vspace{4pt}
\textbf{Answer: Detect suspicious activity and alert.} Intrusion Detection Systems operate out-of-band to inspect traffic, identify threats, and generate administrative alerts.
\end{quote}
\begin{quote}
\textbf{Practice Question 14: Which security system is commonly deployed inline to block malicious traffic?}\\[4pt]
\vspace{4pt}
\textbf{Answer: IPS.} Intrusion Prevention Systems are deployed inline directly in the traffic flow to actively drop malicious packets and terminate hostile connections.
\end{quote}
\begin{quote}
\textbf{Practice Question 15: What does a WAF primarily protect?}\\[4pt]
\vspace{4pt}
\textbf{Answer: Web applications.} Web Application Firewalls (WAF) inspect Layer 7 HTTP/HTTPS traffic to block SQL injection, Cross-Site Scripting (XSS), and web exploits.
\end{quote}
\begin{quote}
\textbf{Practice Question 16: Which principle gives users only the permissions necessary for their work?}\\[4pt]
\vspace{4pt}
\textbf{Answer: Least Privilege.} The Principle of Least Privilege mandates that accounts, processes, and systems receive only the minimum permissions required to complete tasks.
\end{quote}
\begin{quote}
\textbf{Practice Question 17: What is the primary purpose of a DMZ?}\\[4pt]
\vspace{4pt}
\textbf{Answer: Host public-facing services in a separated network segment.} A Demilitarized Zone (DMZ) buffers the internal trusted network by isolating externally accessible servers.
\end{quote}
\begin{quote}
\textbf{Practice Question 18: What does Stop-and-Wait ARQ do when an acknowledgment is lost?}\\[4pt]
\vspace{4pt}
\textbf{Answer: Waits for timeout and retransmits the frame.} In Stop-and-Wait ARQ, sender timer expiration triggers retransmission of the unacknowledged frame.
\end{quote}
\begin{quote}
\textbf{Practice Question 19: How many sequence-number values are needed for basic Stop-and-Wait ARQ?}\\[4pt]
\vspace{4pt}
\textbf{Answer: 2.} Because only a single frame is outstanding, a 1-bit sequence number ($0$ and $1$) completely suffices to distinguish new frames from duplicates.
\end{quote}
\begin{quote}
\textbf{Practice Question 20: If a Stop-and-Wait system has a = 2, what is its idealized utilization?}\\[4pt]
\vspace{4pt}
\textbf{Answer: 20\%.} Utilization $U = \frac{1}{1 + 2a} = \frac{1}{1 + 2(2)} = \frac{1}{5} = 20\%$.
\end{quote}

\subsection{Comprehensive Scenario-Based Engineering Analysis}
\subsubsection*{Scenario 1: Two switches are connected through redundant links}

\textbf{Question:} An organization connects two switches using two Ethernet links for redundancy. After connecting both links, users experience network instability and excessive broadcast traffic. What could be the problem, and how should it be resolved?

\textbf{Answer:}

The likely problem is a Layer 2 switching loop.

Both links may be forwarding simultaneously, allowing frames to circulate through redundant paths.

Recommended solution:

\begin{enumerate}
\item Enable STP or RSTP.
\item Verify the root bridge election.
\item Check Root Port and Designated Port roles.
\item Verify that redundant paths are handled correctly.
\item Inspect switch interface counters for broadcast storms and errors.
\end{enumerate}

RSTP is generally preferable in modern networks because it offers faster convergence than traditional STP.

\subsubsection*{Scenario 2: OSPF routers fail to become neighbors}

\textbf{Question:} Two routers connected to the same network are configured for OSPF but do not form a neighbor relationship. Identify possible causes.

\textbf{Answer:}

Possible causes include:

\begin{itemize}
\item Incorrect OSPF area configuration.
\item Mismatched Hello and Dead intervals.
\item Authentication mismatch.
\item Incorrect subnet configuration.
\item Area or network-type incompatibility.
\item Duplicate router IDs.
\item Interface configured as passive.
\item ACL or firewall blocking OSPF packets.
\item Mismatched stub-area flags where applicable.
\end{itemize}

Useful commands:

\begin{lstlisting}
show ip ospf neighbor
show ip ospf interface
show ip ospf
show ip protocols
\end{lstlisting}

OSPF uses IP protocol 89, so troubleshooting should not focus on TCP or UDP ports.

\subsubsection*{Scenario 3: An organization cannot access an external network after changing its subnet}

\textbf{Question:} After changing the LAN from \texttt{/24} to \texttt{/26}, some devices can communicate locally but cannot reach the Internet. What should the administrator check?

\textbf{Answer:}

Check:

\begin{enumerate}
\item Correct subnet mask on all devices.
\item Correct default gateway.
\item Router interface addressing.
\item DHCP scope and subnet mask options.
\item Routing table entries.
\item ACL and firewall rules.
\item NAT rules, if private addresses need Internet access.
\end{enumerate}

A subnet change affects which destinations devices consider local. An incorrect subnet mask can therefore cause devices to send traffic to the wrong next hop.

\subsubsection*{Scenario 4: Employees receive incorrect IP addresses}

\textbf{Question:} Employees suddenly receive IP addresses from an unknown DHCP server. What security problem may be occurring?

\textbf{Answer:}

A rogue DHCP server may be operating on the network.

It may provide incorrect:

\begin{itemize}
\item IP addresses.
\item Subnet masks.
\item Default gateways.
\item DNS server addresses.
\end{itemize}

This can disrupt connectivity or redirect network traffic.

Recommended controls:

\begin{itemize}
\item Enable DHCP snooping.
\item Mark legitimate DHCP-server-facing switch ports as trusted.
\item Keep client-facing ports untrusted.
\item Investigate the unauthorized DHCP server.
\item Review DHCP server logs and switch configuration.
\end{itemize}

\subsubsection*{Scenario 5: A network becomes slow even though bandwidth utilization appears low}

\textbf{Question:} What factors could explain the problem?

\textbf{Answer:}

Possible causes include:

\begin{itemize}
\item High latency.
\item Packet loss.
\item Excessive jitter.
\item DNS resolution delays.
\item Slow server response.
\item Duplex mismatch.
\item High CPU utilization on network devices.
\item Wireless interference.
\item Congestion on a different network segment.
\end{itemize}

Useful tests:

\begin{lstlisting}
ping destination
tracert destination
pathping destination
\end{lstlisting}

On Linux:

\begin{lstlisting}
ping destination
traceroute destination
mtr destination
\end{lstlisting}

A low average bandwidth utilization does not necessarily mean the network is healthy. Latency, packet loss, and application response times should also be examined.

\subsubsection*{Scenario 6: Internal users can access a web server, but external users cannot}

\textbf{Question:} What should the administrator investigate?

\textbf{Answer:}

Possible causes include:

\begin{itemize}
\item Missing or incorrect public DNS records.
\item Firewall blocking inbound HTTP/HTTPS.
\item Incorrect NAT or port-forwarding configuration.
\item Incorrect public IP address.
\item ISP filtering.
\item Web server not listening on the required port.
\item Incorrect routing.
\item CGNAT or double NAT restrictions.
\end{itemize}

A useful troubleshooting sequence is to verify public DNS resolution, the public address, firewall policy, NAT translation, routing, and the server's listening state.

\subsubsection*{Scenario 7: Two Internet providers advertise the same IP prefix}

\textbf{Question:} Which routing protocol can exchange these routes and choose between different Internet paths based on routing policy?

\textbf{Answer:}

BGP.

BGP can compare:

\begin{itemize}
\item LOCAL\_PREF.
\item AS\_PATH.
\item Origin.
\item MED.
\item NEXT\_HOP.
\item Other vendor-specific or configured path attributes.
\end{itemize}

It is a policy-based path-vector protocol, not simply a shortest-distance routing protocol.

\subsubsection*{Scenario 8: A server is receiving a large number of suspicious connection attempts}

\textbf{Question:} Which security technologies could help?

\textbf{Answer:}

\begin{itemize}
\item Firewall: Restricts unauthorized traffic.
\item IDS: Detects suspicious traffic and generates alerts.
\item IPS: Can block traffic matching prevention policies.
\item SIEM: Correlates events across network and server systems.
\item Rate limiting: Controls excessive connection attempts.
\end{itemize}

The correct combination depends on the traffic pattern and the network's security architecture.



\newpage
\begin{center}
\vspace{0.5em}
{\color{primary}\rule{0.45\textwidth}{1.2pt}\\[4pt]}
{\Large\bfseries\color{primary} Section XXII: Comprehensive Solved Numerical Problems and Subnetting Calculations\\[2pt]}
{\color{primary}\rule{0.45\textwidth}{1.2pt}}
\vspace{0.5em}
\end{center}

\section{Comprehensive Solved Numerical Problems and Subnetting Calculations}

\subsection{Part A --- IPv4 Addressing and Subnetting}
\subsubsection*{Problem 1: Identify the network address, broadcast address, and usable host range}

\textbf{Question:} Given the IPv4 address \texttt{192.168.10.75/24}, determine:

\begin{enumerate}
\item Subnet mask.
\item Network address.
\item Broadcast address.
\item First usable host.
\item Last usable host.
\item Number of usable hosts.
\end{enumerate}

\textbf{Solution}

Step 1: Convert \texttt{/24} into a subnet mask.

A \texttt{/24} prefix means 24 network bits and 8 host bits.



\[
\boxed{255.255.255.0}
\]



Step 2: Determine the network address.

The last octet is the host portion. Setting all host bits to zero gives:



\[
\boxed{192.168.10.0}
\]



Step 3: Determine the broadcast address.

Setting all host bits to one gives:



\[
\boxed{192.168.10.255}
\]



Step 4: Determine the usable host range.



\[
\text{First host}=192.168.10.1
\]





\[
\text{Last host}=192.168.10.254
\]



Step 5: Calculate the number of usable hosts.



\[
H=2^h-2
\]



Here, $h=8$.



\[
H=2^8-2=256-2
\]





\[
\boxed{254\text{ usable hosts}}
\]



\textbf{Final answer}

\begin{center}
\begin{tabularx}{\linewidth}{p{5.5cm} X}
\toprule
\textbf{Property} & \textbf{Result} \\
\midrule
IP address & 192.168.10.75/24 \\
Subnet mask & 255.255.255.0 \\
Network address & 192.168.10.0 \\
Broadcast address & 192.168.10.255 \\
First usable host & 192.168.10.1 \\
Last usable host & 192.168.10.254 \\
Usable hosts & 254 \\
\bottomrule
\end{tabularx}
\end{center}

\subsubsection*{Problem 2: Find the subnet information for a \texttt{/26} network}

\textbf{Question:} Given \texttt{192.168.1.130/26}, determine the subnet mask, network address, broadcast address, usable host range, and number of usable hosts.

\textbf{Solution}

Step 1: Calculate the subnet mask.



\[
/26=255.255.255.192
\]



Step 2: Calculate the block size.

The interesting octet is the fourth octet.



\[
\text{Block size}=256-192=64
\]



Step 3: Identify subnet ranges.

The fourth octet increases in increments of 64:

\begin{center}
\begin{tabularx}{\linewidth}{p{5.5cm} X}
\toprule
\textbf{Subnet} & \textbf{Address range} \\
\midrule
1 & 192.168.1.0--63 \\
2 & 192.168.1.64--127 \\
3 & 192.168.1.128--191 \\
4 & 192.168.1.192--255 \\
\bottomrule
\end{tabularx}
\end{center}

Since 130 falls between 128 and 191:



\[
\text{Network address}=192.168.1.128
\]





\[
\text{Broadcast address}=192.168.1.191
\]



Step 4: Find usable host range.

\begin{itemize}
\item First usable: \texttt{192.168.1.129}
\item Last usable: \texttt{192.168.1.190}
\end{itemize}

Step 5: Calculate usable hosts.



\[
h=32-26=6
\]





\[
H=2^6-2=62
\]





\[
\boxed{62\text{ usable hosts}}
\]



\subsubsection*{Problem 3: Determine the required subnet prefix}

\textbf{Question:} An organization requires at least 100 usable IPv4 host addresses in each subnet. What is the minimum number of host bits and the appropriate CIDR prefix?

\textbf{Solution}

The usable host formula is:



\[
H=2^h-2
\]



We need:



\[
2^h-2\geq100
\]



Try $h=6$:



\[
2^6-2=62
\]



Insufficient.

Try $h=7$:



\[
2^7-2=126
\]



Sufficient.

Therefore:



\[
h=7
\]



The prefix length is:



\[
32-7=25
\]





\[
\boxed{/25}
\]



Subnet mask:



\[
\boxed{255.255.255.128}
\]



Each \texttt{/25} subnet supports 126 usable hosts.

\subsubsection*{Problem 4: Calculate the number of subnets}

\textbf{Question:} A network administrator receives \texttt{192.168.10.0/24} and wants to divide it into eight equal-sized subnets.

Determine:

\begin{enumerate}
\item New prefix length.
\item New subnet mask.
\item Number of host bits.
\item Number of usable hosts per subnet.
\end{enumerate}

\textbf{Solution}

Step 1: Determine borrowed bits.



\[
2^n=8
\]





\[
n=3
\]



Step 2: Calculate the new prefix.



\[
24+3=27
\]





\[
\boxed{/27}
\]



Step 3: Calculate subnet mask.



\[
\boxed{255.255.255.224}
\]



Step 4: Calculate remaining host bits.



\[
h=32-27=5
\]



Step 5: Calculate usable hosts.



\[
H=2^5-2=30
\]



\textbf{Final result}

\begin{itemize}
\item Number of subnets: 8.
\item Prefix: \texttt{/27}.
\item Subnet mask: \texttt{255.255.255.224}.
\item Usable hosts per subnet: 30.
\end{itemize}

\subsubsection*{Problem 5: VLSM subnet allocation}

\textbf{Question:} An organization receives the network \texttt{192.168.10.0/24} and requires four subnets:

\begin{center}
\begin{tabularx}{\linewidth}{p{5.5cm} X}
\toprule
\textbf{Department} & \textbf{Required hosts} \\
\midrule
Administration & 100 \\
Finance & 50 \\
HR & 25 \\
IT & 10 \\
\bottomrule
\end{tabularx}
\end{center}

Allocate the network using VLSM.

\textbf{Solution}

VLSM allocation should begin with the largest subnet requirement.

Step 1: Sort by descending host requirement.

100, 50, 25, 10.

Step 2: Calculate the required prefix for each subnet.

\begin{center}
\begin{tabularx}{\linewidth}{p{3cm} p{3cm} p{2.5cm} X}
\toprule
\textbf{Required hosts} & \textbf{Minimum host bits} & \textbf{Prefix} & \textbf{Usable capacity} \\
\midrule
100 & 7 & /25 & 126 \\
50 & 6 & /26 & 62 \\
25 & 5 & /27 & 30 \\
10 & 4 & /28 & 14 \\
\bottomrule
\end{tabularx}
\end{center}

Step 3: Allocate sequentially.

\textbf{Administration: /25}

\begin{itemize}
\item Network: \texttt{192.168.10.0/25}
\item Usable range: \texttt{192.168.10.1--192.168.10.126}
\item Broadcast: \texttt{192.168.10.127}
\end{itemize}

\textbf{Finance: /26}

\begin{itemize}
\item Network: \texttt{192.168.10.128/26}
\item Usable range: \texttt{192.168.10.129--192.168.10.190}
\item Broadcast: \texttt{192.168.10.191}
\end{itemize}

\textbf{HR: /27}

\begin{itemize}
\item Network: \texttt{192.168.10.192/27}
\item Usable range: \texttt{192.168.10.193--192.168.10.222}
\item Broadcast: \texttt{192.168.10.223}
\end{itemize}

\textbf{IT: /28}

\begin{itemize}
\item Network: \texttt{192.168.10.224/28}
\item Usable range: \texttt{192.168.10.225--192.168.10.238}
\item Broadcast: \texttt{192.168.10.239}
\end{itemize}

\textbf{Final allocation}

\begin{center}
\begin{tabularx}{\linewidth}{p{3cm} p{3cm} p{2.5cm} X}
\toprule
\textbf{Department} & \textbf{Network} & \textbf{Usable range} & \textbf{Broadcast} \\
\midrule
Administration & 192.168.10.0/25 & .1--.126 & .127 \\
Finance & 192.168.10.128/26 & .129--.190 & .191 \\
HR & 192.168.10.192/27 & .193--.222 & .223 \\
IT & 192.168.10.224/28 & .225--.238 & .239 \\
\bottomrule
\end{tabularx}
\end{center}

Remaining unallocated range:

\texttt{192.168.10.240--192.168.10.255}

This provides 16 additional addresses for future allocation, although they would need to be organized into appropriately aligned subnets.

\subsubsection*{Problem 6: Find the number of hosts and addresses in a CIDR block}

\textbf{Question:} How many total IPv4 addresses and usable host addresses exist in a \texttt{/20} network?

\textbf{Solution}

Number of host bits:



\[
h=32-20=12
\]



Total addresses:



\[
2^{12}=4096
\]



Usable host addresses in a conventional subnet:



\[
4096-2=4094
\]





\[
\boxed{4096\text{ total addresses, }4094\text{ usable hosts}}
\]



The usual subtraction of two applies to conventional IPv4 subnets, not every possible prefix. For example, \texttt{/31} is used for point-to-point links under RFC 3021, and \texttt{/32} identifies one host route.

\subsubsection*{Problem 7: Supernetting and route aggregation}

\textbf{Question:} Aggregate the following four contiguous \texttt{/24} networks into one supernet:

\begin{itemize}
\item \texttt{192.168.0.0/24}
\item \texttt{192.168.1.0/24}
\item \texttt{192.168.2.0/24}
\item \texttt{192.168.3.0/24}
\end{itemize}

\textbf{Solution}

There are four contiguous \texttt{/24} networks.



\[
4=2^2
\]



Therefore, two network bits can be removed from the prefix:



\[
24-2=22
\]



The four networks share the first 22 bits.

Aggregated route:



\[
\boxed{192.168.0.0/22}
\]



Subnet mask:



\[
\boxed{255.255.252.0}
\]



Address range:

\texttt{192.168.0.0--192.168.3.255}

\textbf{Important:} Route aggregation requires contiguous, correctly aligned networks. Four arbitrary \texttt{/24} networks cannot necessarily be represented by one \texttt{/22}.

\subsection{Part B --- Binary, IP Addressing, and Routing Mathematics}
\subsubsection*{Problem 8: Convert an IPv4 address into binary}

\textbf{Question:} Convert \texttt{192.168.1.10} into binary notation.

\textbf{Solution}

Convert each octet independently.

\begin{center}
\begin{tabularx}{\linewidth}{p{5.5cm} X}
\toprule
\textbf{Decimal} & \textbf{Binary} \\
\midrule
192 & 11000000 \\
168 & 10101000 \\
1 & 00000001 \\
10 & 00001010 \\
\bottomrule
\end{tabularx}
\end{center}

Therefore:



\[
\boxed{11000000.10101000.00000001.00001010}
\]



Each octet contains eight bits.

\subsubsection*{Problem 9: Perform a bitwise subnet-mask operation}

\textbf{Question:} Find the network address of:

\begin{lstlisting}
IP address: 192.168.1.130
Subnet mask: 255.255.255.192
\end{lstlisting}

\textbf{Solution}

The network address is calculated using a bitwise AND operation:



\[
\text{Network}=\text{IP}\ AND\ \text{Mask}
\]



For the fourth octet:

\begin{lstlisting}
IP:      10000010  (130)
Mask:    11000000  (192)
         --------
AND:     10000000  (128)
\end{lstlisting}

The first three octets remain unchanged.



\[
\boxed{192.168.1.128}
\]



\subsubsection*{Problem 10: Longest Prefix Match}

\textbf{Question:} A router contains the following routing entries:

\begin{center}
\begin{tabularx}{\linewidth}{p{5.5cm} X}
\toprule
\textbf{Destination prefix} & \textbf{Next hop} \\
\midrule
10.0.0.0/8 & Router A \\
10.1.0.0/16 & Router B \\
10.1.2.0/24 & Router C \\
10.1.2.128/25 & Router D \\
0.0.0.0/0 & Router E \\
\bottomrule
\end{tabularx}
\end{center}

Determine the next hop for:

\begin{enumerate}
\item \texttt{10.1.2.200}
\item \texttt{10.1.2.100}
\item \texttt{10.5.1.10}
\item \texttt{8.8.8.8}
\end{enumerate}

\textbf{Solution}

Longest Prefix Match selects the matching route with the largest prefix length.

\begin{enumerate}
\item \texttt{10.1.2.200} matches \texttt{/8}, \texttt{/16}, \texttt{/24}, and \texttt{/25}. Select \texttt{/25} $\to$ \textbf{Router D}.
\item \texttt{10.1.2.100} matches \texttt{/8}, \texttt{/16}, and \texttt{/24}, but not \texttt{10.1.2.128/25} $\to$ \textbf{Router C}.
\item \texttt{10.5.1.10} matches only \texttt{/8} $\to$ \textbf{Router A}.
\item \texttt{8.8.8.8} matches no specific route $\to$ default \texttt{/0} $\to$ \textbf{Router E}.
\end{enumerate}

\subsubsection*{Problem 11: Calculate the administrative distance preference}

\textbf{Question:} A Cisco router learns the same destination through:

\begin{itemize}
\item Connected route: AD 0.
\item Static route: AD 1.
\item OSPF: AD 110.
\item RIP: AD 120.
\item eBGP: AD 20.
\end{itemize}

Which route source is preferred if all routes are otherwise eligible and have the same destination prefix?

\textbf{Solution}

Cisco administrative distance is a lower-is-better route-source preference.

Order:

\begin{enumerate}
\item Connected: 0.
\item Static: 1.
\item eBGP: 20.
\item OSPF: 110.
\item RIP: 120.
\end{enumerate}

The connected route is selected.

\textbf{Important distinction:} Administrative distance compares route sources on a router. It is not the same as longest prefix match. A more-specific route can be selected over a less-specific route even if its source has a higher administrative distance.

\subsection{Part C --- Transmission, Bandwidth, Latency, and Data Rate}
\subsubsection*{Problem 12: Calculate transmission delay}

\textbf{Question:} A device transmits a packet of 12,000 bits over a link with a bandwidth of 6 Mbps. Calculate the transmission delay.

\textbf{Solution}

Formula:



\[
T_t=\frac{L}{R}
\]



Given:

\begin{itemize}
\item Packet size $L=12,000$ bits.
\item Bandwidth $R=6,000,000$ bits/s.
\end{itemize}



\[
T_t=\frac{12,000}{6,000,000}
\]





\[
T_t=0.002s
\]





\[
\boxed{2ms}
\]



\subsubsection*{Problem 13: Calculate propagation delay}

\textbf{Question:} A signal travels through a 1,000 km fiber-optic link at a propagation speed of $2\times10^8$ m/s. Calculate the one-way propagation delay.

\textbf{Solution}

Formula:



\[
T_p=\frac{d}{v}
\]



Convert distance:



\[
1000km=1,000,000m
\]



Substitute:



\[
T_p=\frac{1,000,000}{2\times10^8}
\]





\[
T_p=0.005s
\]





\[
\boxed{5ms}
\]



\subsubsection*{Problem 14: Calculate total network delay}

\textbf{Question:} A packet travels through a network with the following characteristics:

\begin{center}
\begin{tabularx}{\linewidth}{p{5.5cm} X}
\toprule
\textbf{Delay component} & \textbf{Value} \\
\midrule
Processing delay & 0.5 ms \\
Queuing delay & 1.5 ms \\
Transmission delay & 2 ms \\
Propagation delay & 5 ms \\
\bottomrule
\end{tabularx}
\end{center}

Calculate the total nodal delay.

\textbf{Solution}

Total nodal delay:



\[
T_{total}=T_{proc}+T_{queue}+T_t+T_p
\]





\[
=0.5+1.5+2+5
\]





\[
\boxed{9ms}
\]



For a path containing multiple routers and links, the total end-to-end delay is obtained by summing the relevant delays across all nodes and links.

\subsubsection*{Problem 15: Calculate round-trip time (RTT)}

\textbf{Question:} A packet takes 25 ms to travel from a client to a server and 35 ms to return. Calculate the RTT.

\textbf{Solution}



\[
RTT=T_{forward}+T_{return}
\]





\[
RTT=25+35
\]





\[
\boxed{60ms}
\]



The path is asymmetric, so RTT is not twice the forward delay.

\subsubsection*{Problem 16: Calculate bandwidth-delay product (BDP)}

\textbf{Question:} A network connection has a bandwidth of 100 Mbps and an RTT of 40 ms. Calculate the bandwidth-delay product in bits and bytes.

\textbf{Solution}

Formula:



\[
BDP=Bandwidth\times RTT
\]



Convert RTT:



\[
40ms=0.04s
\]



Substitute:



\[
BDP=100,000,000\times0.04
\]





\[
BDP=4,000,000\ bits
\]



Convert into bytes:



\[
\frac{4,000,000}{8}=500,000\ bytes
\]





\[
\boxed{4Mb=500,000\ bytes}
\]



This is the approximate amount of data that can be in flight to fill the link under the given conditions.

\subsubsection*{Problem 17: Calculate throughput}

\textbf{Question:} A file of 150 MB is transferred over a network in 30 seconds. Calculate the average throughput in MB/s and Mbps.

\textbf{Solution}



\[
Throughput=\frac{Data\ transferred}{Time}
\]





\[
=\frac{150MB}{30s}
\]





\[
=5MB/s
\]



Convert to bits per second:



\[
5\times8=40Mbps
\]





\[
\boxed{5MB/s=40Mbps}
\]



This assumes decimal MB and Mbps units. Real application throughput may be lower than the nominal link bandwidth due to protocol overhead and other factors.

\subsubsection*{Problem 18: Calculate transmission time for a file}

\textbf{Question:} How long will it take to transfer a 200 MB file over a 50 Mbps connection, assuming the full bandwidth is available?

\textbf{Solution}

Convert file size into bits using decimal units:



\[
200MB=200\times8\times10^6\ bits
\]





\[
=1,600,000,000\ bits
\]



Transmission time:



\[
T=\frac{1,600,000,000}{50,000,000}
\]





\[
\boxed{32s}
\]



Actual transfer time will generally be higher because of protocol overhead, congestion, and other limitations.

\subsection{Part D --- Modulation and Digital Communication}
\subsubsection*{Problem 19: Calculate bit rate using baud rate}

\textbf{Question:} A digital communication system uses 16-QAM modulation at 2,400 baud. Calculate the bit rate.

\textbf{Solution}

For M-ary modulation:



\[
R_b=R_s\log_2M
\]



Given:

\begin{itemize}
\item $R_s=2400$ baud.
\item $M=16$.
\end{itemize}

Bits per symbol:



\[
\log_2 16=4
\]



Therefore:



\[
R_b=2400\times4
\]





\[
\boxed{9600bps}
\]



\subsubsection*{Problem 20: Determine bits per symbol}

\textbf{Question:} How many bits are represented by each symbol in a 64-QAM communication system?

\textbf{Solution}



\[
Bits/Symbol=\log_2M
\]





\[
=\log_2 64
\]





\[
\boxed{6\ bits/symbol}
\]



\subsubsection*{Problem 21: Calculate PCM bit rate}

\textbf{Question:} A PCM system samples an audio signal at 8,000 samples per second and uses 8 bits per sample. Calculate the bit rate.

\textbf{Solution}



\[
R_b=Sampling\ Rate\times Bits/Sample
\]





\[
=8000\times8
\]





\[
\boxed{64,000bps=64kbps}
\]



\subsubsection*{Problem 22: Apply the Nyquist sampling theorem}

\textbf{Question:} An analog signal contains frequencies up to 4 kHz. What is the minimum theoretical sampling rate required to reconstruct the signal under ideal Nyquist sampling conditions?

\textbf{Solution}

Nyquist sampling condition:



\[
f_s\geq2f_{max}
\]





\[
f_s\geq2\times4kHz
\]





\[
\boxed{f_s\geq8kHz}
\]



In practice, systems may use higher sampling rates and anti-aliasing filters.

\subsubsection*{Problem 23: Calculate AM modulation index}

\textbf{Question:} An AM signal has a maximum envelope amplitude of 12 V and a minimum envelope amplitude of 4 V. Calculate its modulation index.

\textbf{Solution}



\[
m=\frac{E_{max}-E_{min}}{E_{max}+E_{min}}
\]



Substitute:



\[
m=\frac{12-4}{12+4}
\]





\[
m=\frac{8}{16}=0.5
\]



Percentage modulation:



\[
0.5\times100=50\%
\]





\[
\boxed{50\%}
\]



\subsubsection*{Problem 24: Calculate FM bandwidth using Carson's rule}

\textbf{Question:} An FM signal has a maximum frequency deviation of 75 kHz and a maximum modulating frequency of 15 kHz. Estimate the required bandwidth using Carson's rule.

\textbf{Solution}

Carson's rule:



\[
BW\approx2(\Delta f+f_m)
\]





\[
=2(75+15)
\]





\[
=180kHz
\]





\[
\boxed{180kHz}
\]



\subsection{Part E --- Routing Protocols and Switching}
\subsubsection*{Problem 25: STP root bridge election}

\textbf{Question:} Three switches have the following Bridge IDs:

\begin{center}
\begin{tabularx}{\linewidth}{p{3.5cm} p{4.5cm} X}
\toprule
\textbf{Switch} & \textbf{Priority} & \textbf{MAC address} \\
\midrule
S1 & 32768 & 00:00:00:00:00:10 \\
S2 & 4096 & 00:00:00:00:00:20 \\
S3 & 8192 & 00:00:00:00:00:05 \\
\bottomrule
\end{tabularx}
\end{center}

Which switch becomes the root bridge?

\textbf{Solution}

STP selects the switch with the lowest Bridge ID.

Compare bridge priorities:

\begin{itemize}
\item S1: 32768.
\item S2: 4096.
\item S3: 8192.
\end{itemize}

S2 has the lowest priority value.

Therefore:



\[
\boxed{S2\ is\ the\ Root\ Bridge}
\]



The MAC addresses are used to break a tie when the priority values are equal.

\subsubsection*{Problem 26: Select the OSPF shortest path}

\textbf{Question:} Router A has three paths to Router D:

\begin{center}
\begin{tabularx}{\linewidth}{p{5.5cm} X}
\toprule
\textbf{Path} & \textbf{Link costs} \\
\midrule
A $\to$ B $\to$ D & 8 + 12 \\
A $\to$ C $\to$ D & 5 + 7 \\
A $\to$ E $\to$ D & 10 + 10 \\
\bottomrule
\end{tabularx}
\end{center}

Which path will OSPF prefer?

\textbf{Solution}

Calculate each total cost:



\[
C_{ABD}=8+12=20
\]





\[
C_{ACD}=5+7=12
\]





\[
C_{AED}=10+10=20
\]



The lowest cost is 12.



\[
\boxed{A\rightarrow C\rightarrow D}
\]



\subsubsection*{Problem 27: OSPF interface cost}

\textbf{Question:} A router uses a reference bandwidth of 100 Mbps. Calculate the OSPF cost for a 20 Mbps interface.

\textbf{Solution}



\[
Cost=\frac{Reference\ Bandwidth}{Interface\ Bandwidth}
\]





\[
=\frac{100}{20}
\]





\[
\boxed{5}
\]



This is the simplified cost calculation using the stated reference bandwidth.

\subsubsection*{Problem 28: BGP path selection}

\textbf{Question:} A router receives two otherwise comparable routes to the same destination:

\begin{center}
\begin{tabularx}{\linewidth}{p{3.5cm} p{4.5cm} X}
\toprule
\textbf{Attribute} & \textbf{Route A} & \textbf{Route B} \\
\midrule
LOCAL\_PREF & 150 & 100 \\
AS\_PATH length & 3 & 2 \\
\bottomrule
\end{tabularx}
\end{center}

Which route is preferred?

\textbf{Solution}

BGP commonly considers LOCAL\_PREF before AS\_PATH length.

A higher LOCAL\_PREF is preferred.



\[
150>100
\]



Therefore:



\[
\boxed{Route\ A}
\]



A shorter AS\_PATH does not automatically override a higher LOCAL\_PREF.

\subsubsection*{Problem 29: Determine the BGP transport port}

\textbf{Question:} Which transport protocol and port are used to establish a standard BGP session?

\textbf{Solution}



\[
\boxed{TCP\ port\ 179}
\]



BGP uses TCP to provide reliable transport for routing information.

\subsubsection*{Problem 30: Identify the routing protocol from a scenario}

\textbf{Question:} An ISP exchanges routing information with another ISP belonging to a different autonomous system. Which routing protocol and BGP type are appropriate?

\textbf{Solution}

\begin{itemize}
\item Routing protocol: BGP.
\item Type: eBGP.
\item Transport: TCP port 179.
\end{itemize}

If the routers belong to the same AS, the corresponding BGP relationship is iBGP.

\subsection{Part F --- NAT, DHCP, DNS, and Network Services}
\subsubsection*{Problem 31: Determine the NAT translation}

\textbf{Question:} A private computer has the following address:

\begin{lstlisting}
Private IP: 192.168.1.10
Private port: 50000
\end{lstlisting}

It accesses a web server using HTTPS. The NAT router translates the connection to:

\begin{lstlisting}
Public IP: 203.0.113.10
Translated port: 40001
\end{lstlisting}

What translation does PAT perform?

\textbf{Solution}

Before translation:

\begin{lstlisting}
Source:      192.168.1.10:50000
Destination: Web server IP:443
\end{lstlisting}

After translation:

\begin{lstlisting}
Source:      203.0.113.10:40001
Destination: Web server IP:443
\end{lstlisting}

PAT changes the source IP address and source port, maintaining a translation-table entry so return traffic can be delivered to the correct private device.

\subsubsection*{Problem 32: Calculate required public IP addresses for Static NAT}

\textbf{Question:} An organization needs permanent one-to-one Static NAT mappings for 30 internal devices. How many public IPv4 addresses are required?

\textbf{Solution}

Static NAT uses a permanent one-to-one mapping.



\[
Public\ IPs=Internal\ devices
\]





\[
\boxed{30\ public\ IP\ addresses}
\]



This assumes each device receives its own dedicated public IPv4 mapping.

\subsubsection*{Problem 33: Determine the DHCP message sequence}

\textbf{Question:} A client joins a network and requests an IPv4 address. Arrange the DHCP messages in the correct order.

\textbf{Solution}



\[
\boxed{Discover\rightarrow Offer\rightarrow Request\rightarrow ACK}
\]



The sequence is known as DORA.

\subsubsection*{Problem 34: DHCP scope capacity}

\textbf{Question:} A DHCP server has a scope from \texttt{192.168.1.100} to \texttt{192.168.1.150}, inclusive. How many addresses are available for assignment?

\textbf{Solution}



\[
150-100+1=51
\]





\[
\boxed{51\ addresses}
\]



If 8 addresses are reserved or excluded, the remaining dynamic pool contains:



\[
51-8=\boxed{43\ addresses}
\]



\subsubsection*{Problem 35: DNS record identification}

\textbf{Question:} Identify the correct DNS record type for each requirement.

\begin{center}
\begin{tabularx}{\linewidth}{p{5.5cm} X}
\toprule
\textbf{Requirement} & \textbf{Record} \\
\midrule
Map a hostname to an IPv4 address & A \\
Map a hostname to an IPv6 address & AAAA \\
Identify a mail server & MX \\
Identify an alias & CNAME \\
Identify authoritative name servers & NS \\
Reverse DNS lookup & PTR \\
Store domain-related text information & TXT \\
Specify a certificate authority policy & CAA \\
\bottomrule
\end{tabularx}
\end{center}

\subsubsection*{Problem 36: DNS troubleshooting scenario}

\textbf{Question:} A user can access a website using its IP address but cannot access it using its domain name. What is the most likely problem, and how would you troubleshoot it?

\textbf{Solution}

The issue is likely related to DNS resolution, although the application may also depend on hostname-based routing or TLS configuration.

A practical troubleshooting sequence:

1\. Test the domain:

\begin{lstlisting}
nslookup example.com
\end{lstlisting}

2\. Query a specific resolver:

\begin{lstlisting}
nslookup example.com 8.8.8.8
\end{lstlisting}

3\. Inspect DNS records:

\begin{lstlisting}
dig A example.com
dig AAAA example.com
\end{lstlisting}

\begin{enumerate}
\item Check the configured DNS server.
\item Check DNS cache and relevant host-file entries.
\item Verify authoritative DNS records and delegation if necessary.
\end{enumerate}

If the default resolver fails but another resolver returns the correct address, the problem may be with the configured resolver or its network path.

\subsection{Part G --- Firewall, Security, and Network Troubleshooting}
\subsubsection*{Problem 37: Identify the appropriate security technology}

\textbf{Question:} Match each situation with the appropriate security mechanism.

\begin{center}
\begin{tabularx}{\linewidth}{p{5.5cm} X}
\toprule
\textbf{Situation} & \textbf{Appropriate mechanism} \\
\midrule
Block unauthorized incoming TCP connections & Firewall \\
Detect suspicious network traffic and alert administrators & IDS \\
Automatically block malicious network packets & IPS \\
Protect a web application from malicious HTTP requests & WAF \\
Prevent unauthorized DHCP servers & DHCP Snooping \\
Detect and prevent ARP spoofing & Dynamic ARP Inspection \\
Protect remote connections over an untrusted network & VPN \\
Collect and correlate security logs & SIEM \\
\bottomrule
\end{tabularx}
\end{center}

\subsubsection*{Problem 38: Firewall rule analysis}

\textbf{Question:} Consider these firewall rules, evaluated from top to bottom:

\begin{center}
\begin{tabularx}{\linewidth}{c p{2.5cm} p{2.5cm} p{2cm} X}
\toprule
\textbf{Rule} & \textbf{Source} & \textbf{Destination} & \textbf{Port} & \textbf{Action} \\
\midrule
1 & Any & Web Server & TCP 443 & Allow \\
2 & Any & Database Server & TCP 3306 & Deny \\
3 & Internal LAN & Internet & TCP 80 & Allow \\
4 & Any & Any & Any & Deny \\
\bottomrule
\end{tabularx}
\end{center}

Determine the result of the following connections:

\begin{enumerate}
\item External client to web server on TCP 443.
\item External client to database server on TCP 3306.
\item Internal client to Internet on TCP 80.
\item External client to internal server on TCP 22.
\end{enumerate}

\textbf{Solution}

\begin{enumerate}
\item Rule 1 matches $\to$ Allow.
\item Rule 2 matches $\to$ Deny.
\item Rule 3 matches $\to$ Allow.
\item Rule 4 matches $\to$ Deny.
\end{enumerate}

This assumes the rules are applied in the stated order and that no other policy or stateful connection rule changes the outcome.

\subsubsection*{Problem 39: Diagnose a network connectivity problem}

\textbf{Question:} A computer has the following configuration:

\begin{lstlisting}
IP Address:      192.168.10.25
Subnet Mask:     255.255.255.0
Default Gateway: 192.168.20.1
DNS Server:      8.8.8.8
\end{lstlisting}

The computer cannot communicate with devices outside its local subnet. Identify the configuration problem.

\textbf{Solution}

The computer belongs to:



\[
192.168.10.0/24
\]



Its default gateway is:



\[
192.168.20.1
\]



The gateway is outside the computer's local subnet.

The gateway should normally have an address within the same directly connected subnet, such as:



\[
\boxed{192.168.10.1}
\]



Assuming that is the actual router interface address, correcting the default gateway should restore routed connectivity.

\subsubsection*{Problem 40: Diagnose an APIPA address}

\textbf{Question:} A Windows computer displays the address \texttt{169.254.12.30}. It cannot access the Internet. What is the likely cause?

\textbf{Solution}

The address belongs to the IPv4 link-local APIPA range:



\[
169.254.0.0/16
\]



The likely problem is failure to obtain a DHCP lease.

Possible causes:

\begin{itemize}
\item DHCP server unavailable.
\item DHCP scope exhausted.
\item Incorrect VLAN.
\item DHCP relay failure.
\item Physical connectivity problem.
\item DHCP traffic blocked.
\end{itemize}

Useful commands:

\begin{lstlisting}
ipconfig /all
ipconfig /release
ipconfig /renew
\end{lstlisting}

\subsection{Part H --- Stop-and-Wait ARQ and Reliable Transmission}
\subsubsection*{Problem 41: Calculate Stop-and-Wait utilization}

\textbf{Question:} A Stop-and-Wait ARQ system transmits 1,000-bit frames over a 1 Mbps link. The one-way propagation delay is 5 ms. Ignore ACK transmission and processing time. Calculate:

\begin{enumerate}
\item Frame transmission time.
\item Total cycle time.
\item Link utilization.
\item Effective throughput.
\end{enumerate}

\textbf{Solution}

Given:



\[
L=1000\ bits
\]





\[
R=1,000,000\ bps
\]





\[
T_p=5ms
\]



Step 1: Transmission time.



\[
T_t=\frac{L}{R}
\]





\[
=\frac{1000}{1,000,000} =0.001s
\]





\[
\boxed{T_t=1ms}
\]



Step 2: Cycle time.



\[
T_{cycle}=T_t+2T_p
\]





\[
=1+2(5)
\]





\[
\boxed{11ms}
\]



Step 3: Utilization.



\[
U=\frac{T_t}{T_t+2T_p}
\]





\[
=\frac{1}{11}
\]





\[
\boxed{9.09\%}
\]



Step 4: Effective throughput.



\[
Throughput=R\times U
\]





\[
=1,000,000\times0.0909
\]





\[
\boxed{90.9kbps}
\]



\subsubsection*{Problem 42: Calculate Stop-and-Wait efficiency using $a$}

\textbf{Question:} A Stop-and-Wait system has:



\[
a=\frac{T_p}{T_t}=3
\]



Calculate its utilization.

\textbf{Solution}



\[
U=\frac{1}{1+2a}
\]





\[
=\frac{1}{1+2(3)}
\]





\[
=\frac{1}{7}
\]





\[
\boxed{14.29\%}
\]



\subsubsection*{Problem 43: Calculate the minimum window size needed to keep a link busy}

\textbf{Question:} A network has:

\begin{itemize}
\item Bandwidth = 10 Mbps.
\item RTT = 40 ms.
\item Frame size = 1,000 bits.
\end{itemize}

Approximately how many frames must be in flight to fill the link?

\textbf{Solution}

First, calculate the bandwidth-delay product:



\[
BDP=10,000,000\times0.04
\]





\[
BDP=400,000\ bits
\]



Divide by the frame size:



\[
N=\frac{400,000}{1000}
\]





\[
\boxed{N=400\ frames}
\]



Approximately 400 frames must be in flight to fill the link, ignoring headers, ACKs, and other overhead.

This demonstrates why sliding-window protocols are more efficient than Stop-and-Wait on high-bandwidth, high-latency connections.

\subsubsection*{Problem 44: Calculate the number of retransmissions}

\textbf{Question:} A sender transmits 100 frames. Five frames are lost, and each lost frame is successfully retransmitted once. How many total frame transmission attempts occur?

\textbf{Solution}

Original transmissions:



\[
100
\]



Retransmissions:



\[
5
\]



Total attempts:



\[
100+5=\boxed{105}
\]



This assumes no retransmission itself is lost or corrupted.

\subsection{Part I --- Additional High-Value Numerical Problems}
\subsubsection*{Problem 45: Calculate the maximum theoretical data transfer rate}

\textbf{Question:} A network interface has a bandwidth of 1 Gbps. Express its theoretical data rate in MB/s.

\textbf{Solution}



\[
1Gbps=1,000,000,000\ bits/s
\]



Convert bits to bytes:



\[
\frac{1,000,000,000}{8} =125,000,000\ bytes/s
\]





\[
\boxed{125MB/s}
\]



This uses decimal MB, where $1MB=1,000,000$ bytes.

\subsubsection*{Problem 46: Calculate the time needed to transmit a packet over multiple links}

\textbf{Question:} A 10,000-bit packet travels over three store-and-forward links, each operating at 5 Mbps. Ignore propagation, processing, and queuing delays. Calculate the total transmission delay.

\textbf{Solution}

Transmission time over one link:



\[
T_t=\frac{10,000}{5,000,000}
\]





\[
=0.002s=2ms
\]



With store-and-forward switching, the entire packet must be received before being transmitted over the next link.

Therefore:



\[
T=3\times2
\]





\[
\boxed{6ms}
\]



This assumes identical link speeds and no additional delays.

\subsubsection*{Problem 47: Calculate network availability}

\textbf{Question:} A network service has an availability of 99.9\%. Calculate the expected downtime over a 365-day year.

\textbf{Solution}

Unavailable percentage:



\[
100\%-99.9\%=0.1\%
\]



Total annual time:



\[
365\times24\times60=525,600\ minutes
\]



Downtime:



\[
525,600\times0.001=525.6\ minutes
\]



Convert to hours:



\[
\frac{525.6}{60}=8.76\ hours
\]





\[
\boxed{8\ hours\ 45\ minutes\ 36\ seconds}
\]



\subsubsection*{Problem 48: Calculate packet loss percentage}

\textbf{Question:} A network sends 5,000 packets, but 75 packets fail to reach their destination. Calculate the packet loss percentage.

\textbf{Solution}



\[
Packet\ Loss\%= \frac{Lost\ Packets}{Total\ Packets}\times100
\]





\[
=\frac{75}{5000}\times100
\]





\[
\boxed{1.5\%}
\]



\subsubsection*{Problem 49: Calculate subnet capacity for a point-to-point connection}

\textbf{Question:} Two routers need an IPv4 point-to-point link. Which prefix can provide exactly two usable addresses without wasting a conventional four-address \texttt{/30} allocation?

\textbf{Solution}

A \texttt{/31} subnet contains:



\[
2^{32-31}=2
\]



It provides two addresses that can both be used as endpoint addresses on a point-to-point link under RFC 3021.

Example:

\begin{lstlisting}
Network: 10.0.0.0/31
Endpoint 1: 10.0.0.0
Endpoint 2: 10.0.0.1
\end{lstlisting}



\[
\boxed{/31}
\]



A \texttt{/30} is also commonly used for point-to-point links, but it has two network-related addresses under conventional subnetting rules and provides two usable host addresses.

\subsubsection*{Problem 50: Calculate the total number of hosts across multiple subnets}

\textbf{Question:} An organization has 6 identical \texttt{/26} subnets. Calculate the total number of usable host addresses.

\textbf{Solution}

Usable hosts in one \texttt{/26} subnet:



\[
2^{32-26}-2=64-2=62
\]



For six subnets:



\[
6\times62=372
\]





\[
\boxed{372\ usable\ host\ addresses}
\]




\newpage
\begin{center}
\vspace{0.5em}
{\color{primary}\rule{0.45\textwidth}{1.2pt}\\[4pt]}
{\Large\bfseries\color{primary} Section XXIII: Master Networking and Subnetting Revision Sheet\\[2pt]}
{\color{primary}\rule{0.45\textwidth}{1.2pt}}
\vspace{0.5em}
\end{center}

\section{Master Networking and Subnetting Revision Sheet}

\subsection{Comprehensive Quick-Revision Reference Table}

\begin{table}[htbp]
\centering
\caption{Comprehensive Examination Revision Sheet: Networking, Modulation, Routing, and Services}
\label{tab:master_exam_revision_sheet}
\begin{tabularx}{\textwidth}{p{4.2cm} X}
\toprule
\textbf{Core Topic / Concept} & \textbf{Key Examination Rule / Invariant} \\
\midrule
\textbf{Amplitude Modulation (AM)} & Modifies carrier amplitude $A_c$; frequency and phase remain constant; $m = A_m / A_c$ \\
\textbf{Frequency Modulation (FM)} & Modifies carrier frequency $f_c$; amplitude remains constant; Carson's BW $\approx 2(\Delta f + f_m)$ \\
\textbf{Digital Modulation (QAM)} & Modifies both amplitude and phase; bits/symbol $= \log_2 M$ \\
\textbf{Baud vs. Bit Rate} & $\text{Bit Rate} = \text{Baud Rate} \times \text{Bits per Symbol} = R_s \log_2 M$ \\
\textbf{Nyquist Sampling Theorem} & Minimum sampling frequency for perfect reconstruction: $f_s \ge 2 f_{\max}$ \\
\textbf{PCM Bit Rate} & $\text{Bit Rate} = f_s \times n$ (sampling rate $\times$ bits per sample; voice standard = 64 kbps) \\
\textbf{Longest Prefix Match (LPM)} & Hardware forwarding selects the route with the maximum number of leading network bits \\
\textbf{Default IPv4 Route} & \texttt{0.0.0.0/0} (least specific matching prefix in IPv4) \\
\textbf{DNS Port Allocation} & Port 53 UDP (standard client queries) and Port 53 TCP (zone transfers and large responses) \\
\textbf{Primary DNS Records} & A (IPv4), AAAA (IPv6), CNAME (alias), MX (mail), NS (nameserver), PTR (reverse lookup) \\
\textbf{NXDOMAIN Response} & Authoritative signal that the requested domain name does not exist \\
\textbf{Static NAT} & Permanent, fixed 1-to-1 mapping between private IP and public IP \\
\textbf{Dynamic NAT} & Temporary, stateful 1-to-1 mapping assigned from an available public IP pool \\
\textbf{PAT (NAT Overload)} & Many private hosts share a single public IP distinguished by transport source port numbers \\
\textbf{Source NAT (SNAT)} & Rewrites source IP address; used for outbound Internet access from private subnets \\
\textbf{Destination NAT (DNAT)} & Rewrites destination IP address; standard mechanism for Port Forwarding \\
\textbf{Carrier-Grade NAT (CGNAT)} & Double NAT operated by ISPs using reserved shared space \textbf{100.64.0.0/10} (RFC 6598) \\
\bottomrule
\end{tabularx}
\end{table}

\subsection{Comprehensive Formula and Invariant Quick-Reference}

\begin{table}[htbp]
\centering
\caption{IPv4 Addressing and Subnetting Formula Reference}
\label{tab:formulas_ipv4_subnetting}
\begin{tabularx}{\textwidth}{p{4.8cm} X}
\toprule
\textbf{Subnetting Property} & \textbf{Mathematical Formula / Rule} \\
\midrule
Host bits ($h$) & $h = 32 - p$, where $p$ is the CIDR prefix length \\
Total addresses per subnet & $N_{\text{total}} = 2^h = 2^{32-p}$ \\
Usable host addresses & $N_{\text{usable}} = 2^h - 2$ (excluding network and broadcast IDs) \\
Subnets created by borrowing $s$ bits & $N_{\text{subnets}} = 2^s$, where $s = p_{\text{new}} - p_{\text{old}}$ \\
Fourth-octet block size & $\text{Block Size} = 256 - \text{Mask Value} = 2^h$ \\
Supernet aggregate prefix & $p_{\text{super}} = p_{\text{base}} - \log_2(k)$, for $k$ contiguous aligned subnets \\
\bottomrule
\end{tabularx}
\end{table}

\begin{table}[htbp]
\centering
\caption{Transmission and Latency Formula Reference}
\label{tab:formulas_transmission_latency}
\begin{tabularx}{\textwidth}{p{4.8cm} X}
\toprule
\textbf{Performance Metric} & \textbf{Mathematical Formula} \\
\midrule
Transmission delay ($T_t$) & $T_t = \frac{L}{R}$ (Packet size in bits / Transmission rate in bps) \\
Propagation delay ($T_p$) & $T_p = \frac{d}{v}$ (Physical distance / Propagation speed) \\
Total nodal delay ($D_{\text{total}}$) & $D_{\text{total}} = D_{\text{proc}} + D_{\text{queue}} + D_{\text{trans}} + D_{\text{prop}}$ \\
Round-Trip Time (RTT) & $\text{RTT} = D_{\text{forward}} + D_{\text{return}} \approx 2 D_{\text{one-way}}$ \\
Bandwidth-Delay Product (BDP) & $\text{BDP} = R \times \text{RTT}$ (Bits in flight required to saturate pipe) \\
Throughput ($S$) & $S = \frac{\text{Data Delivered}}{\text{Elapsed Time}}$ \\
Annual Downtime ($T_{\text{down}}$) & $T_{\text{down}} = 525,600\text{ min} \times (1 - \text{Availability})$ \\
\bottomrule
\end{tabularx}
\end{table}

\begin{table}[htbp]
\centering
\caption{Digital Modulation and Communication Formula Reference}
\label{tab:formulas_digital_communication}
\begin{tabularx}{\textwidth}{p{4.8cm} X}
\toprule
\textbf{Communication Concept} & \textbf{Mathematical Formula} \\
\midrule
Bits per constellation symbol & $n = \log_2 M$, where $M$ is the number of constellation points \\
Gross Bit Rate ($R_b$) & $R_b = R_s \times \log_2 M$ (Baud rate $\times$ bits per symbol) \\
PCM Bit Rate & $R_b = f_s \times n_{\text{bits}}$ (Sampling frequency $\times$ bits/sample) \\
Nyquist Sampling Criterion & $f_s \ge 2 f_{\max}$ (Minimum sampling rate for reconstruction) \\
AM Modulation Index ($m$) & $m = \frac{A_m}{A_c} = \frac{E_{\max} - E_{\min}}{E_{\max} + E_{\min}}$ \\
Carson's Rule for FM Bandwidth & $BW \approx 2(\Delta f + f_m)$, where $\Delta f$ is frequency deviation \\
\bottomrule
\end{tabularx}
\end{table}

\begin{table}[htbp]
\centering
\caption{Data-Link Reliability and Flow Control Formula Reference}
\label{tab:formulas_reliable_transmission}
\begin{tabularx}{\textwidth}{p{4.8cm} X}
\toprule
\textbf{Reliability Metric} & \textbf{Mathematical Formula} \\
\midrule
Stop-and-Wait cycle time ($T_{\text{cycle}}$) & $T_{\text{cycle}} = T_t + 2 T_p$ (neglecting ACK transmission and processing) \\
Delay-bandwidth ratio parameter ($a$) & $a = \frac{T_p}{T_t}$ (Ratio of propagation delay to transmission delay) \\
Stop-and-Wait utilization ($U$) & $U = \frac{T_t}{T_t + 2 T_p} = \frac{1}{1 + 2a}$ \\
Stop-and-Wait throughput ($S$) & $S = R \times U = \frac{L}{T_{\text{cycle}}}$ \\
Optimal sliding window size ($W_{\text{opt}}$) & $W_{\text{opt}} \ge 1 + 2a$ (Pipelining required to achieve 100\% utilization) \\
\bottomrule
\end{tabularx}
\end{table}

\subsection{Examination Priority Matrix and Revision Methodology}

\begin{table}[htbp]
\centering
\caption{Technical Recruitment Examination Priority Matrix for Networking}
\label{tab:priority_matrix_networking}
\begin{tabularx}{\textwidth}{p{2.5cm} p{3.8cm} X}
\toprule
\textbf{Priority Tier} & \textbf{Core Topic Area} & \textbf{Critical Competencies Tested} \\
\midrule
\textbf{Critical (Top 1)} & IPv4 Subnetting \& VLSM & Binary conversions, usable hosts ($2^h-2$), prefix calculation, route summarization \\
\textbf{Critical (Top 2)} & OSI \& TCP/IP Models & Layer encapsulation, protocol mappings, TCP handshake, TCP vs. UDP \\
\textbf{Critical (Top 3)} & Routing \& Protocols & OSPF (Dijkstra, areas, cost), BGP (attributes, path selection), STP root election \\
\textbf{High} & DNS, DHCP, \& NAT & DORA sequence, port numbers, static/dynamic/PAT, DNS record types, diagnostics \\
\textbf{High} & Communication Latency & $T_t = L/R$, $T_p = d/v$, total nodal delay, Bandwidth-Delay Product (BDP) \\
\textbf{Medium--High} & Network Security & Firewall rules, DMZ, IDS vs. IPS, DHCP Snooping, Dynamic ARP Inspection \\
\textbf{Medium} & Modulation \& Stop-and-Wait & Baud rate vs. bit rate, Nyquist sampling, Stop-and-Wait efficiency ($1/(1+2a)$) \\
\bottomrule
\end{tabularx}
\end{table}

\subsubsection{Recommended Final Revision Strategy}

\begin{enumerate}
    \item \textbf{Master subnetting first:} Practice prefix lengths \texttt{/24}, \texttt{/25}, \texttt{/26}, \texttt{/27}, \texttt{/28}, \texttt{/29}, \texttt{/30}, \texttt{/31}, and \texttt{/32} until address ranges, host capacities, and boundary calculations become second nature.
    \item \textbf{Memorize protocol numbers and transport port allocations:} Memorize standard ports: TCP 20/21 (FTP), 22 (SSH), 23 (Telnet), 25 (SMTP), 53 (DNS), 80 (HTTP), 110 (POP3), 143 (IMAP), 443 (HTTPS), 179 (BGP); UDP 53 (DNS), 67/68 (DHCP), 123 (NTP), and 161 (SNMP); along with IP protocol numbers 1 (ICMP), 6 (TCP), 17 (UDP), and 89 (OSPF).
    \item \textbf{Practice routing and switching scenarios:} Understand STP Root Bridge election (lowest Bridge ID), port states, OSPF Dijkstra cost calculation, BGP path attributes (\texttt{WEIGHT}, \texttt{LOCAL\_PREF}, \texttt{AS\_PATH}, \texttt{MED}), and Longest Prefix Match lookup mechanics.
    \item \textbf{Practice quantitative calculations:} Solve latency, propagation, bandwidth-delay product, bit rate versus baud rate, Nyquist sampling frequency, PCM voice bit rate, and Stop-and-Wait ARQ utilization formulas repeatedly.
    \item \textbf{Master systematic troubleshooting workflows:} Be prepared to diagnose end-to-end connectivity failures across physical media, Layer 2 switching loops, DHCP APIPA allocation issues, DNS NXDOMAIN resolution failures, default gateway routing reachability, and stateful firewall/NAT drops.
\end{enumerate}

